% El deporte, el ocio y el bienestar (continuación) : expresión oral y escrita — Espagnol 1ère A — Cours signé OnBuch+
% Programme officiel MINESEC. Compiler avec : tectonic EA08-deporte-ocio-bienestar-continuacion.tex
\newcommand{\DOCMATIERE}{Espagnol}
\newcommand{\DOCNIVEAU}{1ère A}
\newcommand{\DOCMODULE}{Módulo 2 : Environnement, santé et bien-être}
\newcommand{\DOCLECON}{EA08}
\newcommand{\DOCTITRE}{El deporte, el ocio y el bienestar (continuación) : expresión oral y escrita}
% OnBuch+ — préambule « cours signés » ENRICHI (compilé avec tectonic / XeLaTeX)
% Système visuel « Pop » d'OnBuch+ : fond crème (jamais blanc), bordures encre
% épaisses, ombres « dures » décalées, titres Archivo Black en capitales.
% Version MATHÉMATIQUES : graphes (pgfplots/tikz), boîtes pédagogiques
% (définition, propriété, méthode, attention, le savais-tu, exercice résolu,
% exercice, TP, bilan), page de garde et signature OnBuch+ ancrée en bas.
%
% Macros attendues dans main.tex AVANT \input{preamble} :
%   \DOCMATIERE  \DOCNIVEAU  \DOCMODULE  \DOCLECON (numéro)  \DOCTITRE
\documentclass[11pt]{article}

\usepackage{fontspec}
\usepackage[spanish,french]{babel} % français = langue principale ; l'espagnol via \esp{..} / espagnol
\usepackage[a4paper,margin=2cm,top=3.4cm,bottom=2.8cm,headheight=1.4cm,footskip=1.3cm]{geometry}
\usepackage{amsmath,amssymb,amsfonts}
\usepackage{mathtools} % \xmapsto, \coloneqq... souvent utilisés par les modèles
\usepackage{cancel} % simplifications barrées (bilans de forces)
% \abs{z} (module, valeur absolue) et \norm{u} (norme), utilisés par les modèles
\providecommand{\abs}{}\let\abs\relax\DeclarePairedDelimiter{\abs}{\lvert}{\rvert}
\providecommand{\norm}{}\let\norm\relax\DeclarePairedDelimiter{\norm}{\lVert}{\rVert}
\usepackage[table]{xcolor} % [table] : fournit aussi \rowcolors (alternance de lignes)
\usepackage{pagecolor}
\usepackage{tikz}
\usetikzlibrary{arrows.meta,positioning,calc,shapes.geometric,shapes.misc,shapes.arrows,decorations.pathmorphing,decorations.pathreplacing,decorations.markings,patterns,shadows,mindmap,trees,fit,backgrounds,matrix,intersections,angles,quotes}
\usepackage{pgfplots}
\usepackage[version=4]{mhchem} % équations nucléaires \ce{^{235}_{92}U}
\usepackage[european,siunitx]{circuitikz} % schémas électriques
\pgfplotsset{compat=1.18}
\usepgfplotslibrary{fillbetween,groupplots} % aires entre courbes (calcul intégral)
\usepackage{enumitem}
\usepackage{fancyhdr}
% [most] charge toutes les bibliothèques tcolorbox (listings, minted, raster,
% documentation...) et gonfle inutilement la mémoire TeX sur un document long
% (~300 boîtes) : on ne charge que ce qui est réellement utilisé.
\usepackage[skins,breakable]{tcolorbox}
\usepackage{graphicx}
\usepackage{colortbl}
\usepackage{booktabs}
\usepackage{tabularx}
\usepackage{diagbox} % case d'en-tête diagonale des tableaux à double entrée
% Colonnes extensibles centrée / gauche / droite (C, L, R), souvent utilisées par les modèles
\newcolumntype{C}{>{\centering\arraybackslash}X}
\newcolumntype{L}{>{\raggedright\arraybackslash}X}
\newcolumntype{R}{>{\raggedleft\arraybackslash}X}
\usepackage{array}
\usepackage{multicol}
\usepackage{siunitx}
% parse-numbers=false : accepte aussi \SI{2,5 \times 10^{-3}}{...}
\sisetup{locale=FR,output-decimal-marker={,},group-minimum-digits=4,parse-numbers=false}
\DeclareSIUnit\ohm{\ensuremath{\symup{\Omega}}} % U+2126 absent de la police
\DeclareSIUnit\division{div} % réglages d'oscilloscope (ms/div, V/div)
\DeclareSIUnit\tr{tr} % tours (vitesse de rotation, ex. \SI{1500}{\tr\per\minute})
\DeclareSIUnit\molar{M} % concentration molaire (ex. \SI{3}{\molar}, \SI{1}{\micro\molar})
\DeclareSIUnit\an{an} % année (le modèle confond parfois avec \year, un compteur TeX) — ex. \SI{5}{\kilo\meter\per\an}
\DeclareSIUnit\FCFA{FCFA} % franc CFA (ex. \SI{500}{\FCFA\per\kilo\gram})
\DeclareSIUnit\degreeBrix{\textdegree Brix} % teneur en sucre d'un sirop/jus (réfractométrie)
\DeclareSIUnit\month{mois} % le modèle utilise parfois \month, un compteur TeX, comme unité de durée
\usepackage{qrcode}
% \degree : commande fréquemment utilisée par les modèles pour un angle ou une
% température en dehors de \SI{}{} (non fournie par les paquets chargés).
\providecommand{\degree}{^{\circ}}
% \qed (fin de démonstration), utilisé par les modèles sans amsthm
\providecommand{\qed}{\hfill$\square$}
% \barre : double barre des valeurs interdites dans un tableau de variations (tkz-tab)
\providecommand{\barre}{\ensuremath{\|}}
% environnement proof (amsthm non chargé) utilisé par les modèles
\ifdefined\proof\else\newenvironment{proof}[1][Démonstration]{\par\noindent\textit{#1.}\ }{\qed\par}\fi
\usepackage{lastpage}
\usepackage{hyperref}
\hypersetup{hidelinks}

% ============================================================
% Palette « Pop » OnBuch+ — mêmes hex que la classe P (plus_ui.dart)
% ============================================================
\definecolor{poporange}{HTML}{F59321}
\definecolor{popdark}{HTML}{E07A0C}
\definecolor{popcream}{HTML}{FFF6E8}
\definecolor{popink}{HTML}{1C1714}
\definecolor{poppurple}{HTML}{7A5AE0}
\definecolor{popgreen}{HTML}{2E9E6A}
\definecolor{poppink}{HTML}{E0457B}
\definecolor{popblue}{HTML}{2F7DE1}
\definecolor{popgold}{HTML}{C98A12}
\definecolor{popmuted}{HTML}{8A7F76}
\definecolor{popline}{HTML}{EADFD2}
\definecolor{popgray}{HTML}{8A7F76}
\colorlet{popgrey}{popgray}
% Alias tolérants (le modèle nomme parfois les couleurs différemment)
\colorlet{popwhite}{white}
\colorlet{popblack}{popink}
\colorlet{popred}{poppink}
\colorlet{popyellow}{popgold}
% Teintes claires (fonds de tableaux/encadrés) — le modèle nomme parfois
% les couleurs avec un suffixe L (ex. popblueL) sans les avoir définies.
\colorlet{poporangeL}{poporange!20}
\colorlet{popdarkL}{popdark!20}
\colorlet{poppurpleL}{poppurple!20}
\colorlet{popgreenL}{popgreen!20}
\colorlet{poppinkL}{poppink!20}
\colorlet{popblueL}{popblue!20}
\colorlet{popgoldL}{popgold!20}
\colorlet{popredL}{popred!20}
\colorlet{popgrayL}{popgray!20}
% Teintes claires pour fonds de tableaux / zones
\colorlet{popblueL}{popblue!10!white}
\colorlet{poporangeL}{poporange!14!white}
\colorlet{popgreenL}{popgreen!12!white}
\colorlet{poppurpleL}{poppurple!12!white}
\colorlet{poppinkL}{poppink!12!white}
\colorlet{popgoldL}{popgold!14!white}

\pagecolor{popcream}

% ============================================================
% Typographie
% ============================================================
\defaultfontfeatures{Ligatures=TeX}
\setmainfont{PlusJakartaSans-Medium.ttf}[Path=../../../fonts/,BoldFont=PlusJakartaSans-Bold.ttf]
% Maths en sans-serif (Fira Math) avec chiffres et lettres droites de Plus
% Jakarta, pour que les formules se fondent dans le texte.
\usepackage[math-style=ISO,bold-style=ISO]{unicode-math}
\setmathfont{FiraMath-Regular.otf}[Path=../../../fonts/]
\setmathfont{PlusJakartaSans-Medium.ttf}[Path=../../../fonts/,range={up/num,up/{Latin,latin},\mathup}]
\usepackage{icomma} % virgule décimale sans espace en mode math
% Caractères mathématiques Unicode absents des polices de texte (Plus Jakarta,
% Archivo Black) : rendus par leur équivalent mathématique, en texte comme en
% formule — sinon ils s'affichent en carré vide (ex. « Arithmétique dans ℤ »).
\usepackage{newunicodechar}
\newunicodechar{ℝ}{\ensuremath{\mathbb{R}}}
\newunicodechar{ℤ}{\ensuremath{\mathbb{Z}}}
\newunicodechar{ℂ}{\ensuremath{\mathbb{C}}}
\newunicodechar{ℕ}{\ensuremath{\mathbb{N}}}
\newunicodechar{ℚ}{\ensuremath{\mathbb{Q}}}
\newunicodechar{𝔻}{\ensuremath{\mathbb{D}}}
\newunicodechar{α}{\ensuremath{\alpha}}
\newunicodechar{β}{\ensuremath{\beta}}
\newunicodechar{γ}{\ensuremath{\gamma}}
\newunicodechar{δ}{\ensuremath{\delta}}
\newunicodechar{ε}{\ensuremath{\varepsilon}}
\newunicodechar{θ}{\ensuremath{\theta}}
\newunicodechar{λ}{\ensuremath{\lambda}}
\newunicodechar{μ}{\ensuremath{\mu}}
\newunicodechar{σ}{\ensuremath{\sigma}}
\newunicodechar{τ}{\ensuremath{\tau}}
\newunicodechar{φ}{\ensuremath{\varphi}}
\newunicodechar{ψ}{\ensuremath{\psi}}
\newunicodechar{ω}{\ensuremath{\omega}}
\newunicodechar{Σ}{\ensuremath{\Sigma}}
% majuscules produites par \MakeUppercase dans les titres (α-aminés -> Α)
\newunicodechar{Α}{\ensuremath{\alpha}}
\newunicodechar{Β}{\ensuremath{\beta}}
\newunicodechar{Γ}{\ensuremath{\Gamma}}
\newunicodechar{Δ}{\ensuremath{\Delta}}
\newunicodechar{Ε}{\ensuremath{\varepsilon}}
\newunicodechar{Θ}{\ensuremath{\Theta}}
\newunicodechar{Λ}{\ensuremath{\Lambda}}
\newunicodechar{Μ}{\ensuremath{\mu}}
\newunicodechar{Τ}{\ensuremath{\tau}}
\newunicodechar{Φ}{\ensuremath{\Phi}}
\newunicodechar{Ψ}{\ensuremath{\Psi}}
\newunicodechar{Ω}{\ensuremath{\Omega}}
\newunicodechar{↦}{\ensuremath{\mapsto}}
\newunicodechar{∈}{\ensuremath{\in}}
\newunicodechar{∉}{\ensuremath{\notin}}
\newunicodechar{∧}{\ensuremath{\wedge}}
\newunicodechar{⇔}{\ensuremath{\Leftrightarrow}}
\newunicodechar{⟺}{\ensuremath{\Longleftrightarrow}}
\newunicodechar{⇒}{\ensuremath{\Rightarrow}}
\newunicodechar{⟹}{\ensuremath{\Longrightarrow}}
\newunicodechar{∘}{\ensuremath{\circ}}
\newunicodechar{∪}{\ensuremath{\cup}}
\newunicodechar{∩}{\ensuremath{\cap}}
\newunicodechar{≡}{\ensuremath{\equiv}}
\newunicodechar{⊂}{\ensuremath{\subset}}
\newunicodechar{⊥}{\ensuremath{\perp}}
\newunicodechar{∃}{\ensuremath{\exists}}
\newunicodechar{∀}{\ensuremath{\forall}}
\newunicodechar{∛}{\ensuremath{\sqrt[3]{\,}}}
\newunicodechar{∜}{\ensuremath{\sqrt[4]{\,}}}
\newunicodechar{⃗}{\ensuremath{{}^{\to}}}
\newunicodechar{ⁿ}{\ifmmode^{n}\else\textsuperscript{\ensuremath{n}}\fi}
\newunicodechar{ˣ}{\ifmmode^{x}\else\textsuperscript{\ensuremath{x}}\fi}
\newunicodechar{ᵇ}{\ifmmode^{b}\else\textsuperscript{\ensuremath{b}}\fi}
\newunicodechar{ʳ}{\ifmmode^{r}\else\textsuperscript{\ensuremath{r}}\fi}
\newunicodechar{ᵘ}{\ifmmode^{u}\else\textsuperscript{\ensuremath{u}}\fi}
\newunicodechar{ʸ}{\ifmmode^{y}\else\textsuperscript{\ensuremath{y}}\fi}
\newunicodechar{ᵃ}{\ifmmode^{a}\else\textsuperscript{\ensuremath{a}}\fi}
\newunicodechar{ᵉ}{\ifmmode^{e}\else\textsuperscript{\ensuremath{e}}\fi}
\newunicodechar{ᵏ}{\ifmmode^{k}\else\textsuperscript{\ensuremath{k}}\fi}
\newunicodechar{ᵖ}{\ifmmode^{p}\else\textsuperscript{\ensuremath{p}}\fi}
\newunicodechar{ᴺ}{\ifmmode^{N}\else\textsuperscript{\ensuremath{N}}\fi}
\newunicodechar{ᶜ}{\ifmmode^{c}\else\textsuperscript{\ensuremath{c}}\fi}
\newunicodechar{ʰ}{\ifmmode^{h}\else\textsuperscript{\ensuremath{h}}\fi}
\newunicodechar{ᵠ}{\ifmmode^{q}\else\textsuperscript{\ensuremath{q}}\fi}
\newunicodechar{ₙ}{\ifmmode_{n}\else\textsubscript{\ensuremath{n}}\fi}
\newunicodechar{ₐ}{\ifmmode_{a}\else\textsubscript{\ensuremath{a}}\fi}
\newunicodechar{ᵢ}{\ifmmode_{i}\else\textsubscript{\ensuremath{i}}\fi}
\newunicodechar{ᵣ}{\ifmmode_{r}\else\textsubscript{\ensuremath{r}}\fi}
\newunicodechar{ₖ}{\ifmmode_{k}\else\textsubscript{\ensuremath{k}}\fi}
\newunicodechar{ᵦ}{\ifmmode_{\beta}\else\textsubscript{\ensuremath{\beta}}\fi}
\newunicodechar{ₓ}{\ifmmode_{x}\else\textsubscript{\ensuremath{x}}\fi}
\newfontfamily\popdisplay{ArchivoBlack-Regular.ttf}[Path=../../../fonts/]
\newfontfamily\poplogo{SpaceGrotesk-Bold.ttf}[Path=../../../fonts/]
\newfontfamily\popsemi{PlusJakartaSans-SemiBold.ttf}[Path=../../../fonts/]
\newfontfamily\popxbold{PlusJakartaSans-ExtraBold.ttf}[Path=../../../fonts/]
\newfontfamily\popxboldls{PlusJakartaSans-ExtraBold.ttf}[Path=../../../fonts/,LetterSpace=3.0]

\setlist[enumerate]{leftmargin=1.7em,itemsep=5pt,topsep=4pt}
\setlist[itemize]{leftmargin=1.7em,itemsep=4pt,topsep=4pt,label={\color{poporange}\textbullet}}
\setlength{\parindent}{0pt}
\setlength{\parskip}{5pt}
\renewcommand{\arraystretch}{1.25}

% pgfplots : style Pop par défaut
\pgfplotsset{
  every axis/.append style={axis line style={popink,line width=1pt},tick style={popink},
    grid style={popline,line width=0.5pt},label style={font=\small},tick label style={font=\footnotesize}},
  popcurve/.style={poporange,line width=1.8pt,no markers,smooth},
}
% Même style disponible dans un \draw TikZ simple (hors pgfplots)
\tikzset{popcurve/.style={poporange,line width=1.8pt}}
\tikzset{popgrid/.style={popline,very thin}}
\colorlet{popcurve}{poporange} % le nom du style est parfois utilisé comme couleur

% ============================================================
% Pill / squiggle
% ============================================================
\newcommand{\poppill}[3]{%
  \tikz[baseline=-0.55ex]\node[draw=popink,line width=1.2pt,rounded corners=2.6pt,%
    fill=#2,inner xsep=6.5pt,inner ysep=3pt]{%
    \popxboldls\fontsize{7}{7}\selectfont\color{#3}\MakeUppercase{#1}};%
}
\newcommand{\popsquiggle}[1]{%
  \tikz[baseline=-0.4ex]{
    \draw[#1,line width=2.6pt,line cap=round]
      plot [smooth] coordinates {(0,0.09) (0.34,0.01) (0.68,0.21) (1.02,0.01) (1.36,0.10)};
  }%
}
% Mot-clé surligné (marqueur orange sous le texte)
\newcommand{\cle}[1]{{\popxbold\color{popdark}#1}}

% ============================================================
% En-tête / pied
% ============================================================
\pagestyle{fancy}
\fancyhf{}
\renewcommand{\headrulewidth}{0pt}
\renewcommand{\footrulewidth}{0pt}
\lhead{\raisebox{-0.15ex}{\poplogo\fontsize{13}{13}\selectfont\color{popink}OnBuch\color{poporange}+}}
\rhead{\poppill{\DOCMATIERE\ \textbullet\ \DOCNIVEAU\ \textbullet\ Leçon \DOCLECON}{white}{popink}}
\lfoot{\popxbold\fontsize{7.6}{7.6}\selectfont\color{popmuted}\MakeUppercase{Cours signé OnBuch+}\ \textbullet\ {\popsemi\DOCTITRE}}
\rfoot{\tikz[baseline=-0.6ex]\node[rounded corners=8.5pt,fill=popink,minimum height=17pt,minimum width=17pt,inner xsep=5pt,inner sep=0]{\popxbold\fontsize{8}{8}\selectfont\color{popcream}\thepage\,/\,\pageref*{LastPage}};}

% ============================================================
% Titres
% ============================================================
\newcommand{\chaptitre}[2]{%
  \begin{tcolorbox}[enhanced,colback=poporange,colframe=popink,boxrule=2.6pt,arc=11pt,
      drop shadow=popink,left=15pt,right=15pt,top=12pt,bottom=11pt,before skip=3pt,after skip=18pt]
    \poppill{#2}{popink}{popcream}\par\vspace{9pt}
    {\raggedright\hyphenpenalty=10000\exhyphenpenalty=10000\popdisplay\fontsize{22}{24}\selectfont\color{popcream}\MakeUppercase{#1}\par}
  \end{tcolorbox}
}
% Section numérotée : pastille orange avec le numéro + titre Archivo
\newcounter{popsec}
\newcommand{\coursec}[1]{%
  \stepcounter{popsec}\setcounter{popsub}{0}%
  \par\vspace{16pt}\noindent
  \begin{minipage}{\linewidth}
  \tikz[baseline=-0.7ex]\node[draw=popink,line width=1.6pt,fill=poporange,rounded corners=4pt,minimum size=22pt,inner sep=0]{\popdisplay\fontsize{12}{12}\selectfont\color{popcream}\thepopsec};%
  \hspace{8pt}{\popdisplay\fontsize{15}{17}\selectfont\color{popink}\MakeUppercase{#1}}\par
  \vspace{4pt}\noindent{\color{poporange}\rule{36pt}{3.6pt}}
  \end{minipage}\par\nopagebreak\vspace{8pt}
}
\newcounter{popsub}
\newcommand{\courssub}[1]{%
  \stepcounter{popsub}%
  \par\vspace{9pt}\noindent{\popxbold\fontsize{11.5}{13}\selectfont\color{popink}{\color{poporange}\thepopsec.\thepopsub}\hspace{6pt}#1}\par\nopagebreak\vspace{4pt}
}

% ============================================================
% Boîtes pédagogiques — même recette : fond blanc, bordure épaisse colorée,
% ombre dure encre, badge plein. Titre optionnel [#1].
% ============================================================
\tcbset{popbase/.style={breakable,enhanced,colback=white,boxrule=2.3pt,arc=9pt,
  shadow={3pt}{-3pt}{0pt}{popink},left=14pt,right=14pt,top=11pt,bottom=11pt,before skip=11pt,after skip=15pt,
  fontupper=\popsemi\fontsize{10.6}{14.5}\selectfont\color{popink},
  fontlower=\popsemi\fontsize{10.6}{14.5}\selectfont\color{popink}}}
% Coche dessinée (le glyphe ✓ n'existe pas dans les polices du document)
\newcommand{\popcheck}[1]{\tikz[baseline=-0.5ex]\draw[#1,line width=1.6pt,line cap=round,line join=round](-0.13,0.0)--(-0.04,-0.09)--(0.13,0.1);}
\providecommand{\coche}{\popcheck{popgreen}}
\providecommand{\resultat}[1]{\par\smallskip\noindent\fcolorbox{popgreen}{popgreenL}{\parbox{\dimexpr\linewidth-2\fboxsep-2\fboxrule\relax}{\textbf{Résultat.} #1}}\par\smallskip}
\newcommand{\popboxhead}[3]{\poppill{#1}{#2}{white}\ifx\relax#3\relax\else\hspace{6pt}{\popxbold\fontsize{10.6}{12}\selectfont\color{popink}#3}\fi\par\vspace{7pt}}

\newtcolorbox{aretenir}[1][]{popbase,colframe=popgreen,before upper={\popboxhead{À retenir}{popgreen}{#1}}}
% Texte en espagnol (document de lecture, dialogue, poème, extrait) : cadre
% d'étude. Un titre optionnel (source, auteur) s'affiche dans l'en-tête.
\newtcolorbox{textebox}[1][]{popbase,colframe=popblue,before upper={\popboxhead{Texto}{popblue}{#1}\begin{otherlanguage}{spanish}},after upper={\end{otherlanguage}}}
% Marques ✓ / ✗ (forme correcte / fautive) souvent utilisées par les modèles
\providecommand{\cross}{\textcolor{poppink}{\ensuremath{\times}}}
\providecommand{\xmark}{\cross}\providecommand{\crossmark}{\cross}\providecommand{\cmark}{\ensuremath{\checkmark}}
% Ligne de réponse / justification à compléter (inventée par les modèles)
\providecommand{\justification}[1]{\par\noindent\textit{Justification~:} \dotfill\par}
% \textipa (package tipa non chargé) : la police ne couvre pas l'API -> simple passage
\providecommand{\textipa}[1]{#1}
% Soulignement (ulem non chargé) : \uline, \ul
\providecommand{\uline}[1]{\underline{#1}}\providecommand{\ul}[1]{\underline{#1}}
% Ordinaux français écrits en macro par les modèles : 1\ière, 3\ième
\newcommand{\ière}{\textsuperscript{re}}\newcommand{\ième}{\textsuperscript{e}}
% Mot ou phrase en espagnol à l'intérieur d'un paragraphe français
\newcommand{\esp}[1]{\ifmmode\text{\textit{#1}}\else\begingroup\selectlanguage{spanish}\itshape #1\endgroup\fi}
\newtcolorbox{exemplebox}[1][]{popbase,colframe=poporange,before upper={\popboxhead{Exemple}{poporange}{#1}}}
\newtcolorbox{definition}[1][]{popbase,colframe=popblue,before upper={\popboxhead{Définition}{popblue}{#1}}}
\newtcolorbox{propriete}[1][]{popbase,colframe=poppurple,before upper={\popboxhead{Propriété}{poppurple}{#1}}}
\newtcolorbox{methode}[1][]{popbase,colframe=popgold,before upper={\popboxhead{Méthode}{popgold}{#1}}}
\newtcolorbox{attention}[1][]{popbase,colframe=poppink,before upper={\popboxhead{Attention}{poppink}{#1}}}
\newtcolorbox{experience}[1][]{popbase,colframe=popblue,colback=popblueL,before upper={\popboxhead{Expérience}{popblue}{#1}}}
% « Le savais-tu ? » : carte violette pleine (contexte, culture, Cameroun)
\newtcolorbox{savaistu}[1][]{popbase,colback=poppurple,colframe=popink,
  fontupper=\popsemi\fontsize{10.4}{14.2}\selectfont\color{white},
  before upper={\poppill{Le savais-tu\,?}{white}{poppurple}\ifx\relax#1\relax\else\hspace{6pt}{\popxbold\color{white}#1}\fi\par\vspace{7pt}}}
% Exercice résolu : énoncé en haut, \tcblower puis la solution sur fond crème
\newcounter{popexo}\newcounter{popexr}
\newtcolorbox{exoresolu}[1][]{popbase,colframe=poporange,colbacklower=popcream,
  segmentation style={popink,line width=1.2pt,dashed},
  before upper={\stepcounter{popexr}\popboxhead{Exercice résolu \thepopexr}{poporange}{#1}},
  before lower={\poppill{Solution}{popink}{popcream}\par\vspace{6pt}}}
% Exercice (non résolu en place) — la correction est dans la partie Corrigés
\newtcolorbox{exercice}[1][]{popbase,colframe=popink,
  before upper={\stepcounter{popexo}\popboxhead{Exercice \thepopexo}{popink}{#1}}}
% Corrigé
\newtcolorbox{corrige}[1][]{popbase,colframe=popgreen,colback=popgreenL,
  before upper={\popboxhead{Corrigé}{popgreen}{#1}}}
% Objectifs / prérequis / bilan
\newtcolorbox{objectifs}{popbase,colframe=popink,colback=poporangeL,
  before upper={\popboxhead{Objectifs de la leçon}{popink}{}}}
\newtcolorbox{prerequis}{popbase,colframe=popink,colback=popblueL,
  before upper={\popboxhead{Prérequis}{popblue}{}}}
\newtcolorbox{bilan}{popbase,colframe=popink,colback=white,boxrule=2.8pt,
  before upper={\popboxhead{Fiche bilan}{poporange}{L'essentiel à connaître pour l'examen}}}
% Figure encadrée avec légende
\newtcolorbox{popfigure}[1][]{enhanced,breakable=false,colback=white,colframe=popline,boxrule=1.4pt,arc=8pt,
  left=10pt,right=10pt,top=10pt,bottom=8pt,before skip=10pt,after skip=12pt,halign=center,
  lower separated=false,halign lower=center,
  fontlower=\popsemi\fontsize{9}{11.5}\selectfont\color{popmuted},
  #1}
\newcommand{\legende}[1]{\par\vspace{6pt}{\popsemi\fontsize{9}{11.5}\selectfont\color{popmuted}#1}}

% ============================================================
% Signature OnBuch+
% ============================================================
\newcommand{\poplogoinline}[1]{{\poplogo\fontsize{#1}{#1}\selectfont\color{popink}OnBuch\color{poporange}+}}
\newcommand{\signatureonbuch}{%
  \begin{tcolorbox}[enhanced,colback=popink,colframe=popink,boxrule=2.2pt,arc=10pt,
      drop shadow=poporange,left=14pt,right=14pt,top=10pt,bottom=10pt,before skip=0pt,after skip=0pt]
    \tikz[baseline=-0.7ex]\node[circle,fill=popgreen,draw=popcream,line width=1.3pt,minimum size=22pt,inner sep=0]{\popcheck{white}};%
    \hspace{9pt}\begin{minipage}[c]{0.62\linewidth}
      {\popxbold\fontsize{12}{13}\selectfont\color{popcream}Signé\ \textbullet\ {\poplogo OnBuch\color{poporange}+}}\par\vspace{2pt}
      {\raggedright\popsemi\fontsize{8}{10.5}\selectfont\color{popline}\MakeUppercase{Équipe pédagogique OnBuch+}\\\MakeUppercase{Conforme au programme officiel MINESEC}\par}
    \end{minipage}\hfill
    {\poplogo\fontsize{22}{22}\selectfont\color{popcream}OnBuch\color{poporange}+}
  \end{tcolorbox}%
}

% ============================================================
% Page de garde
% #1 titre · #2 matière · #3 niveau · #4 module · #5 n° leçon · #6 durée ·
% #7 description / objectifs (texte ou itemize)
% ============================================================
\newcommand{\pagegarde}[7]{%
  \thispagestyle{empty}
  \newgeometry{a4paper,margin=1.7cm,top=1.6cm,bottom=1.4cm}%
  \begin{tikzpicture}[remember picture,overlay]
    \fill[poporange!18] ([xshift=-3.2cm,yshift=-3.6cm]current page.north east) circle (4.6cm);
    \fill[poppurple!14] ([xshift=2.4cm,yshift=7.6cm]current page.south west) circle (3.8cm);
  \end{tikzpicture}%
  {\poplogo\fontsize{30}{30}\selectfont\color{popink}OnBuch\color{poporange}+}\hfill
  \raisebox{0.6ex}{\poppill{Cours signé \textbullet\ Premium}{popink}{popcream}}\par
  \vspace{1pt}\popsquiggle{poppurple}\par
  \vspace{8pt}
  \begin{tcolorbox}[enhanced,colback=poporange,colframe=popink,boxrule=2.8pt,arc=13pt,
      drop shadow=popink,left=18pt,right=18pt,top=12pt,bottom=13pt,after skip=10pt]
    \poppill{#2 \textbullet\ #3}{popink}{popcream}\hspace{5pt}\poppill{#4}{white}{popink}\par\vspace{8pt}
    {\popdisplay\fontsize{14}{14}\selectfont\color{popink}LEÇON #5}\par\vspace{4pt}
    {\raggedright\hyphenpenalty=10000\exhyphenpenalty=10000\popdisplay\fontsize{25}{27}\selectfont\color{popcream}\MakeUppercase{#1}\par}
    \vspace{8pt}
    {\popxbold\fontsize{9.5}{9.5}\selectfont\color{popink}\MakeUppercase{Durée conseillée : #6}}
  \end{tcolorbox}
  \begin{tcolorbox}[enhanced,colback=white,colframe=popink,boxrule=2.2pt,arc=10pt,
      drop shadow=popink,left=16pt,right=16pt,top=10pt,bottom=10pt,after skip=8pt]
    \poppill{Dans cette leçon}{poppurple}{white}\par\vspace{5pt}
    \popsemi\fontsize{9.3}{12.6}\selectfont\color{popink}#7
  \end{tcolorbox}
  \noindent\begin{minipage}[t]{0.487\linewidth}
    \begin{tcolorbox}[enhanced,colback=popgreen,colframe=popink,boxrule=2.2pt,arc=10pt,
        drop shadow=popink,left=11pt,right=11pt,top=10pt,bottom=10pt,height=3.1cm,valign=center]
      \tcbox[colback=white,colframe=popink,boxrule=1.3pt,arc=5pt,boxsep=0pt,nobeforeafter,
        left=3pt,right=3pt,top=3pt,bottom=3pt,valign=center]{\qrcode[height=1.9cm]{https://play.google.com/store/apps/details?id=com.luvvix.onbuch&hl=fr}}%
      \hspace{8pt}\begin{minipage}[c]{0.5\linewidth}
        {\popdisplay\fontsize{11}{12.5}\selectfont\color{white}\MakeUppercase{Télécharge OnBuch}}\par\vspace{4pt}
        {\raggedright\popsemi\fontsize{8.4}{11}\selectfont\color{white}Gratuit \textbullet\ Google Play\\Tuteur IA, annales, concours, résultats\par}
      \end{minipage}
    \end{tcolorbox}
  \end{minipage}\hfill
  \begin{minipage}[t]{0.487\linewidth}
    \begin{tcolorbox}[enhanced,colback=poppurple,colframe=popink,boxrule=2.2pt,arc=10pt,
        drop shadow=popink,left=11pt,right=11pt,top=10pt,bottom=10pt,height=3.1cm,valign=center]
      \tikz[baseline=-0.7ex]\node[circle,fill=white,draw=popink,line width=1.3pt,minimum size=1.6cm,inner sep=0]{\popdisplay\fontsize{24}{24}\selectfont\color{poppurple}+};%
      \hspace{8pt}\begin{minipage}[c]{0.6\linewidth}
        {\popdisplay\fontsize{11}{12.5}\selectfont\color{white}\MakeUppercase{Passe à OnBuch+}}\par\vspace{4pt}
        {\raggedright\popsemi\fontsize{8.4}{11}\selectfont\color{white}Tous les cours signés, Tuteur IA illimité, fiches PDF — onglet OnBuch+ de l'app\par}
      \end{minipage}
    \end{tcolorbox}
  \end{minipage}
  \vspace{8pt}
  \popcontenu
  \vfill
  \signatureonbuch
  \restoregeometry
  \newpage
}

% Rangée « Ce cours contient » (page de garde)
\newcommand{\poptile}[3]{% #1 couleur #2 gros texte #3 légende
  \begin{tcolorbox}[enhanced,colback=#1,colframe=popink,boxrule=2pt,arc=9pt,shadow={2.5pt}{-2.5pt}{0pt}{popink},
      left=6pt,right=6pt,top=9pt,bottom=9pt,halign=center,valign=center,height=1.75cm,width=\linewidth,nobeforeafter]
    {\popdisplay\fontsize{12.5}{13}\selectfont\color{white}\MakeUppercase{#2}}\par\vspace{3pt}
    {\centering\popsemi\fontsize{7.8}{9.5}\selectfont\color{white}#3\par}
  \end{tcolorbox}}
\newcommand{\popcontenu}{%
  \par\noindent{\popxboldls\fontsize{7.5}{7.5}\selectfont\color{popmuted}CE COURS SIGNÉ CONTIENT}\par\vspace{7pt}
  \noindent\begin{minipage}[t]{0.235\linewidth}\poptile{popblue}{Cours}{complet, illustré, pas à pas}\end{minipage}\hfill
  \begin{minipage}[t]{0.235\linewidth}\poptile{poporange}{Méthodes}{et exercices résolus}\end{minipage}\hfill
  \begin{minipage}[t]{0.235\linewidth}\poptile{poppink}{Exercices}{type examen + corrigés}\end{minipage}\hfill
  \begin{minipage}[t]{0.235\linewidth}\poptile{popgreen}{Fiche bilan}{l'essentiel à retenir}\end{minipage}\par}

% Sommaire
\newtcolorbox{sommaire}{enhanced,colback=white,colframe=popink,boxrule=2.2pt,arc=10pt,
  drop shadow=popink,left=18pt,right=18pt,top=13pt,bottom=13pt,before skip=6pt,after skip=20pt,
  before upper={\poppill{Sommaire}{poppurple}{white}\par\vspace{9pt}\popsemi\fontsize{10.6}{14.5}\selectfont\color{popink}}}

% Signature de fin de cours — poussée tout en bas de la dernière page
\newcommand{\findecours}{%
  \par\vfill\nopagebreak
  \begin{center}\popsquiggle{poporange}\end{center}\vspace{4pt}
  \signatureonbuch
}
\AtBeginDocument{\def\llbracket{\mathopen{[\mkern-3.5mu[}}\def\rrbracket{\mathclose{]\mkern-3.5mu]}}} % intervalles d'entiers (glyphe ⟦ absent de la police)
% Glyphes absents de Fira Math : redéfinis à partir de glyphes disponibles
\AtBeginDocument{%
  \def\setminus{\mathbin{\backslash}}\let\smallsetminus\setminus
  \def\vdots{\mathord{\vbox{\baselineskip3.2pt\lineskiplimit0pt\kern4pt\hbox{.}\hbox{.}\hbox{.}}}}%
  \def\ddots{\mathinner{\mkern1mu\raise6.5pt\hbox{.}\mkern2mu\raise3.6pt\hbox{.}\mkern2mu\raise0.7pt\hbox{.}\mkern1mu}}%
  \def\triangle{\mathord{\scalebox{0.9}{$\symup{\Delta}$}}}%
  \def\nabla{\mathord{\raisebox{\depth}{\scalebox{1}[-1]{$\symup{\Delta}$}}}}%
  \def\checkmark{\popcheck{popdark}}%
  \def\star{\mathbin{\ast}}\def\bigstar{\mathord{\ast}}%
  \def\diamond{\mathbin{\scalebox{0.6}{$\Diamond$}}}%
  \def\bigcup{\mathop{\mathchoice{\raisebox{-0.3em}{\scalebox{1.7}{$\cup$}}}{\scalebox{1.25}{$\cup$}}{\cup}{\cup}}\displaylimits}%
  \def\bigcap{\mathop{\mathchoice{\raisebox{-0.3em}{\scalebox{1.7}{$\cap$}}}{\scalebox{1.25}{$\cap$}}{\cap}{\cap}}\displaylimits}%
  \def\lesseqqgtr{\mathrel{\lessgtr}}\def\bigoplus{\mathop{\oplus}\displaylimits}%
}
\providecommand{\ssi}{si et seulement si} % abréviation fréquente
\AtBeginDocument{\providecommand{\wideparen}{\overparen}} % arc de cercle

\begin{document}

\pagegarde{El deporte, el ocio y el bienestar (continuación) : expresión oral y escrita}{Espagnol}{1ère A}{Módulo 2 : Environnement, santé et bien-être}{EA08}{2 h}{Cette leçon mixte poursuit l'étude du sport, du loisir et du bien-être : à partir d'un dialogue authentique sur un match, elle développe le lexique sportif, la grammaire de la comparaison (comparatifs et superlatifs) et conduit à la rédaction d'un texte sur son sport préféré ainsi qu'à la traduction de phrases sur le sport.
\vspace{4pt}
\begin{multicols}{2}\small
\begin{itemize}
\item À la fin de cette leçon, je sais comprendre un dialogue sur un match, un entraînement ou un loisir
\item À la fin de cette leçon, je sais utiliser le lexique du sport, des loisirs et du bien-être en contexte
\item À la fin de cette leçon, je sais former et employer les comparatifs de supériorité, d'infériorité et d'égalité (más… que, menos… que, tan… como)
\item À la fin de cette leçon, je sais former et employer les superlatifs relatifs et absolus (el más…, el mejor, -ísimo)
\item À la fin de cette leçon, je sais comparer deux sports ou deux loisirs à l'oral et à l'écrit
\item À la fin de cette leçon, je sais rédiger un texte de 80 mots pour présenter mon sport préféré ou un champion
\end{itemize}\end{multicols}}

\begin{sommaire}
\begin{multicols}{2}
\begin{enumerate}
\item Texte support : un dialogue après le match
\item Lexique : sport, loisirs et bien-être
\item Grammaire 1 : les comparatifs
\item Grammaire 2 : les superlatifs
\item Traduction : phrases sur le sport (thème et version)
\item Production guidée : mon sport préféré ou mon champion (80 mots)
\item Activité d'intégration
\item Méthodes et exercices résolus
\item Exercices
\item Corrigés des exercices
\item Fiche bilan
\end{enumerate}
\end{multicols}
\end{sommaire}

\begin{objectifs}
\begin{itemize}
\item À la fin de cette leçon, je sais comprendre un dialogue sur un match, un entraînement ou un loisir
\item À la fin de cette leçon, je sais utiliser le lexique du sport, des loisirs et du bien-être en contexte
\item À la fin de cette leçon, je sais former et employer les comparatifs de supériorité, d'infériorité et d'égalité (más… que, menos… que, tan… como)
\item À la fin de cette leçon, je sais former et employer les superlatifs relatifs et absolus (el más…, el mejor, -ísimo)
\item À la fin de cette leçon, je sais comparer deux sports ou deux loisirs à l'oral et à l'écrit
\item À la fin de cette leçon, je sais rédiger un texte de 80 mots pour présenter mon sport préféré ou un champion
\item À la fin de cette leçon, je sais traduire des phrases sur le sport du français vers l'espagnol et inversement
\item À la fin de cette leçon, je sais éviter les pièges fréquents des francophones dans la comparaison
\end{itemize}
\end{objectifs}

\begin{prerequis}
\begin{itemize}
\item Les adjectifs qualificatifs et leur accord (Seconde)
\item Le présent de l'indicatif des verbes réguliers et de jugar, practicar, hacer, ir (Seconde)
\item Le lexique de base du sport vu en EA07 (deporte, equipo, partido, entrenamiento)
\item Les articles définis el, la, los, las et les possessifs
\item Les connecteurs simples : y, pero, porque, además
\end{itemize}
\end{prerequis}

\newpage

\coursec{Texte support : un dialogue après le match}

C'est samedi au stade Ahmadou Ahidjo de Yaoundé : les Lions Indomptables viennent de gagner et, dans les rues, tout le monde compare les joueurs. « Aboubakar est meilleur que l'attaquant d'hier », « ce match était moins intéressant que la finale de la CAN », « Eto'o reste le plus grand ». En espagnol, pour débattre ainsi du sport, des loisirs et du bien-être, il faut maîtriser les \cle{comparatifs} et les \cle{superlatifs}, et savoir raconter un match à l'oral comme à l'écrit. C'est exactement ce que cette leçon va vous permettre de faire.

\textbf{Problématique :}\par
Comment raconter un match et donner son opinion en espagnol en comparant les joueurs, les équipes et les performances ?

\courssub{Lecture globale du dialogue}

Avant de lire, observez la situation : deux lycéens camerounais, Andrés et Chantal, sortent du stade après un match de championnat scolaire. Ils commentent la rencontre, comparent les joueurs et ne sont pas toujours d'accord. Lisez le dialogue une première fois \emph{en silence}, puis une deuxième fois \emph{à voix haute}, sans vous arrêter sur les mots inconnus : le glossaire vous aidera ensuite.

\begin{textebox}[Después del partido]
\textbf{Andrés :} ¡Qué partido, Chantal! El equipo de nuestro barrio jugó mejor que el año pasado. El marcador final fue dos a uno: ¡ganamos!\par
\textbf{Chantal :} Sí, pero el portero rival fue tan bueno como el nuestro. Paró tres balones difíciles. Sin él, el resultado habría sido peor para su equipo.\par
\textbf{Andrés :} Para mí, el delantero Mbarga es el mejor jugador del campeonato: corre más que todos y marca goles en cada partido. Creo que va a jugar en un equipo profesional.\par
\textbf{Chantal :} No estoy de acuerdo: el capitán es menos rápido que él, pero es el más inteligente del campo. Siempre sabe dónde está el balón y ayuda a sus compañeros.\par
\textbf{Andrés :} Bueno, los dos son buenos. ¿Y la afición? ¡Increíble! Hoy había más gente que el sábado pasado. Todos cantaban y bailaban.\par
\textbf{Chantal :} Es verdad. Me parece que el ambiente fue mejor que en el partido contra el Lycée de Douala. Hubo casi un empate al final, pero nuestro entrenador cambió la táctica.\par
\textbf{Andrés :} El entrenamiento de esta semana fue más duro que el anterior, ¡pero valió la pena! El lunes volvemos a entrenar.\par
\textbf{Chantal :} ¡Perfecto! Yo también quiero entrenar más: el deporte es tan importante como los estudios para estar bien.
\end{textebox}

\textbf{Glossaire des mots nouveaux :}

\begin{center}
\begin{tabularx}{\linewidth}{|X|X|}\hline
\rowcolor{popblueL} \textbf{Español} & \textbf{Français} \\ \hline
el marcador & le score, le tableau d'affichage \\ \hline
el empate & le match nul \\ \hline
el portero & le gardien de but \\ \hline
el delantero & l'attaquant \\ \hline
el entrenador & l'entraîneur \\ \hline
la afición & les supporters, le public passionné \\ \hline
el campeonato & le championnat \\ \hline
duro, -a & difficile, dur \\ \hline
valer la pena & valoir la peine \\ \hline
el ambiente & l'ambiance \\ \hline
parar (un balón) & arrêter (un ballon) \\ \hline
\end{tabularx}
\end{center}

\begin{savaistu}[Portero ou arquero ?]
En Amérique latine, on dit souvent \esp{el arquero} pour « gardien de but », alors qu'en Espagne on dit \esp{el portero}. Au Bac, les deux sont acceptés. De même, les Latino-Américains disent \esp{el equipo} comme les Espagnols, mais prononcent souvent le \emph{c} et le \emph{z} comme un « s » (seseo) : \esp{afición} se dit « a-fi-sion » à Mexico et « a-fi-thion » à Madrid (ceceo).
\end{savaistu}

\courssub{Repérage des informations}

Après une première lecture, vous devez pouvoir répondre aux questions essentielles : \emph{qui ? où ? quoi ? quel résultat ?}

\begin{itemize}
\item \textbf{Les personnages :} Andrés et Chantal, deux lycéens qui reviennent du stade.
\item \textbf{Le lieu :} les abords du stade, juste après le match (\esp{después del partido}).
\item \textbf{Les équipes :} l'équipe du quartier (\esp{el equipo de nuestro barrio}) contre une équipe rivale ; on mentionne aussi un ancien match contre le Lycée de Douala.
\item \textbf{Le score :} \esp{dos a uno} (deux à un) ; l'équipe du quartier a gagné (\esp{¡ganamos!}).
\item \textbf{Les opinions :} Andrés admire le delantero Mbarga ; Chantal préfère le capitán, plus intelligent selon elle ; les deux s'accordent sur la bonne ambiance et l'importance du sport.
\end{itemize}

\begin{methode}[Comprendre un dialogue en quatre étapes]
\begin{enumerate}
\item \textbf{Identifier qui parle et où} : repérez les noms des personnages et les indices de lieu (\esp{después del partido}, \esp{el campo}, \esp{el estadio}).
\item \textbf{Repérer les faits} : score, événements, actions objectives (\esp{el marcador final fue dos a uno}).
\item \textbf{Distinguer faits et opinions} : une opinion est souvent introduite par \esp{me parece que}, \esp{creo que}, \esp{para mí}.
\item \textbf{Relever les mots de la comparaison} : soulignez \esp{más... que}, \esp{menos... que}, \esp{tan... como}, \esp{el mejor}, \esp{el más} : ils annoncent la grammaire de la leçon.
\end{enumerate}
\end{methode}

\courssub{Questions de compréhension}

\begin{exercice}[Vrai ou faux ? Justifiez par le texte]
\begin{enumerate}
\item El equipo del barrio perdió el partido.
\item El portero rival paró tres balones.
\item Chantal piensa que Mbarga es el mejor jugador.
\item Había menos gente que el sábado pasado.
\item El entrenamiento de esta semana fue más duro que el anterior.
\end{enumerate}
\end{exercice}

\begin{exercice}[Questions ouvertes]
\begin{enumerate}
\item ¿Cuál fue el marcador final del partido?
\item ¿Por qué Chantal prefiere al capitán?
\item ¿Qué hizo el entrenador al final del partido?
\item ¿Qué opina Chantal del deporte y de los estudios?
\end{enumerate}
\end{exercice}

\begin{corrige}[Exercices 1 et 2]
\textbf{Vrai ou faux :}
\begin{enumerate}
\item \textbf{Falso} : \esp{El marcador final fue dos a uno: ¡ganamos!} — l'équipe du quartier a gagné.
\item \textbf{Verdadero} : \esp{Paró tres balones difíciles.}
\item \textbf{Falso} : c'est Andrés qui admire Mbarga ; Chantal dit \esp{No estoy de acuerdo} et préfère le capitán.
\item \textbf{Falso} : \esp{Hoy había más gente que el sábado pasado.}
\item \textbf{Verdadero} : \esp{El entrenamiento de esta semana fue más duro que el anterior.}
\end{enumerate}
\textbf{Questions ouvertes :}
\begin{enumerate}
\item \esp{El marcador final fue dos a uno.}
\item \esp{Porque es el más inteligente del campo: siempre sabe dónde está el balón y ayuda a sus compañeros.}
\item \esp{Cambió la táctica.}
\item \esp{Piensa que el deporte es tan importante como los estudios para estar bien.}
\end{enumerate}
\end{corrige}

\courssub{Relevé des comparaisons et des expressions d'opinion}

Relevez dans le dialogue toutes les phrases qui comparent. Observez-les avant de découvrir la règle dans la section grammaire :

\begin{center}
\begin{tabularx}{\linewidth}{|X|X|X|}\hline
\rowcolor{popblueL} \textbf{Phrase du dialogue} & \textbf{Traduction} & \textbf{Structure} \\ \hline
\esp{jugó mejor que el año pasado} & a mieux joué que l'année dernière & mejor + que \\ \hline
\esp{tan bueno como el nuestro} & aussi bon que le nôtre & tan + adj. + como \\ \hline
\esp{corre más que todos} & il court plus que tout le monde & más + adv. + que \\ \hline
\esp{el mejor jugador del campeonato} & le meilleur joueur du championnat & el mejor + sustantivo + de \\ \hline
\esp{menos rápido que él} & moins rapide que lui & menos + adj. + que \\ \hline
\esp{el más inteligente del campo} & le plus intelligent du terrain & el más + adj. + de \\ \hline
\esp{más gente que el sábado pasado} & plus de monde que samedi dernier & más + sustantivo + que \\ \hline
\esp{más duro que el anterior} & plus dur que le précédent & más + adj. + que \\ \hline
\esp{tan importante como los estudios} & aussi important que les études & tan + adj. + como \\ \hline
\end{tabularx}
\end{center}

\begin{exemplebox}[Deux exemples à retenir]
\esp{El portero rival fue tan bueno como el nuestro.} « Le gardien adverse a été aussi bon que le nôtre. »\par
\esp{Corre más que todos.} « Il court plus que tout le monde. »
\end{exemplebox}

\begin{aretenir}[Exprimer son opinion]
Pour donner son avis dans un débat sportif, utilisez : \esp{me parece que...} (il me semble que...), \esp{creo que...} (je pense que...), \esp{para mí...} (pour moi...), \esp{estoy de acuerdo} (je suis d'accord) et \esp{no estoy de acuerdo} (je ne suis pas d'accord). Ces expressions transforment un fait en opinion personnelle.
\end{aretenir}

\begin{attention}[Piège du francophone]
Ne traduisez jamais « je suis d'accord » par \esp{soy de acuerdo} : c'est une erreur de calque du français. On dit \esp{\textbf{estoy} de acuerdo}, avec le verbe \esp{estar}. De même, « pour moi » se dit \esp{para mí} (avec accent sur \emph{mí}), et non \esp{para yo}.
\end{attention}

\courssub{Carte mentale du vocabulaire du match}

\begin{popfigure}
\begin{tikzpicture}[mindmap, grow cyclic, every node/.style={concept, align=center, font=\small}, root concept/.append style={concept color=popblue, font=\bfseries}, level 1/.append style={level distance=3.2cm, sibling angle=72}]
\node[root concept]{El partido}
  child[concept color=poporange] { node {equipos} child[concept color=poporangeL]{ node {el portero} } child[concept color=poporangeL]{ node {el delantero} } }
  child[concept color=popgreen] { node {marcador} child[concept color=popgreenL]{ node {dos a uno} } child[concept color=popgreenL]{ node {el empate} } }
  child[concept color=poppurple] { node {jugadores} child[concept color=poppurpleL]{ node {el capitán} } child[concept color=poppurpleL]{ node {el entrenador} } }
  child[concept color=poppink] { node {opiniones} child[concept color=poppinkL]{ node {me parece que} } child[concept color=poppinkL]{ node {creo que} } }
  child[concept color=popgold] { node {afición} child[concept color=popgoldL]{ node {cantar} } child[concept color=popgoldL]{ node {bailar} } };
\end{tikzpicture}
\legende{Carte mentale : le vocabulaire du match de football}
\end{popfigure}

\courssub{Exercice résolu et jeu de rôles}

\begin{exoresolu}[Repérage dans le texte]
Énoncé : relevez dans le dialogue une phrase de fait et une phrase d'opinion, puis traduisez-les.
\tcblower
Solution : \textbf{Fait :} \esp{El marcador final fue dos a uno.} « Le score final a été de deux à un. » C'est une information objective, vérifiable. \textbf{Opinion :} \esp{Para mí, el delantero Mbarga es el mejor jugador del campeonato.} « Pour moi, l'attaquant Mbarga est le meilleur joueur du championnat. » L'expression \esp{para mí} signale qu'il s'agit du point de vue personnel d'Andrés, que Chantal conteste d'ailleurs : \esp{No estoy de acuerdo}.
\end{exoresolu}

\begin{experience}[Jeu de rôles : lecture dramatisée]
Par binômes, lisez le dialogue à voix haute : un élève joue Andrés, l'autre Chantal. Mettez le ton : enthousiasme pour \esp{¡Qué partido!}, désaccord ferme mais poli pour \esp{No estoy de acuerdo}. Puis échangez les rôles. Variante : remplacez Mbarga par un joueur de votre choix (Aboubakar, Eto'o, un joueur de votre quartier) et improvisez la fin du dialogue en utilisant au moins deux comparaisons et une expression d'opinion.
\end{experience}

\begin{aretenir}[Ce qu'il faut retenir de cette section]
Un dialogue sportif mélange \textbf{faits} (score, actions) et \textbf{opinions} (\esp{me parece que}, \esp{creo que}, \esp{para mí}). Pour comparer, l'espagnol utilise \esp{más... que} (plus... que), \esp{menos... que} (moins... que), \esp{tan... como} (aussi... que) et les superlatifs \esp{el mejor} (le meilleur), \esp{el más... de} (le plus... de). Ces structures, observées ici dans le dialogue, seront étudiées en détail dans les sections de grammaire.
\end{aretenir}

\coursec{Lexique : sport, loisirs et bien-être}

Dans la section précédente, nous avons lu et commenté le dialogue entre Carlos et André après le match de football de leur lycée. Ce dialogue regorgeait de mots nouveaux : \esp{el partido}, \esp{el marcador}, \esp{empatar}, \esp{el entrenador}... Pour pouvoir parler librement d'un match, d'un entraînement ou d'un loisir, il faut maintenant organiser et mémoriser ce vocabulaire. C'est l'objectif de cette deuxième section : construire, champ par champ, un lexique riche et précis sur le sport, les loisirs et le bien-être, que tu réutiliseras dans la traduction (section 5) et dans ta rédaction (section 6).

\begin{definition}[El ocio]
\esp{El ocio} (prononcé « o-sio ») désigne le \cle{temps libre} et l'ensemble des \cle{activités de loisir} que l'on pratique pour se détendre : le sport, la musique, le cinéma, les sorties entre amis. Attention : \esp{ocio} est un mot positif en espagnol ; il ne faut pas le confondre avec \esp{la pereza} (la paresse), qui a un sens négatif. \esp{El ocio es necesario para la salud.} « Le loisir est nécessaire pour la santé. »
\end{definition}

\courssub{Champ 1 : les sports}

Observons d'abord cette phrase tirée de notre dialogue : \esp{El fútbol es muy popular en Camerún, pero también me gusta el baloncesto.} « Le football est très populaire au Cameroun, mais j'aime aussi le basket-ball. » On remarque que les noms de sports sont précédés de l'article \esp{el} : c'est une règle générale.

\begin{center}
\begin{tabularx}{\linewidth}{|X|X|X|}
\hline
\rowcolor{popblueL} Español & Français & Exemple \\
\hline
el fútbol & le football & \esp{Juego al fútbol los sábados.} (Je joue au foot le samedi.) \\
\hline
el baloncesto & le basket-ball & \esp{El baloncesto es muy rápido.} \\
\hline
el vóleibol & le volley-ball & \esp{Las chicas juegan al vóleibol.} \\
\hline
el atletismo & l'athlétisme & \esp{El atletismo exige mucho esfuerzo.} \\
\hline
la natación & la natation & \esp{La natación es buena para la salud.} \\
\hline
el boxeo & la boxe & \esp{El boxeo es un deporte fuerte.} \\
\hline
el ciclismo & le cyclisme & \esp{El ciclismo es popular en España.} \\
\hline
\end{tabularx}
\end{center}

\begin{propriete}[Genre des noms de sports et verbes associés]
Presque tous les noms de sports sont \cle{masculins} et s'emploient avec l'article défini : \esp{el fútbol}, \esp{el baloncesto}, \esp{el ciclismo}. Exception fréquente : \esp{la natación} (féminin, comme tous les mots en \esp{-ción}). Deux constructions à connaître :
\begin{itemize}
\item Avec \esp{jugar} (jouer) : \esp{jugar a + article} ; devant \esp{el}, la contraction est obligatoire : \esp{jugar \textbf{al} fútbol}, \esp{jugar \textbf{al} vóleibol}.
\item Avec \esp{practicar} ou \esp{hacer} (pratiquer, faire) : \esp{practicar el judo}, \esp{hacer atletismo}, \esp{hacer natación}.
\end{itemize}
\esp{Los domingos juego al fútbol y practico la natación.} « Le dimanche, je joue au football et je fais de la natation. »
\end{propriete}

\begin{attention}[Le faux ami « foot »]
En espagnol, \esp{el fútbol} (prononcé « fut-bol ») signifie « le football » (le sport que joue Samuel Eto'o). Le mot anglais \emph{foot} n'existe pas en espagnol : ne traduis jamais « je joue au foot » par \esp{juego al foot}, mais par \esp{juego al fútbol}. Et si tu parles du sport américain, il faut préciser : \esp{el fútbol americano} (le football américain). De même, \esp{el balón} est le ballon, \esp{el baloncesto} le basket : ne confonds pas \esp{balón} et \esp{ballet} !
\end{attention}

\courssub{Champ 2 : les acteurs et les lieux}

Pour raconter un match, il faut nommer les personnes et les lieux. Observons : \esp{El portero paró el balón y el árbitro pitó el final del partido.} « Le gardien a arrêté le ballon et l'arbitre a sifflé la fin du match. »

\begin{center}
\begin{tabularx}{\linewidth}{|X|X|}
\hline
\rowcolor{popgreenL} Personas & Lugares \\
\hline
el/la jugador(a) : le/la joueur(-euse) & el estadio : le stade \\
\hline
el portero : le gardien de but & la cancha / el campo : le terrain (de sport) \\
\hline
el entrenador / la entrenadora : l'entraîneur(-euse) & el gimnasio : le gymnase, la salle de sport \\
\hline
el árbitro / la árbitra : l'arbitre & la piscina : la piscine \\
\hline
el aficionado / la aficionada : le/la supporter, passionné(e) & el campo de fútbol : le terrain de foot \\
\hline
el equipo : l'équipe & la pista : la piste (d'athlétisme) \\
\hline
\end{tabularx}
\end{center}

Remarque la formation du féminin : \esp{el jugador} devient \esp{la jugadora}, comme \esp{el entrenador} / \esp{la entrenadora}. En revanche, \esp{el árbitro} / \esp{la árbitra} (accent sur le \emph{á}) et \esp{el aficionado} / \esp{la aficionada} suivent la règle normale en \esp{-o}/\esp{-a}. Note : \esp{la cancha} est très utilisé en Amérique latine et en Guinée équatoriale ; en Espagne on dit souvent \esp{el campo}.

\begin{exemplebox}[Exemples traduits]
\esp{El entrenador habla con los jugadores en el vestuario.} « L'entraîneur parle avec les joueurs dans le vestiaire. »

\esp{Hay muchos aficionados en el estadio de Yaoundé.} « Il y a beaucoup de supporters dans le stade de Yaoundé. »

\esp{Mi hermana es una jugadora de vóleibol muy fuerte.} « Ma sœur est une joueuse de volley très forte. »
\end{exemplebox}

\courssub{Champ 3 : le match et l'entraînement}

Voici le cœur du vocabulaire du dialogue de la section 1. Retenons les noms et les verbes essentiels :

\begin{center}
\begin{tabularx}{\linewidth}{|X|X|}
\hline
\rowcolor{poporangeL} Español & Français \\
\hline
el partido & le match \\
\hline
el marcador & le score (tableau d'affichage) \\
\hline
el empate & le match nul \\
\hline
ganar & gagner \\
\hline
perder & perdre \\
\hline
empatar & faire match nul, égaliser \\
\hline
marcar un gol & marquer un but \\
\hline
entrenarse & s'entraîner \\
\hline
el campeonato & le championnat \\
\hline
la liga & la ligue, le championnat \\
\hline
\end{tabularx}
\end{center}

\begin{exemplebox}[Exemples traduits]
\esp{Nuestro equipo empató dos a dos.} « Notre équipe a fait match nul deux partout. »

\esp{El delantero marcó un gol en el último minuto.} « L'attaquant a marqué un but à la dernière minute. »

\esp{Los jugadores se entrenan todos los días en el gimnasio.} « Les joueurs s'entraînent tous les jours à la salle de sport. »

\esp{Camerún ganó el campeonato africano.} « Le Cameroun a gagné le championnat africain. »
\end{exemplebox}

\begin{attention}[Pièges de traduction]
\begin{itemize}
\item « Gagner un match » se dit \esp{ganar un partido}, jamais \esp{ganar un juego} (\esp{el juego} = le jeu).
\item \esp{Perder} est un verbe irrégulier à changement radical \esp{e > ie} : \esp{pierdo, pierdes, pierde, perdemos, perdéis, pierden}. \esp{Mi equipo nunca pierde en casa.} « Mon équipe ne perd jamais à domicile. »
\item \esp{Entrenarse} est pronominal : n'oublie pas le pronom réfléchi : \esp{me entreno, te entrenas, se entrena...}
\item Le score se lit avec \esp{a} : \esp{tres a uno} = « trois à un ».
\end{itemize}
\end{attention}

\courssub{Champ 4 : les loisirs}

Le sport n'est pas le seul loisir ! Voici le vocabulaire des activités du temps libre :

\begin{center}
\begin{tabularx}{\linewidth}{|X|X|}
\hline
\rowcolor{poppurpleL} Español & Français \\
\hline
la música & la musique \\
\hline
el cine & le cinéma \\
\hline
la lectura & la lecture \\
\hline
los videojuegos & les jeux vidéo \\
\hline
pasear / dar un paseo & se promener / faire une promenade \\
\hline
bailar & danser \\
\hline
salir con amigos & sortir avec des amis \\
\hline
ver la televisión & regarder la télévision \\
\hline
\end{tabularx}
\end{center}

\esp{Los fines de semana, me gusta salir con amigos y bailar.} « Les week-ends, j'aime sortir avec des amis et danser. » — \esp{Mi hermano prefiere los videojuegos, pero yo prefiero la lectura.} « Mon frère préfère les jeux vidéo, mais moi je préfère la lecture. »

\courssub{Champ 5 : le bien-être}

Le sport et les loisirs ont un but : le bien-être, \esp{el bienestar}. Ce dernier champ lexical relie notre leçon au module « santé et bien-être ».

\begin{center}
\begin{tabularx}{\linewidth}{|X|X|}
\hline
\rowcolor{popgoldL} Español & Français \\
\hline
la salud & la santé \\
\hline
estar en forma & être en forme \\
\hline
descansar & se reposer \\
\hline
relajarse & se détendre, se relaxer \\
\hline
una vida sana & une vie saine \\
\hline
hacer ejercicio & faire de l'exercice \\
\hline
dormir bien & bien dormir \\
\hline
\end{tabularx}
\end{center}

\begin{exemplebox}[Exemples traduits]
\esp{Hago ejercicio para estar en forma.} « Je fais de l'exercice pour rester en forme. »

\esp{Después del partido, los jugadores descansan y se relajan.} « Après le match, les joueurs se reposent et se détendent. »

\esp{Una vida sana es importante para la salud.} « Une vie saine est importante pour la santé. »
\end{exemplebox}

\courssub{Adjectifs utiles pour comparer}

Pour préparer la grammaire des sections 3 et 4 (comparatifs et superlatifs), mémorise ces paires d'adjectifs contraires :

\begin{center}
\begin{tabularx}{\linewidth}{|X|X|}
\hline
\rowcolor{poppinkL} Español & Français \\
\hline
rápido / lento & rapide / lent \\
\hline
fuerte / débil & fort / faible \\
\hline
fácil / difícil & facile / difficile \\
\hline
interesante / aburrido & intéressant / ennuyeux \\
\hline
divertido / aburrido & amusant / ennuyeux \\
\hline
sano & sain \\
\hline
popular & populaire \\
\hline
\end{tabularx}
\end{center}

\esp{El atletismo es más rápido que el ciclismo.} « L'athlétisme est plus rapide que le cyclisme. » — \esp{La lectura es menos aburrida que los videojuegos.} « La lecture est moins ennuyeuse que les jeux vidéo. »

\courssub{Expressions d'opinion et de préférence}

Pour exprimer tes goûts, trois formules à connaître par cœur :

\begin{itemize}
\item \esp{Me gusta más...} : « J'aime davantage... » → \esp{Me gusta más el baloncesto que el boxeo.} « J'aime plus le basket que la boxe. »
\item \esp{Prefiero...} : « Je préfère... » (verbe irrégulier, radical \esp{e > ie} : \esp{prefiero, prefieres, prefiere, preferimos, preferís, prefieren}) → \esp{Prefiero la natación.} « Je préfère la natation. »
\item \esp{Mi deporte favorito es...} : « Mon sport préféré est... » → \esp{Mi deporte favorito es el fútbol.} « Mon sport préféré est le football. »
\end{itemize}

\begin{popfigure}
\begin{center}
\begin{tikzpicture}
\node[draw=popblue, fill=popblueL, rounded corners, align=center, minimum width=3.4cm, minimum height=2.8cm] (dep) at (0,0) {\textbf{Deportes}\\[2pt] el fútbol\\ el baloncesto\\ la natación\\ el atletismo\\ el ciclismo};
\node[draw=poppurple, fill=poppurpleL, rounded corners, align=center, minimum width=3.4cm, minimum height=2.8cm] (ocio) at (5.2,0) {\textbf{Ocio}\\[2pt] la música\\ el cine\\ la lectura\\ bailar\\ salir con amigos};
\node[draw=popgreen, fill=popgreenL, rounded corners, align=center, minimum width=3.4cm, minimum height=2.8cm] (bien) at (10.4,0) {\textbf{Bienestar}\\[2pt] la salud\\ estar en forma\\ descansar\\ relajarse\\ hacer ejercicio};
\node[draw=popdark, fill=popcream, rounded corners, align=center, font=\bfseries] (titre) at (5.2,3.2) {El deporte, el ocio y el bienestar};
\draw[-{Stealth}, thick, popmuted] (titre.south) -- (dep.north);
\draw[-{Stealth}, thick, popmuted] (titre.south) -- (ocio.north);
\draw[-{Stealth}, thick, popmuted] (titre.south) -- (bien.north);
\end{tikzpicture}
\end{center}
\legende{Carte mentale : trois champs lexicaux de la leçon, cinq mots chacun.}
\end{popfigure}

\begin{exoresolu}[Classer et compléter]
Énoncé. 1) Classe les mots suivants dans les trois colonnes \esp{Deportes / Ocio / Bienestar} : \esp{el boxeo, la lectura, descansar, el vóleibol, la salud, bailar, el ciclismo, el cine, relajarse, la música}. 2) Traduis en espagnol : « Je joue au volley-ball le samedi et je fais de l'exercice pour être en forme. »
\tcblower
Solution détaillée. 1) \esp{Deportes} : el boxeo, el vóleibol, el ciclismo. \esp{Ocio} : la lectura, bailar, el cine, la música. \esp{Bienestar} : descansar, la salud, relajarse. 2) On utilise \esp{jugar a + el = al} pour le sport avec ballon, et \esp{hacer ejercicio} pour « faire de l'exercice » ; « pour » de but se traduit par \esp{para} : \esp{Los sábados juego al vóleibol y hago ejercicio para estar en forma.}
\end{exoresolu}

\begin{exercice}[À toi de jouer]
1) Complète avec l'article ou la contraction correcte : \esp{Juego ... fútbol ; practico ... judo ; ... natación es buena para ... salud.} 2) Traduis : « Mon équipe a fait match nul deux partout hier au stade. » 3) Écris trois phrases avec \esp{me gusta más}, \esp{prefiero} et \esp{mi deporte favorito es}.
\end{exercice}

\begin{corrige}[Exercice « À toi de jouer »]
1) \esp{Juego \textbf{al} fútbol ; practico \textbf{el} judo ; \textbf{la} natación es buena para \textbf{la} salud.} 2) \esp{Mi equipo empató dos a dos ayer en el estadio.} 3) Réponses possibles : \esp{Me gusta más el ciclismo que el boxeo. Prefiero la lectura. Mi deporte favorito es el baloncesto.}
\end{corrige}

\begin{savaistu}[Le fútbol en Guinée équatoriale]
Le Cameroun a un voisin hispanophone : la Guinée équatoriale, ancienne colonie espagnole, où l'espagnol est langue officielle. Là-bas comme ici, \esp{el fútbol} est une passion nationale : l'équipe nationale de Guinée équatoriale est surnommée \esp{Nzalang Nacional}. Quand les Lions Indomptables rencontrent leurs voisins, les commentateurs espagnols parlent de \esp{un partido muy emocionante} : « un match très émouvant » !
\end{savaistu}

\begin{aretenir}[L'essentiel du lexique]
Retiens les cinq champs : les \cle{sports} (masculins, avec article : \esp{el fútbol}, sauf \esp{la natación}), les \cle{acteurs et lieux} (\esp{el jugador, el árbitro, el estadio}), le \cle{match} (\esp{ganar, perder, empatar, marcar un gol}), les \cle{loisirs} (\esp{el cine, bailar, salir con amigos}) et le \cle{bien-être} (\esp{la salud, estar en forma, hacer ejercicio}). Avec \esp{jugar}, n'oublie jamais la contraction : \esp{jugar \textbf{al} fútbol}. Ce vocabulaire te servira pour comparer les sports (sections 3 et 4), traduire (section 5) et rédiger ton texte sur ton sport préféré (section 6).
\end{aretenir}

\coursec{Grammaire 1 : les comparatifs}

Dans la section précédente, nous avons enrichi notre vocabulaire du sport, des loisirs et du bien-être : \esp{el entrenamiento}, \esp{la afición}, \esp{el marcador}, \esp{estar en forma}… Or, dès que l'on parle de sport, on compare sans cesse : une équipe joue \emph{mieux} qu'une autre, un champion court \emph{plus vite} que son rival, la natation est \emph{aussi saine que} l'athlétisme. C'est exactement ce que faisaient les deux amis du dialogue après le match. Reprenons leurs phrases pour observer comment l'espagnol construit la comparaison.

\courssub{Observation : la comparaison dans le dialogue}

\begin{exemplebox}[Phrases relevées dans le dialogue]
\begin{itemize}
\item \esp{El equipo jugó mejor que el año pasado.} « L'équipe a joué mieux que l'année dernière. »
\item \esp{Su delantero es tan bueno como el nuestro.} « Leur attaquant est aussi bon que le nôtre. »
\item \esp{Nuestro portero es menos rápido que él.} « Notre gardien est moins rapide que lui. »
\item \esp{El fútbol es más popular que el boxeo.} « Le football est plus populaire que la boxe. »
\end{itemize}
\end{exemplebox}

Observons attentivement. Chaque phrase contient \cle{trois blocs} : d'abord un mot qui annonce le sens de la comparaison (\esp{más}, \esp{menos}, \esp{tan}, ou une forme spéciale comme \esp{mejor}), ensuite un adjectif ou un adverbe (\esp{popular}, \esp{bueno}, \esp{rápido}), enfin un petit mot de liaison (\esp{que} ou \esp{como}) qui introduit le second terme de la comparaison. C'est exactement la même architecture qu'en français : « plus… que », « moins… que », « aussi… que ».

\begin{propriete}[Structure générale de la comparaison]
La comparaison se construit en \cle{trois blocs} :
\begin{center}
[más / menos / tan] + [adjectif, adverbe ou nom] + [que / como]
\end{center}
Deux règles fondamentales :
\begin{enumerate}
\item L'adjectif s'\cle{accorde en genre et en nombre avec le premier terme comparé} : \esp{María es más rápida que Juan.} « María est plus rapide que Juan. » (\esp{rápida}, féminin singulier, comme \esp{María})
\item Après \esp{más} et \esp{menos}, on emploie toujours \esp{que} (jamais \esp{como}) ; après \esp{tan}, on emploie toujours \esp{como}.
\end{enumerate}
\end{propriete}

\begin{popfigure}
\begin{tikzpicture}[node distance=0.6cm, every node/.style={align=center, text width=4.2cm, rounded corners}]
\node[draw=popgreen, fill=popgreenL, align=center] (sup){\textbf{Supériorité}\\ \esp{más\ldots que} $\uparrow$\\[2pt] \footnotesize\esp{El fútbol es más popular que el boxeo.}};
\node[draw=popblue, fill=popblueL, right=0.8cm of sup, align=center] (ega){\textbf{Égalité}\\ \esp{tan\ldots como} $=$\\[2pt] \footnotesize\esp{La natación es tan sana como el atletismo.}};
\node[draw=poporange, fill=poporangeL, right=0.8cm of ega, align=center] (inf){\textbf{Infériorité}\\ \esp{menos\ldots que} $\downarrow$\\[2pt] \footnotesize\esp{El ciclismo es menos peligroso que el boxeo.}};
\end{tikzpicture}
\legende{Les trois degrés de la comparaison en espagnol}
\end{popfigure}

\courssub{Le comparatif de supériorité : más\ldots que}

\begin{propriete}[Comparatif de supériorité]
Forme : \cle{más} + adjectif / adverbe / nom + \cle{que}.
\begin{itemize}
\item Avec un adjectif : \esp{El baloncesto es más espectacular que el golf.} « Le basket est plus spectaculaire que le golf. »
\item Avec un adverbe : \esp{Corro más rápido que mi hermano.} « Je cours plus vite que mon frère. »
\item Avec un nom : \esp{En Douala hay más aficionados que en Bafoussam.} « À Douala, il y a plus de supporters qu'à Bafoussam. »
\end{itemize}
Remarque : \esp{más} est invariable ; c'est l'adjectif qui s'accorde : \esp{La piscina es más grande que el gimnasio.} « La piscine est plus grande que le gymnase. »
\end{propriete}

\courssub{Le comparatif d'infériorité : menos\ldots que}

\begin{propriete}[Comparatif d'infériorité]
Forme : \cle{menos} + adjectif / adverbe / nom + \cle{que}.
\begin{itemize}
\item \esp{El ciclismo es menos peligroso que el boxeo.} « Le cyclisme est moins dangereux que la boxe. »
\item \esp{Entreno menos a menudo que antes.} « Je m'entraîne moins souvent qu'avant. »
\item \esp{Este estadio tiene menos asientos que el otro.} « Ce stade a moins de places que l'autre. »
\end{itemize}
\end{propriete}

\courssub{Le comparatif d'égalité : tan\ldots como et tanto\ldots como}

\begin{propriete}[Comparatif d'égalité]
Deux constructions selon la nature du mot comparé :
\begin{enumerate}
\item \cle{tan} + adjectif ou adverbe + \cle{como} : \esp{La natación es tan sana como el atletismo.} « La natation est aussi saine que l'athlétisme. » ; \esp{Juega tan bien como su hermana.} « Il joue aussi bien que sa sœur. »
\item \cle{tanto, tanta, tantos, tantas} + nom + \cle{como} : \esp{tanto} s'accorde en genre et en nombre avec le nom : \esp{Hay tantos aficionados como jugadores.} « Il y a autant de supporters que de joueurs. » ; \esp{Hago tanto deporte como tú.} « Je fais autant de sport que toi. » ; \esp{Gana tantas medallas como su entrenador en su juventud.} « Elle gagne autant de médailles que son entraîneur dans sa jeunesse. »
\end{enumerate}
\end{propriete}

\begin{attention}[tan ou tanto ?]
\esp{tan} ne s'emploie qu'avec les \cle{adjectifs} et les \cle{adverbes} : \esp{tan alto como}, \esp{tan rápidamente como}. Devant un \cle{nom}, on utilise \esp{tanto}, qui s'accorde : \esp{tanto tiempo como} « autant de temps que », \esp{tantas veces como} « autant de fois que ». Ne dites jamais \esp{tan tiempo como}.
\end{attention}

\courssub{Les comparatifs irréguliers : mejor, peor, mayor, menor}

Comme le français dit « meilleur » et non « plus bon », l'espagnol possède quatre formes irrégulières à connaître par cœur :

\begin{center}
\begin{tabularx}{\linewidth}{|X|X|X|X|}
\hline
\rowcolor{popblueL} Adjectif & Comparatif & Exemple & Traduction \\
\hline
\esp{bueno} (bon) & \esp{mejor} & \esp{El tenis es mejor que el golf.} & Le tennis est meilleur que le golf. \\
\hline
\esp{malo} (mauvais) & \esp{peor} & \esp{Su defensa es peor que la nuestra.} & Leur défense est pire que la nôtre. \\
\hline
\esp{grande} (grand) & \esp{mayor} & \esp{Mi hermano es mayor que yo.} & Mon frère est plus âgé que moi. \\
\hline
\esp{pequeño} (petit) & \esp{menor} & \esp{Mi hermana es menor que yo.} & Ma sœur est plus jeune que moi. \\
\hline
\end{tabularx}
\end{center}

Précisions d'emploi :
\begin{itemize}
\item \esp{mejor} et \esp{peor} sont invariables en genre mais prennent la marque du pluriel : \esp{Estas zapatillas son mejores que esas.} « Ces chaussures sont meilleures que celles-là. »
\item \esp{mayor} et \esp{menor} désignent surtout l'\cle{âge} (ou l'importance/la hiérarchie) : \esp{Pedro es mayor que Ana.} Pour la taille physique, on garde la forme régulière : \esp{El estadio es más grande que el gimnasio.} « Le stade est plus grand que le gymnase. »
\item L'adverbe \esp{bien} a aussi son irrégulier : \esp{juega mejor que antes} « il joue mieux qu'avant » (et non \esp{más bien}).
\end{itemize}

\begin{attention}[Ne jamais dire « más bueno »]
Erreur très fréquente des francophones : on ne dit jamais \esp{más bueno}, on dit \cle{mejor} ; de même \esp{más malo} devient \cle{peor}. De la même façon, pour l'âge, « plus âgé que » se traduit \esp{mayor que}, et non \esp{más viejo que} dans un contexte soigné.
\end{attention}

\courssub{que, como ou de ? Les petits mots de la comparaison}

\begin{propriete}[Le mot de liaison après más / menos]
\begin{enumerate}
\item Devant un nom, un pronom ou un infinitif : \cle{que}. \esp{Tengo más camisetas que tú.} « J'ai plus de maillots que toi. »
\item Devant un \cle{chiffre} ou une quantité précise : \cle{de}. \esp{Hay más de veinte jugadores en el campo.} « Il y a plus de vingt joueurs sur le terrain. » ; \esp{Tengo más de diez camisetas.} « J'ai plus de dix maillots. »
\item Après \esp{tan} : toujours \cle{como}. \esp{Es tan alto como su padre.} « Il est aussi grand que son père. »
\end{enumerate}
\end{propriete}

\begin{center}
\begin{tabularx}{\linewidth}{|X|X|X|}
\hline
\rowcolor{popblueL} Français & Español & Remarque \\
\hline
plus… que & \esp{más… que} & adjectif accordé avec le 1er terme \\
\hline
moins… que & \esp{menos… que} & jamais \esp{como} après \esp{menos} \\
\hline
aussi… que & \esp{tan… como} & avec adjectifs et adverbes \\
\hline
autant de… que & \esp{tanto, tanta, tantos, tantas… como} & accord avec le nom \\
\hline
plus de (+ chiffre) & \esp{más de} & \esp{más de veinte jugadores} \\
\hline
meilleur que & \esp{mejor que} & forme irrégulière de \esp{bueno} \\
\hline
pire que & \esp{peor que} & forme irrégulière de \esp{malo} \\
\hline
plus âgé / plus jeune & \esp{mayor / menor} & pour l'âge \\
\hline
\end{tabularx}
\end{center}

\begin{methode}[Traduire « plus… que » en espagnol]
\begin{enumerate}
\item Repérer l'adjectif ou l'adverbe de la phrase française.
\item Vérifier s'il s'agit de \emph{bon}, \emph{mauvais}, \emph{grand} (âge) ou \emph{petit} (âge) : si oui, utiliser la forme irrégulière \esp{mejor}, \esp{peor}, \esp{mayor}, \esp{menor}.
\item Sinon, appliquer le schéma \esp{más} + adjectif + \esp{que}.
\item Accorder l'adjectif avec le premier terme comparé (genre et nombre).
\item Devant un chiffre, remplacer \esp{que} par \esp{de}.
\end{enumerate}
Exemple : « Mon équipe est meilleure que la tienne. » → adjectif « meilleure » = irrégulier → \esp{Mi equipo es mejor que el tuyo.}
\end{methode}

\begin{exoresolu}[Traduction guidée]
Énoncé : traduis en espagnol les phrases suivantes. 1) « Le basket est plus spectaculaire que le golf. » 2) « Il y a autant de supporters que de joueurs. » 3) « J'ai plus de dix maillots. » 4) « María court plus vite que Juan. » 5) « Ce champion est meilleur que l'autre. »
\tcblower
Solution détaillée :
\begin{enumerate}
\item Adjectif « spectaculaire » : régulier → \esp{El baloncesto es más espectacular que el golf.}
\item « Autant de + nom » → \esp{tanto} accordé avec \esp{aficionados} (masculin pluriel) → \esp{Hay tantos aficionados como jugadores.}
\item Devant un chiffre, \esp{más de} → \esp{Tengo más de diez camisetas.}
\item Accord de l'adverbe : invariable, mais attention à l'ordre → \esp{María corre más rápido que Juan.}
\item « Meilleur » = irrégulier → \esp{Este campeón es mejor que el otro.} (jamais \esp{más bueno}).
\end{enumerate}
\end{exoresolu}

\begin{exemplebox}[Pour s'entraîner : exemples traduits]
\begin{itemize}
\item \esp{El atletismo es más exigente que la natación.} « L'athlétisme est plus exigeant que la natation. »
\item \esp{Las campeonas son tan fuertes como los campeones.} « Les championnes sont aussi fortes que les champions. »
\item \esp{Este año entrenamos menos horas que el año pasado.} « Cette année, nous nous entraînons moins d'heures que l'année dernière. »
\item \esp{El partido fue peor que el anterior.} « Le match a été pire que le précédent. »
\item \esp{En el estadio de Yaoundé hay más de mil aficionados.} « Dans le stade de Yaoundé, il y a plus de mille supporters. »
\item \esp{Mi prima es menor que yo, pero juega mejor.} « Ma cousine est plus jeune que moi, mais elle joue mieux. »
\end{itemize}
\end{exemplebox}

\begin{aretenir}[L'essentiel des comparatifs]
\begin{itemize}
\item Trois blocs : \esp{más / menos / tan} + adjectif + \esp{que / como}.
\item L'adjectif s'accorde avec le premier terme : \esp{María es más rápida que Juan.}
\item Irréguliers : \esp{bueno → mejor}, \esp{malo → peor}, \esp{grande → mayor} (âge), \esp{pequeño → menor} (âge).
\item Devant un chiffre : \esp{más de veinte}.
\item Devant un nom : \esp{tanto, tanta, tantos, tantas\ldots como}.
\end{itemize}
\end{aretenir}

\begin{exercice}[À toi de comparer]
\begin{enumerate}
\item Traduis : « La natation est moins dangereuse que la boxe. » ; « Mon frère est plus âgé que moi. » ; « Il y a plus de trente équipes dans le tournoi. » ; « Elle nage aussi bien que sa sœur. » ; « Ce match est pire que l'autre. »
\item Complète avec \esp{que}, \esp{como} ou \esp{de} : \esp{más rápido \ldots\ tú} ; \esp{tan alto \ldots\ su padre} ; \esp{más \ldots\ cien espectadores} ; \esp{menos tiempo \ldots\ antes}.
\item Corrige la faute : \esp{Este jugador es más bueno que aquel.}
\end{enumerate}
\end{exercice}

\begin{corrige}[Exercice « À toi de comparer »]
1) \esp{La natación es menos peligrosa que el boxeo.} ; \esp{Mi hermano es mayor que yo.} ; \esp{Hay más de treinta equipos en el torneo.} ; \esp{Ella nada tan bien como su hermana.} ; \esp{Este partido es peor que el otro.}
2) \esp{que} ; \esp{como} ; \esp{de} ; \esp{que}.
3) On dit \esp{Este jugador es mejor que aquel.} : \esp{bueno} possède un comparatif irrégulier, \esp{mejor}.
\end{corrige}

\coursec{Grammaire 2 : les superlatifs}

Dans la section précédente, nous avons appris à \cle{comparer} deux éléments : \esp{más rápido que}, \esp{menos caro que}, \esp{tan bueno como}. Mais que se passe-t-il quand un élément se détache de \emph{tout} le groupe ? Quand un joueur n'est pas seulement « plus rapide qu'un autre », mais « le plus rapide de l'équipe » ? C'est exactement le rôle du \cle{superlatif}, que nous allons maintenant étudier en deux temps : le superlatif relatif (le plus\ldots{} de) et le superlatif absolu en \esp{-ísimo} (très, extrêmement).

\courssub{Observation : partons des exemples}

Reprenons le dialogue après le match de la section 1. Trois phrases ont attiré notre attention :

\begin{exemplebox}[Phrases observées dans le dialogue]
\esp{Para mí, Eto'o es el mejor jugador del campeonato.}\\
« Pour moi, Eto'o est le meilleur joueur du championnat. »

\esp{Y Andrés es el más inteligente del campo: siempre sabe dónde está el balón.}\\
« Et Andrés est le plus intelligent du terrain : il sait toujours où est le ballon. »

\esp{Fue el partido más emocionante del año.}\\
« Ce fut le match le plus passionnant de l'année. »
\end{exemplebox}

Observons bien la construction. Dans chaque phrase, on trouve : un \textbf{article défini} (\esp{el}), parfois un \textbf{nom} (\esp{jugador}, \esp{partido}), \esp{más} ou \esp{mejor}, un \textbf{adjectif}, puis la préposition \esp{de} qui introduit le \textbf{groupe de référence} (\esp{del campeonato}, \esp{del campo}, \esp{del año}). C'est exactement la structure française « le meilleur joueur \emph{du} championnat », mais avec une différence capitale que nous verrons : l'espagnol utilise toujours \esp{de}, jamais \esp{que} ni \esp{en}.

\courssub{Le superlatif relatif : le plus\ldots{} de}

\begin{propriete}[Formation du superlatif relatif]
Le superlatif relatif désigne l'élément qui possède une qualité au \textbf{degré maximal} (ou minimal) à l'intérieur d'un groupe. Sa structure est :

\begin{center}
\textbf{article défini + (nom) + más / menos + adjectif + de + groupe de référence}
\end{center}

\begin{itemize}
\item L'\textbf{article} s'accorde en genre et en nombre avec le nom : \esp{el}, \esp{la}, \esp{los}, \esp{las}.
\item \esp{más} = « le plus » ; \esp{menos} = « le moins ».
\item La préposition \esp{de} introduit le groupe de référence : \textbf{jamais} \esp{que} (qui sert au comparatif) ni \esp{en} (calque du français « dans »).
\item Attention aux contractions : \esp{de + el = del} ; \esp{de + la} reste \esp{de la}.
\end{itemize}
\end{propriete}

\begin{exemplebox}[Exemples traduits]
\esp{Messi es el jugador más famoso del mundo.}\\
« Messi est le joueur le plus célèbre du monde. »

\esp{Eto'o es el mejor delantero de la historia de Camerún.}\\
« Eto'o est le meilleur attaquant de l'histoire du Cameroun. »

\esp{Es la jugadora menos conocida del equipo.}\\
« C'est la joueuse la moins connue de l'équipe. »

\esp{El Ahmadou Ahidjo es el estadio más grande de Yaoundé.}\\
« Le stade Ahmadou Ahidjo est le plus grand stade de Yaoundé. »

\esp{Los lunes son los días menos buenos para entrenar.}\\
« Les lundis sont les jours les moins bons pour s'entraîner. »
\end{exemplebox}

\begin{attention}[« de » et non « que » ni « en »]
Le francophone commet deux erreurs typiques :
\begin{itemize}
\item Utiliser \esp{que} : \esp{el más rápido que el equipo} est \textbf{faux}. \esp{que} sert uniquement à comparer deux éléments (\esp{más rápido que su hermano}). Pour le superlatif, il faut \esp{de} : \esp{el más rápido del equipo}.
\item Traduire « le meilleur \emph{du} monde » par \esp{en el mundo} : c'est un calque. On dit \esp{el mejor del mundo}. La préposition espagnole du superlatif est toujours \esp{de}.
\end{itemize}
\end{attention}

\courssub{La place de l'adjectif et le nom sous-entendu}

En espagnol, l'adjectif superlatif se place normalement \textbf{après le nom} :

\esp{el jugador más rápido} = « le joueur le plus rapide » ; \esp{la chica más alta} = « la fille la plus grande ».

Mais quand le nom est évident, on peut le \textbf{sous-entendre} : l'article porte alors seul le superlatif.

\begin{exemplebox}[Nom exprimé ou sous-entendu]
\esp{¿Quién es el jugador más rápido del equipo? — Pedro es el más rápido.}\\
« Qui est le joueur le plus rapide de l'équipe ? — Pedro est le plus rapide. »

\esp{De todas las camisetas, esta es la más bonita.}\\
« De tous les maillots, celui-ci est le plus beau. »

\esp{María es la más joven de la clase.}\\
« María est la plus jeune de la classe. »
\end{exemplebox}

Remarquez que le français fait exactement la même chose : « le plus rapide » peut se dire sans répéter « joueur ». C'est un point de convergence utile entre les deux langues.

\courssub{Les superlatifs irréguliers : el mejor, el peor, el mayor, el menor}

Comme pour les comparatifs, quatre adjectifs possèdent des formes superlatives \textbf{irrégulières} qu'il faut connaître par cœur. Ce sont les mêmes mots que les comparatifs irréguliers (\esp{mejor}, \esp{peor}, \esp{mayor}, \esp{menor}), précédés de l'article défini.

\begin{center}
\begin{tabularx}{\linewidth}{|X|X|X|X|}
\hline
\rowcolor{popblueL} Adjectif & Comparatif & Superlatif & Français \\
\hline
bueno (bon) & mejor & el mejor & le meilleur \\
\hline
malo (mauvais) & peor & el peor & le pire \\
\hline
grande (grand) & mayor & el mayor & l'aîné / le plus important \\
\hline
pequeño (petit) & menor & el menor & le cadet / le moindre \\
\hline
\end{tabularx}
\end{center}

\begin{popfigure}
\begin{center}
\begin{tikzpicture}[node distance=0.55cm and 1.1cm,
  base/.style={rectangle, rounded corners=3pt, draw=popblue, fill=popblueL, minimum width=2.3cm, minimum height=0.85cm, align=center, font=\small},
  comp/.style={rectangle, rounded corners=3pt, draw=poporange, fill=poporangeL, minimum width=2.3cm, minimum height=0.85cm, align=center, font=\small},
  sup/.style={rectangle, rounded corners=3pt, draw=popgreen, fill=popgreenL, minimum width=2.5cm, minimum height=0.85cm, align=center, font=\small\bfseries},
  arr/.style={-{Stealth[length=2.5mm]}, thick, popmuted}]
  \node[base, align=center] (b1){bueno\\ \emph{bon}};
  \node[comp, right=of b1, align=center] (b2){mejor\\ \emph{meilleur}};
  \node[sup, right=of b2, align=center] (b3){el mejor\\ \emph{le meilleur}};
  \draw[arr] (b1) -- (b2); \draw[arr] (b2) -- (b3);
  \node[base, below=of b1, align=center] (m1){malo\\ \emph{mauvais}};
  \node[comp, right=of m1, align=center] (m2){peor\\ \emph{pire}};
  \node[sup, right=of m2, align=center] (m3){el peor\\ \emph{le pire}};
  \draw[arr] (m1) -- (m2); \draw[arr] (m2) -- (m3);
  \node[base, below=of m1, align=center] (g1){grande\\ \emph{grand}};
  \node[comp, right=of g1, align=center] (g2){mayor\\ \emph{plus âgé}};
  \node[sup, right=of g2, align=center] (g3){el mayor\\ \emph{l'aîné}};
  \draw[arr] (g1) -- (g2); \draw[arr] (g2) -- (g3);
  \node[base, below=of g1, align=center] (p1){pequeño\\ \emph{petit}};
  \node[comp, right=of p1, align=center] (p2){menor\\ \emph{plus jeune}};
  \node[sup, right=of p2, align=center] (p3){el menor\\ \emph{le cadet}};
  \draw[arr] (p1) -- (p2); \draw[arr] (p2) -- (p3);
\end{tikzpicture}
\end{center}
\legende{Les quatre superlatifs irréguliers : de l'adjectif au superlatif}
\end{popfigure}

\begin{exemplebox}[Exemples traduits avec les formes irrégulières]
\esp{Eto'o es el mejor delantero de la historia de Camerún.}\\
« Eto'o est le meilleur attaquant de l'histoire du Cameroun. »

\esp{Fue el peor partido de la temporada.}\\
« Ce fut le pire match de la saison. »

\esp{Mi hermano mayor es el mayor de la familia y juega al baloncesto.}\\
« Mon frère aîné est l'aîné de la famille et joue au basket-ball. »

\esp{Soy el menor de mis hermanos, pero el más rápido.}\\
« Je suis le cadet de mes frères, mais le plus rapide. »
\end{exemplebox}

\begin{attention}[Ne dis jamais « el más bueno »]
Les francophones traduisent souvent « le meilleur » par \esp{el más bueno} : c'est une \textbf{erreur grave} en espagnol standard. La forme correcte est \esp{el mejor}. De même, « le pire » se dit \esp{el peor}, jamais \esp{el más malo}.

Nuance utile : \esp{el más grande} existe, mais il signifie « le plus grand \emph{en taille} » (\esp{el edificio más grande de Douala} = « le plus grand bâtiment de Douala »). \esp{el mayor} renvoie à l'âge (l'aîné) ou à l'importance (\esp{el mayor problema} = « le plus grand problème »). De même, \esp{el más pequeño} = « le plus petit (en taille) », \esp{el menor} = « le cadet, le moindre ».
\end{attention}

\courssub{Le superlatif absolu en -ísimo : très, extrêmement}

Il existe un second superlatif, dit \cle{absolu}, qui ne compare plus rien : il exprime simplement un degré \textbf{très élevé} d'une qualité. En français, on le traduit par « très », « vraiment », « extrêmement ». En espagnol, on ajoute le suffixe \esp{-ísimo} à l'adjectif (ou on utilise \esp{muy} + adjectif).

\begin{propriete}[Formation du superlatif absolu en -ísimo]
\begin{itemize}
\item Adjectif terminé par une voyelle : on supprime la voyelle finale et on ajoute \esp{-ísimo / -ísima / -ísimos / -ísimas} : \esp{rápido} $\rightarrow$ \esp{rapidísimo} ; \esp{bonita} $\rightarrow$ \esp{bonitísima}.
\item Adjectif terminé par une consonne : on ajoute directement \esp{-ísimo} : \esp{fácil} $\rightarrow$ \esp{facilísimo} ; \esp{feliz} $\rightarrow$ \esp{felicísimo}.
\item Changements orthographiques pour conserver le son : \esp{z} $\rightarrow$ \esp{c} (\esp{feliz} $\rightarrow$ \esp{felicísimo}) ; \esp{c} $\rightarrow$ \esp{qu} (\esp{rico} $\rightarrow$ \esp{riquísimo}) ; \esp{g} $\rightarrow$ \esp{gu} (\esp{largo} $\rightarrow$ \esp{larguísimo}).
\item L'accent porte toujours sur le \esp{í} du suffixe : \esp{emocionantísimo}, \esp{buenísimo}.
\item L'accord se fait avec le nom : \esp{una final emocionantísima}, \esp{unos goles buenísimos}.
\end{itemize}
\end{propriete}

\begin{exemplebox}[Exemples traduits]
\esp{El partido fue emocionantísimo.}\\
« Le match était vraiment passionnant. »

\esp{Corre rapidísimo: nadie puede seguirlo.}\\
« Il court extrêmement vite : personne ne peut le suivre. »

\esp{El examen de español fue facilísimo.}\\
« L'examen d'espagnol était très facile. »

\esp{Estoy felicísimo porque mi equipo ganó la copa.}\\
« Je suis extrêmement heureux parce que mon équipe a gagné la coupe. »
\end{exemplebox}

\begin{center}
\begin{tabularx}{\linewidth}{|X|X|X|}
\hline
\rowcolor{popblueL} Français & Español & Remarque \\
\hline
le plus\ldots{} de & el más\ldots{} de & \esp{de} introduit le groupe, jamais \esp{en} \\
\hline
le moins\ldots{} de & el menos\ldots{} de & même structure \\
\hline
le meilleur & el mejor & jamais \esp{más bueno} \\
\hline
le pire & el peor & jamais \esp{más malo} \\
\hline
très / extrêmement\ldots{} & muy\ldots{} ou -ísimo & \esp{muy rápido} = \esp{rapidísimo} \\
\hline
\end{tabularx}
\end{center}

\begin{savaistu}[Le -ísimo des commentateurs sportifs]
Le suffixe \esp{-ísimo} est extrêmement vivant dans le langage des commentateurs sportifs hispanophones. Pendant un match des Lions Indomptables ou du Real Madrid, vous entendrez : \esp{¡Golazo buenísimo!} (« Quel but magnifique ! »), \esp{¡Una final emocionantísima!} (« Une finale vraiment palpitante ! »), \esp{El portero está altísimo hoy} (« Le gardien est en très grande forme aujourd'hui »). À la radio espagnole, on entend même des superlatifs créés à la volée, comme \esp{golazo} (de \esp{gol} + suffixe augmentatif \esp{-azo}), preuve que la langue sportive est un formidable laboratoire de créativité.
\end{savaistu}

\begin{exoresolu}[Transforme les phrases au superlatif]
Énoncé : transforme chaque phrase comparative en phrase superlative, comme dans le modèle.\\
Modèle : \esp{Pedro es más rápido que Juan.} $\rightarrow$ \esp{Pedro es el más rápido del equipo.}
\begin{enumerate}
\item \esp{Esta jugadora es más buena que las otras.} (de l'équipe)
\item \esp{Este partido es más emocionante que los otros.} (de l'année)
\item \esp{Mi hermana es menor que mis otros hermanos.} (de la famille)
\end{enumerate}
\tcblower
Solution détaillée :
\begin{enumerate}
\item \esp{Esta jugadora es la mejor del equipo.} — « Cette joueuse est la meilleure de l'équipe. » Attention : « la meilleure » ne se dit pas \esp{la más buena} mais \esp{la mejor} (forme irrégulière), et l'article devient \esp{la} pour l'accord avec \esp{jugadora}.
\item \esp{Este partido es el más emocionante del año.} — « Ce match est le plus passionnant de l'année. » Adjectif régulier : article \esp{el} + \esp{más} + adjectif + \esp{de} + groupe ; contraction \esp{de + el año = del año}.
\item \esp{Mi hermana es la menor de la familia.} — « Ma sœur est la cadette de la famille. » Superlatif irrégulier de \esp{pequeño} pour l'âge : \esp{la menor} ; pas de contraction devant \esp{la familia}.
\end{enumerate}
\end{exoresolu}

\begin{exercice}[À toi de jouer]
\begin{enumerate}
\item Traduis en espagnol : « Samuel Eto'o est le meilleur joueur de l'histoire du football camerounais. »
\item Traduis : « C'est la joueuse la moins connue de l'équipe, mais elle court extrêmement vite. »
\item Complète avec la forme correcte : \esp{Este estadio es \ldots{} (le plus grand) de la ciudad.}
\item Mets au superlatif absolu : \esp{La final fue muy emocionante.} $\rightarrow$ \esp{La final fue \ldots}
\end{enumerate}
\end{exercice}

\begin{corrige}[Exercice « À toi de jouer »]
\begin{enumerate}
\item \esp{Samuel Eto'o es el mejor jugador de la historia del fútbol camerunés.} — « le meilleur » = \esp{el mejor} (jamais \esp{más bueno}) ; contraction \esp{de + el fútbol = del fútbol}.
\item \esp{Es la jugadora menos conocida del equipo, pero corre rapidísimo.} — superlatif avec \esp{menos} ; \esp{rapidísimo} avec accent sur le \esp{í} (invariable comme adverbe de manière).
\item \esp{Este estadio es el más grande de la ciudad.} — pour la taille, \esp{el más grande} est correct ; \esp{el mayor} conviendrait pour l'âge ou l'importance.
\item \esp{La final fue emocionantísima.} — accord féminin singulier avec \esp{la final}.
\end{enumerate}
\end{corrige}

\begin{aretenir}[L'essentiel sur les superlatifs]
\begin{itemize}
\item Superlatif relatif : \textbf{article + (nom) + más/menos + adjectif + de + groupe} : \esp{el jugador más famoso del mundo}.
\item La préposition est toujours \esp{de}, jamais \esp{que} ni \esp{en}.
\item Formes irrégulières : \esp{el mejor} (le meilleur), \esp{el peor} (le pire), \esp{el mayor} (l'aîné), \esp{el menor} (le cadet).
\item Superlatif absolu en \esp{-ísimo} = « très, extrêmement » : \esp{rapidísimo}, \esp{felicísimo}, \esp{emocionantísimo} ; attention aux changements \esp{z} $\rightarrow$ \esp{c}, \esp{c} $\rightarrow$ \esp{qu}, \esp{g} $\rightarrow$ \esp{gu}.
\end{itemize}
\end{aretenir}

\coursec{Traduction : phrases sur le sport (thème et version)}

Dans la section précédente, nous avons étudié les \cle{superlatifs} : \esp{el más alto}, \esp{la mejor jugadora}, \esp{el partido más importante del año}. Nous savons donc maintenant comparer et situer un élément au sommet d'un groupe. Il est temps de mettre ces outils en pratique dans l'exercice roi de l'apprentissage des langues : la \cle{traduction}. Traduire, ce n'est pas remplacer mot à mot : c'est \textbf{reconnaître une structure} dans la langue de départ, choisir la \textbf{structure équivalente} dans la langue d'arrivée, puis vérifier les accords et les formes irrégulières. C'est exactement ce que nous allons faire avec les phrases comparatives sur le thème du sport.

\courssub{Observer avant de traduire}

Lisons d'abord ce court texte, qui rassemble presque toutes les structures de comparaison que nous devrons traduire.

\begin{textebox}[Después del torneo escolar]
El torneo de fútbol del liceo terminó ayer. Nuestro equipo jugó mejor que el año pasado, pero el equipo de Douala fue el mejor del campeonato. Su delantero es más rápido que todos nuestros defensores y marca tantos goles como todo nuestro equipo al completo. El partido final fue más emocionante que la semifinal, y hubo más de mil aficionados en las gradas. El deporte escolar es cada vez más popular entre los jóvenes cameruneses. Para mí, el fútbol es tan importante como los estudios: entreno tanto como estudio. Mi entrenador dice que soy el jugador más disciplinado del grupo, aunque no el más técnico. El año que viene queremos ser mejores que este año.
\end{textebox}

\textbf{Glossaire :} \esp{torneo} : tournoi ; \esp{delantero} : attaquant ; \esp{al completo} : au complet ; \esp{gradas} : tribunes ; \esp{aunque} : bien que ; \esp{entrenador} : entraîneur ; \esp{disciplinado} : discipliné.

\textbf{Questions de repérage :}
\begin{enumerate}
\item Repère dans le texte deux comparatifs de supériorité, un comparatif d'égalité et deux superlatifs.
\item Comment l'auteur traduit-il « de plus en plus populaire » ?
\item Quelle structure exprime « plus de mille » ?
\end{enumerate}

\courssub{Version : de l'espagnol vers le français}

\begin{methode}[Traduire une phrase comparative (version)]
\begin{enumerate}
\item \textbf{Repérer} la structure comparative espagnole : \esp{más\ldots que}, \esp{menos\ldots que}, \esp{tan\ldots como}, \esp{tanto\ldots como}, \esp{el más\ldots de}, \esp{cada vez más}.
\item \textbf{Identifier} la nature du mot comparé : adjectif, adverbe, nom ou chiffre.
\item \textbf{Choisir} l'équivalent français exact : « plus\ldots que », « aussi\ldots que », « autant de\ldots que », « le plus\ldots de », « de plus en plus ».
\item \textbf{Traduire} le reste de la phrase en français naturel, sans calquer l'ordre espagnol.
\item \textbf{Relire} : le français doit être correct et idiomatique.
\end{enumerate}
\end{methode}

\begin{exoresolu}[Version : quatre phrases à traduire]
Traduisez en français :\\
1. \esp{El balonmano es menos popular que el fútbol en Camerún.}\\
2. \esp{Mi hermana nada tan rápido como yo.}\\
3. \esp{Samuel Eto'o es el jugador más famoso de la historia del fútbol camerunés.}\\
4. \esp{Hay más de veinte equipos en el torneo y cada vez más espectadores.}
\tcblower
\textbf{Solutions détaillées :}\\
1. Structure : \esp{menos\ldots que} + adjectif = comparatif d'infériorité. → « Le handball est moins populaire que le football au Cameroun. »\\
2. Structure : \esp{tan\ldots como} + adverbe = comparatif d'égalité. → « Ma sœur nage aussi vite que moi. »\\
3. Structure : \esp{el más\ldots de} = superlatif relatif. → « Samuel Eto'o est le joueur le plus célèbre de l'histoire du football camerounais. »\\
4. Deux structures : \esp{más de} + chiffre = « plus de » ; \esp{cada vez más} = « de plus en plus ». → « Il y a plus de vingt équipes dans le tournoi et de plus en plus de spectateurs. »
\end{exoresolu}

\courssub{Thème : du français vers l'espagnol}

Le thème est plus exigeant : c'est vous qui devez \textbf{construire} la structure espagnole. Voici la démarche complète.

\begin{methode}[Réussir le thème : la méthode en 4 étapes]
\begin{enumerate}
\item \textbf{Souligner} la structure comparative française : « plus\ldots que », « aussi\ldots que », « autant de\ldots que », « le meilleur », « de plus en plus », « plus de + chiffre ».
\item \textbf{Choisir} l'équivalent espagnol exact : \esp{más\ldots que}, \esp{tan\ldots como}, \esp{tanto\ldots como}, \esp{el mejor}, \esp{cada vez más}, \esp{más de}.
\item \textbf{Vérifier l'accord} de l'adjectif (genre et nombre) et celui de \esp{tanto} : \esp{tanto, tanta, tantos, tantas}.
\item \textbf{Contrôler} les formes irrégulières : \esp{mejor} (meilleur/mieux), \esp{peor} (pire), \esp{mayor} (plus grand, plus âgé), \esp{menor} (plus jeune).
\end{enumerate}
\end{methode}

\begin{exemplebox}[Exemples de thème commentés]
\begin{itemize}
\item « Mon frère est plus grand que moi. » = \esp{Mi hermano es más alto que yo.} (comparatif de supériorité + pronom sujet \esp{yo})
\item « Ce match est le plus important de la saison. » = \esp{Este partido es el más importante de la temporada.} (superlatif relatif : article + \esp{más} + adjectif)
\item « Elle court autant que son frère. » = \esp{Ella corre tanto como su hermano.} (\esp{tanto} adverbe, invariable, avec un verbe)
\item « Il y a plus de vingt équipes. » = \esp{Hay más de veinte equipos.} (\esp{más de} devant un chiffre)
\end{itemize}
\end{exemplebox}

\begin{center}
\begin{tabularx}{\linewidth}{|X|X|X|}
\hline
\rowcolor{popblueL} Français & Español & Remarque \\
\hline
plus\ldots que & más\ldots que & supériorité \\
\hline
moins\ldots que & menos\ldots que & infériorité \\
\hline
aussi\ldots que & tan\ldots como & égalité, adjectif ou adverbe \\
\hline
autant de\ldots que & tanto/a/os/as\ldots como & égalité + nom, accord obligatoire \\
\hline
le plus\ldots de & el más\ldots de & superlatif relatif \\
\hline
le meilleur / mieux & el mejor / mejor & adjectif / adverbe \\
\hline
de plus en plus & cada vez más & progression \\
\hline
plus de + chiffre & más de + cifra & quantité numérique \\
\hline
\end{tabularx}
\end{center}

\begin{popfigure}
\begin{tikzpicture}[
  node distance=0.55cm and 0.9cm,
  every node/.style={font=\small},
  decision/.style={diamond, draw=popdark, fill=poppurpleL, align=center, inner sep=1pt, aspect=2},
  action/.style={rectangle, rounded corners, draw=popdark, fill=popblueL, align=center, inner sep=4pt},
  result/.style={rectangle, rounded corners, draw=popdark, fill=popgreenL, align=center, inner sep=4pt},
  fleche/.style={-{Stealth[length=2.5mm]}, thick, popdark}
]
\node[decision] (q1) {Que comparer ?};
\node[action, below left=0.6cm and 2.2cm of q1, align=center] (adj){adjectif\\ou adverbe};
\node[action, below=0.6cm of q1] (nom) {un nom};
\node[action, below right=0.6cm and 2.2cm of q1] (chiffre) {un chiffre};
\node[result, below=0.6cm of adj, align=center] (resadj){más\ldots que\\menos\ldots que\\tan\ldots como};
\node[result, below=0.6cm of nom, align=center] (resnom){tanto/a/os/as\\\ldots como};
\node[result, below=0.6cm of chiffre, align=center] (resnum){más de\ldots\\menos de\ldots};
\draw[fleche] (q1) -- (adj);
\draw[fleche] (q1) -- (nom);
\draw[fleche] (q1) -- (chiffre);
\draw[fleche] (adj) -- (resadj);
\draw[fleche] (nom) -- (resnom);
\draw[fleche] (chiffre) -- (resnum);
\end{tikzpicture}
\legende{Organigramme de décision : choisir la bonne structure comparative en espagnol.}
\end{popfigure}

\courssub{Les cinq pièges classiques des francophones}

\begin{attention}[Piège 1 : « aussi\ldots que » n'est pas « también\ldots que »]
Le français « aussi\ldots que » se traduit par \esp{tan\ldots como}. \esp{También} signifie « également, aussi » au sens d'addition, jamais dans une comparaison.\\
« Le volley-ball est aussi amusant que le football. » = \esp{El voleibol es tan divertido como el fútbol.}\\
Erreur à bannir : \esp{El voleibol es también divertido que el fútbol.}
\end{attention}

\begin{attention}[Piège 2 : « le meilleur » et « mieux » ne se traduisent pas pareil]
« Le meilleur » est un \textbf{adjectif} : \esp{el mejor, la mejor, los mejores, las mejores}. « Mieux » est un \textbf{adverbe} : \esp{mejor}, invariable.\\
« Il joue mieux que moi. » = \esp{Él juega mejor que yo.}\\
« C'est le meilleur joueur de l'équipe. » = \esp{Es el mejor jugador del equipo.}
\end{attention}

\begin{aretenir}[L'accord de « mejor »]
\esp{Mejor} est \textbf{invariable} quand il est adverbe : \esp{Ella canta mejor.} (« Elle chante mieux. »)\\
Mais il \textbf{s'accorde en nombre} quand il est adjectif : \esp{los mejores jugadores} (« les meilleurs joueurs »), \esp{las mejores atletas} (« les meilleures athlètes »).
\end{aretenir}

\begin{attention}[Piège 3 : « de plus en plus » = cada vez más]
Ne traduisez jamais mot à mot par \esp{más y más}. La structure correcte est \esp{cada vez más} + adjectif ou nom.\\
« Le sport est de plus en plus populaire. » = \esp{El deporte es cada vez más popular.}
\end{attention}

\begin{attention}[Piège 4 : « autant de\ldots que » exige l'accord de « tanto »]
Devant un nom, \esp{tanto} s'accorde en genre et en nombre : \esp{tanto dinero como tiempo}, \esp{tanta energía como tú}, \esp{tantos partidos como entrenamientos} (« autant de matchs que d'entraînements »), \esp{tantas chicas como chicos}.\\
Devant un verbe, \esp{tanto} reste invariable : \esp{Entreno tanto como mi hermano.} (« Je m'entraîne autant que mon frère. »)
\end{attention}

\begin{attention}[Piège 5 : « plus de + chiffre » = más de (et non « más que »)]
Devant une quantité chiffrée, on emploie \esp{más de} : \esp{más de mil aficionados} (« plus de mille supporters »), \esp{más de tres horas de entrenamiento}. De même : \esp{menos de diez minutos} (« moins de dix minutes »).
\end{attention}

\begin{attention}[Piège 6 : « que moi, que toi, que lui » = pronom sujet]
Après \esp{que} dans une comparaison, l'espagnol exige le \textbf{pronom sujet} : \esp{que yo, que tú, que él, que ella, que nosotros, que ellos}. N'écrivez jamais \esp{que mí} ou \esp{que ti} dans ce contexte.\\
« Il est plus rapide que moi. » = \esp{Él es más rápido que yo.}\\
« Tu cours plus vite que lui. » = \esp{Tú corres más rápido que él.}
\end{attention}

\courssub{Entraînement guidé}

\begin{exoresolu}[Thème : quatre phrases à traduire]
Traduisez en espagnol :\\
1. « Le basket-ball est plus difficile que le football. »\\
2. « Ma sœur court aussi vite que moi. »\\
3. « Cette joueuse est la meilleure de l'équipe. »\\
4. « Nous avons autant de matchs que d'entraînements, et le sport est de plus en plus populaire. »
\tcblower
\textbf{Solutions détaillées :}\\
1. Structure « plus\ldots que » → \esp{más\ldots que} ; adjectif masculin singulier. → \esp{El baloncesto es más difícil que el fútbol.}\\
2. « aussi\ldots que » → \esp{tan\ldots como} ; « que moi » → \esp{que yo} (pronom sujet). → \esp{Mi hermana corre tan rápido como yo.}\\
3. Superlatif : article féminin + \esp{mejor} ; « de l'équipe » → \esp{del equipo} (contraction \esp{de + el}). → \esp{Esta jugadora es la mejor del equipo.}\\
4. « autant de\ldots que » → \esp{tantos\ldots como} (accord masculin pluriel) ; « de plus en plus » → \esp{cada vez más}. → \esp{Tenemos tantos partidos como entrenamientos, y el deporte es cada vez más popular.}
\end{exoresolu}

\begin{exercice}[À vous de jouer]
Traduisez en espagnol :\\
1. « Le stade est plus grand que le gymnase. »\\
2. « Il y a plus de cinq cents spectateurs. »\\
3. « Mon champion préféré joue mieux que tous les autres. »\\
4. « Les filles s'entraînent autant que les garçons. »\\
Traduisez en français :\\
5. \esp{La natación es menos peligrosa que el boxeo.}\\
6. \esp{Este torneo es el más importante del año.}
\end{exercice}

\begin{corrige}[Exercice « À vous de jouer »]
1. \esp{El estadio es más grande que el gimnasio.}\\
2. \esp{Hay más de quinientos espectadores.} (\esp{más de} + chiffre)\\
3. \esp{Mi campeón favorito juega mejor que todos los demás.} (\esp{mejor} adverbe, invariable)\\
4. \esp{Las chicas entrenan tanto como los chicos.} (\esp{tanto} adverbe avec un verbe)\\
5. « La natation est moins dangereuse que la boxe. »\\
6. « Ce tournoi est le plus important de l'année. »
\end{corrige}

\begin{aretenir}[La stratégie de traduction en une phrase]
Identifier le type de comparaison (supériorité, infériorité, égalité, superlatif, quantité chiffrée, progression) → choisir la structure espagnole exacte → accorder l'adjectif et \esp{tanto} → vérifier les formes irrégulières \esp{mejor, peor, mayor, menor} → relire pour chasser les calques du français.
\end{aretenir}

Maintenant que nous savons comparer et traduire avec précision, nous possédons tous les outils pour la dernière étape de la leçon : la production guidée, où vous rédigerez un court texte pour présenter votre sport préféré ou votre champion favori.

\coursec{Production guidée : mon sport préféré ou mon champion (80 mots)}

Dans la section précédente, nous avons traduit des phrases sur le sport, du français vers l'espagnol et de l'espagnol vers le français. Nous disposons maintenant de tout le matériel nécessaire : le lexique du sport et des loisirs, les comparatifs (\esp{más... que}, \esp{menos... que}, \esp{tan... como}) et les superlatifs (\esp{el mejor}, \esp{el más... de}). Il est temps de mobiliser ces acquis dans une \cle{production écrite} complète : rédiger un texte de 80 mots pour présenter son sport préféré ou son champion favori. C'est exactement le type de sujet que l'on rencontre à l'évaluation et au baccalauréat.

\courssub{Observer d'abord : un modèle de production}

Avant d'écrire, lisons attentivement le texte d'une élève de Première, Marlyse, qui présente son sport préféré. Observons comment il est construit.

\begin{textebox}[Mi deporte favorito, par Marlyse, élève de Première à Yaoundé]
Mi deporte favorito es el balonmano. Lo practico desde hace cuatro años en el club de mi colegio, los martes y los jueves por la tarde. Empecé gracias a mi hermana mayor, que jugaba en el equipo de la ciudad.

Para mí, el balonmano es más emocionante que el fútbol porque hay más goles y el juego es más rápido. Además, es menos peligroso que el boxeo y tan colectivo como el baloncesto: sin mis compañeras, no puedo ganar. Mi jugadora favorita es Carmen Martín; para mí, es la mejor jugadora del mundo y la más generosa con sus aficionados.

El balonmano me ayuda a estar en forma y me relaja después de los deberes. Cuando juego, olvido el estrés de los exámenes. Por eso animo a todos los jóvenes de mi barrio a practicar un deporte. En conclusión, el balonmano es el deporte más completo para mí: me da salud, amigas y alegría.
\end{textebox}

\textbf{Glossaire}

\begin{center}
\begin{tabularx}{\linewidth}{|X|X|}
\hline
\rowcolor{popblueL} \textbf{Español} & \textbf{Français} \\
\hline
el balonmano & le handball \\
\hline
por la tarde & l'après-midi \\
\hline
empecé & j'ai commencé \\
\hline
gracias a & grâce à \\
\hline
los goles & les buts \\
\hline
peligroso & dangereux \\
\hline
las compañeras & les camarades (d'équipe) \\
\hline
generosa con sus aficionados & généreuse avec ses supporters \\
\hline
estar en forma & être en forme \\
\hline
me relaja & me détend \\
\hline
los deberes & les devoirs \\
\hline
animo a... a & j'encourage... à \\
\hline
la alegría & la joie \\
\hline
\end{tabularx}
\end{center}

\textbf{Questions de compréhension et d'observation}

\begin{enumerate}
\item ¿Cuál es el deporte favorito de Marlyse? ¿Desde cuándo lo practica?
\item ¿Qué comparativos usa en el segundo párrafo? Subráyalos.
\item ¿Qué superlativos aparecen en el texto?
\item ¿Qué conectores utiliza para organizar sus ideas?
\item ¿Cuántas partes tiene el texto? ¿Qué dice en cada parte?
\end{enumerate}

\begin{corrige}[Questions d'observation]
\begin{enumerate}
\item Su deporte favorito es el balonmano; lo practica desde hace cuatro años.
\item Los comparativos : \esp{más emocionante que}, \esp{más rápido}, \esp{menos peligroso que}, \esp{tan colectivo como}.
\item Los superlativos : \esp{la mejor jugadora del mundo}, \esp{la más generosa}, \esp{el deporte más completo}.
\item Los conectores : \esp{además}, \esp{porque}, \esp{por eso}, \esp{en conclusión}.
\item Tres partes : presentación del deporte, comparaciones y argumentos, bienestar y conclusión.
\end{enumerate}
\end{corrige}

\courssub{La méthode : un plan en cinq étapes pour trois paragraphes}

Le texte de Marlyse n'est pas écrit au hasard : il suit un plan rigoureux que l'on peut visualiser en cinq étapes logiques, regroupées en \cle{trois paragraphes}. Le schéma ci-dessous montre la progression du général vers le personnel.

\begin{popfigure}
\begin{center}
\begin{tikzpicture}[
  etape/.style={draw=popblue, fill=popblueL, rounded corners=3pt, align=center, text width=9cm, inner sep=5pt, font=\small},
  fleche/.style={-{Stealth[length=3mm]}, thick, poporange}
]
\node[etape] (p1) {\textbf{1. Presentación} : \esp{Mi deporte favorito es... / Lo practico desde hace...}};
\node[etape, fill=poporangeL, draw=poporange] (p2) at (0,-1.4) {\textbf{2. Comparación} : \esp{más... que / menos... que / tan... como}};
\node[etape, fill=popgreenL, draw=popgreen] (p3) at (0,-2.8) {\textbf{3. Superlatif} : \esp{el mejor / la más... del mundo}};
\node[etape, fill=poppurpleL, draw=poppurple] (p4) at (0,-4.2) {\textbf{4. Bienestar} : \esp{me ayuda a estar en forma, me relaja}};
\node[etape, fill=popgoldL, draw=popgold] (p5) at (0,-5.6) {\textbf{5. Conclusión} : \esp{Por eso... / En conclusión...}};
\draw[fleche] (p1) -- (p2);
\draw[fleche] (p2) -- (p3);
\draw[fleche] (p3) -- (p4);
\draw[fleche] (p4) -- (p5);
\end{tikzpicture}
\end{center}
\legende{Le plan en entonnoir : 5 étapes pour 3 paragraphes (1 | 2-3 | 4-5)}
\end{popfigure}

\begin{methode}[Rédiger ses 80 mots en cinq étapes]
\begin{enumerate}
\item \textbf{Paragraphe 1 — Présentation (environ 20 mots).} Nomme le sport ou le champion, dis depuis combien de temps tu le pratiques ou le suis, où et avec qui. Débute par : \esp{Mi deporte favorito es...} ou \esp{Mi campeón favorito es...} \textit{(Étape 1 du schéma)}.
\item \textbf{Paragraphe 2 — Arguments et comparaison (environ 35 mots).} Justifie ton choix avec \textbf{au moins deux comparatifs} et \textbf{un superlatif} : \esp{es más emocionante que...}, \esp{es menos caro que...}, \esp{es tan popular como...}, \esp{es el mejor jugador del mundo}. \textit{(Étapes 2 et 3 du schéma)}.
\item \textbf{Paragraphe 3 — Bien-être et conclusion (environ 25 mots).} Explique ce que ce sport t'apporte : \esp{me ayuda a estar en forma}, \esp{me relaja}, puis conclus avec \esp{por eso...} ou \esp{en conclusión...}. \textit{(Étapes 4 et 5 du schéma)}.
\item \textbf{Compter les mots.} Vise 80 mots, avec une marge de plus ou moins 10 \%, donc entre 72 et 88 mots. Compte-les ligne par ligne.
\item \textbf{Relire.} Vérifie les accords (masculin/féminin, singulier/pluriel), les accents, la conjugaison au présent et la présence obligatoire des deux comparatifs et du superlatif.
\end{enumerate}
\end{methode}

\begin{exemplebox}[Un début de production modèle]
\esp{Mi deporte favorito es el baloncesto. Es más dinámico que el fútbol y tan popular como él en mi barrio. Para mí, Joel Embiid es el mejor jugador del mundo.}

« Mon sport préféré est le basket. Il est plus dynamique que le football et aussi populaire que lui dans mon quartier. Pour moi, Joel Embiid est le meilleur joueur du monde. »

Observons : en trois phrases, l'élève a déjà placé un comparatif de supériorité (\esp{más dinámico que}), un comparatif d'égalité (\esp{tan popular como}) et un superlatif (\esp{el mejor jugador del mundo}). C'est un modèle d'efficacité.
\end{exemplebox}

\courssub{La banque de phrases toutes faites}

Pour rédiger vite et bien, il faut connaître par cœur des phrases toutes faites que l'on peut adapter à n'importe quel sujet. Voici la banque organisée selon le plan.

\begin{center}
\begin{tabularx}{\linewidth}{|X|X|}
\hline
\rowcolor{popblueL} \textbf{Español} & \textbf{Français} \\
\hline
Mi deporte favorito es... & Mon sport préféré est... \\
\hline
Mi campeón favorito es... & Mon champion préféré est... \\
\hline
Lo practico desde hace... años. & Je le pratique depuis... ans. \\
\hline
Empecé gracias a mi hermano / mi amigo. & J'ai commencé grâce à mon frère / mon ami. \\
\hline
Es más emocionante que... & C'est plus passionnant que... \\
\hline
Es menos peligroso que... & C'est moins dangereux que... \\
\hline
Es tan popular como... & C'est aussi populaire que... \\
\hline
Para mí, es el deporte más completo. & Pour moi, c'est le sport le plus complet. \\
\hline
Es el mejor jugador del mundo. & C'est le meilleur joueur du monde. \\
\hline
Además, me ayuda a estar en forma. & De plus, cela m'aide à rester en forme. \\
\hline
Me relaja después de los deberes. & Cela me détend après les devoirs. \\
\hline
Por eso animo a todos a practicarlo. & C'est pourquoi j'encourage tous à le pratiquer. \\
\hline
En conclusión, ... & En conclusion, ... \\
\hline
\end{tabularx}
\end{center}

\begin{attention}[« Depuis... ans » : le piège du francophone]
En français, on dit « je le pratique \emph{depuis} trois ans ». En espagnol, on ne dit jamais \esp{desde tres años}. La structure correcte est : \esp{desde hace + durée}, avec le verbe au \textbf{présent} :

\esp{Lo practico desde hace tres años.} = Je le pratique depuis trois ans.

\esp{Juego al fútbol desde hace diez años.} = Je joue au football depuis dix ans.

Littéralement, \esp{desde hace tres años} signifie « depuis que cela fait trois ans ». Ne traduis jamais mot à mot : \esp{hace} est obligatoire.
\end{attention}

\begin{definition}[Boîte à outils : les connecteurs de la rédaction]
Les \cle{connecteurs} donnent du rythme à ton texte et montrent au correcteur que tu sais organiser tes idées. Utilise-en au moins trois :

\begin{itemize}
\item \esp{primero} : d'abord (pour annoncer le premier argument) ;
\item \esp{además} : de plus, en outre (pour ajouter une idée) ;
\item \esp{sin embargo} : cependant (pour nuancer, opposer) ;
\item \esp{por eso} : c'est pourquoi (pour exprimer la conséquence) ;
\item \esp{en conclusión} : en conclusion (pour fermer le texte).
\end{itemize}

Exemple : \esp{El tenis es caro; sin embargo, es muy bueno para la salud.} = Le tennis est cher ; cependant, il est très bon pour la santé.
\end{definition}

\courssub{Exercice résolu : corriger et compléter un texte}

\begin{exoresolu}[À la recherche des quatre erreurs]
\textbf{Énoncé.} Voici le brouillon d'un élève. Il contient \textbf{quatre erreurs} (un calque du français, un comparatif mal formé, un superlatif mal formé, un accord verbal). Trouve-les, corrige-les, puis compte les mots du texte corrigé.

\esp{Mi deporte favorito es el atletismo. Lo practico desde tres años. Es más bueno que el ciclismo y tan sano como la natación. Para mí, es el más mejor deporte del mundo. Me ayuda a estar en forma y me relajo. En conclusión, el atletismo es mi vida.}

\tcblower
\textbf{Solution détaillée.}

\begin{enumerate}
\item \esp{desde tres años} est un calque du français « depuis trois ans ». Correction : \esp{desde hace tres años}.
\item \esp{más bueno que} est incorrect : le comparatif de \esp{bueno} est irrégulier. Correction : \esp{mejor que} (meilleur que).
\item \esp{el más mejor} est un double superlatif impossible. Correction : \esp{el mejor deporte del mundo}.
\item \esp{me relajo} : le verbe est à la 1\ière personne du singulier (\esp{yo}), mais le sujet est \esp{el atletismo} (ou l'activité implicite), 3\ième personne. Correction : \esp{me relaja}.
\end{enumerate}

\textbf{Texte corrigé :}

\esp{Mi deporte favorito es el atletismo. Lo practico desde hace tres años. Es mejor que el ciclismo y tan sano como la natación. Para mí, es el mejor deporte del mundo. Me ayuda a estar en forma y me relaja. En conclusión, el atletismo es mi vida.}

\textbf{Comptage :} 48 mots. Le texte est trop court pour la consigne de 80 mots : il faut développer le paragraphe 2 (ajouter un deuxième comparatif et des détails sur l'entraînement) et le paragraphe 3 (préciser les bienfaits).
\end{exoresolu}

\begin{aretenir}[Les points qui rapportent au Bac]
Un texte de 80 mots réussi contient obligatoirement :
\begin{itemize}
\item \textbf{2 comparatifs} corrects (dont éventuellement un irrégulier : \esp{mejor que}, \esp{peor que}, \esp{mayor que}) ;
\item \textbf{1 superlatif} correct (\esp{el mejor... de}, \esp{el más... de}) ;
\item \textbf{3 connecteurs} variés (\esp{además}, \esp{sin embargo}, \esp{por eso}, \esp{en conclusión}) ;
\item du \textbf{lexique du sport} précis (\esp{el entrenamiento}, \esp{el equipo}, \esp{ganar}, \esp{estar en forma}) ;
\item \textbf{80 mots ± 10 \%}, soit entre 72 et 88 mots ;
\item des accords et des accents relus.
\end{itemize}
\end{aretenir}

\courssub{La grille d'auto-évaluation}

Avant de rendre ta copie, évalue-toi toi-même avec cette grille. Chaque critère vaut un point : vise 6 sur 6.

\begin{center}
\begin{tabularx}{\linewidth}{|X|c|c|}
\hline
\rowcolor{popblueL} \textbf{Critère} & \textbf{Sí} & \textbf{No} \\
\hline
Mon texte compte entre 72 et 88 mots. & & \\
\hline
J'ai utilisé au moins 2 comparatifs corrects. & & \\
\hline
J'ai utilisé au moins 1 superlatif correct. & & \\
\hline
J'ai utilisé au moins 3 connecteurs différents. & & \\
\hline
J'ai employé au moins 5 mots du lexique du sport. & & \\
\hline
J'ai relu les accords, les accents et la conjugaison. & & \\
\hline
\end{tabularx}
\end{center}

\courssub{L'étape orale : présenter son texte en une minute}

L'écrit prépare l'oral. Une fois ton texte rédigé et auto-évalué, tu vas le présenter à la classe.

\begin{experience}[Mi deporte favorito : présentation orale en una hora... ¡no, en un minuto!]
\textbf{Déroulement.}
\begin{enumerate}
\item \textbf{Préparation (5 min).} Relis ton texte, souligne les mots difficiles à prononcer, mémorise le plan (pas le texte mot à mot).
\item \textbf{Présentation (1 min par élève).} Parle sans lire : regarde la classe, annonce ton sport, donne tes deux comparaisons et ton superlatif, termine par ta conclusion. Chronomètre en main, le professeur arrête à une minute.
\item \textbf{Questions des camarades (2 min).} Les auditeurs posent deux questions avec les structures de la leçon : \esp{¿Por qué es más interesante que...?}, \esp{¿Desde cuándo lo practicas?}, \esp{¿Quién es el mejor jugador para ti?}
\item \textbf{Conseils de prononciation.} \esp{baloncesto} se prononce « ba-lon-sès-to » (seseo) ou « ba-lon-thès-to » (ceceo) ; \esp{ejercicio} : « é-jèr-si-sio » (jota aspirée) ; \esp{mejor} : « mé-ror » (jota aspirée) ; \esp{desde} : « dès-de ».
\end{enumerate}
\end{experience}

\begin{exercice}[À toi d'écrire]
Rédige ton texte de 80 mots (± 10 \%) sur l'un des deux sujets suivants :
\begin{enumerate}
\item \esp{Mi deporte favorito} : présente ton sport préféré, compare-le à deux autres sports, place un superlatif et termine par ses bienfaits pour ta santé.
\item \esp{Mi campeón favorito} : présente un sportif camerounais ou international (football, basket, athlétisme...), compare-le à d'autres champions et explique pourquoi il est, pour toi, \esp{el mejor}.
\end{enumerate}
Utilise la banque de phrases, au moins trois connecteurs, puis remplis la grille d'auto-évaluation avant de présenter ton texte à la classe.
\end{exercice}

\begin{savaistu}[Le sujet star du baccalauréat]
Au baccalauréat camerounais, la production écrite sur « ton sport préféré » ou « un champion que tu admires » revient très régulièrement, car elle permet d'évaluer à la fois le lexique, les comparatifs et les superlatifs. Les correcteurs attribuent des points précis à la \emph{richesse de langue} : un texte qui contient deux comparatifs et un superlatif corrects, avec des connecteurs variés, obtient systématiquement une meilleure note qu'un texte plat, même sans faute. Et petit clin d'œil culturel : au Cameroun, parler de Samuel Eto'o ou de Joel Embiid avec un superlatif correct — \esp{Eto'o es el mejor delantero de la historia de Camerún} —, c'est allier la passion du sport et l'excellence en espagnol.
\end{savaistu}

\newpage

\coursec{Activité d'intégration}

\courssub{Objectifs de l'activité}

À la fin de cette activité d'intégration, tu seras capable de :
\begin{itemize}
\item défendre oralement un sport ou un champion devant la classe, en argumentant avec clarté ;
\item mobiliser le \cle{lexique du sport, des loisirs et du bien-être} dans une situation de communication réelle ;
\item employer correctement au moins \cle{deux comparatifs} (\esp{más... que}, \esp{menos... que}, \esp{tan... como}) et \cle{un superlatif} (\esp{el mejor}, \esp{el más...}) ;
\item exprimer ton \cle{opinion} avec des structures variées (\esp{creo que}, \esp{en mi opinión}, \esp{para mí}, \esp{pienso que}) ;
\item rédiger en groupe un \cle{article de 80 mots} pour le journal du lycée, structuré et correct ;
\item écouter les autres, comparer les argumentaires et voter de manière justifiée.
\end{itemize}

\courssub{Situation}

Le journal du lycée bilingue de Yaoundé, \emph{La Voz del Estudiante}, prépare son prochain numéro spécial « Sport et jeunesse ». La rédaction lance un concours entre les classes de Première : chaque groupe doit envoyer un article intitulé \esp{« El mejor deporte para los jóvenes cameruneses »}. Le meilleur article sera publié en première page, et le groupe gagnant représentera le lycée au forum des sports de la ville, où des élèves de Guinée équatoriale présenteront aussi leurs sports favoris en espagnol.

Voici l'appel à contribution publié par la rédaction :

\begin{textebox}[Appel à contribution — La Voz del Estudiante]
¡Atención, alumnos de Primera!\\
El journal del liceo organiza un gran debate: ¿cuál es el mejor deporte para los jóvenes cameruneses?\\
¿El fútbol, la pasión nacional? ¿El baloncesto, cada vez más popular en Douala y Yaoundé? ¿La carrera, el deporte más barato y más sencillo? ¿O quizás el voleibol, tan divertido como el fútbol?\\
En grupos de cuatro, preparen un mini-debate: cada alumno defiende un deporte o un campeón con argumentos claros. Después, escriban un artículo de 80 palabras comparando los deportes y concluyendo con un superlativo.\\
Los mejores artículos serán publicados y expuestos en el aula. ¡Que gane el mejor argumento!
\end{textebox}

\begin{definition}[Glossaire du document]
\esp{el debate} : le débat ; \esp{cada vez más popular} : de plus en plus populaire ; \esp{barato} : bon marché ; \esp{defender} : défendre ; \esp{escriban} : écrivez (subjonctif impératif) ; \esp{expuestos} : affichés ; \esp{¡Que gane el mejor argumento!} : que le meilleur argument gagne.
\end{definition}

\courssub{Consignes}

\begin{enumerate}
\item \textbf{Compréhension du document.} Lis l'appel à contribution et réponds par écrit en espagnol : \esp{¿Qué deportes menciona el texto? ¿Qué hay que escribir? ¿Cuántas palabras tiene el artículo?}
\item \textbf{Choix et préparation individuelle (10 min).} Dans ton groupe de quatre, chaque élève choisit un sport ou un champion différent (par exemple : le football, le basketball, la course, le volleyball ; ou un champion comme Samuel Eto'o, Rigobert Song, Lionel Messi, Serena Williams). Prépare au brouillon au moins quatre phrases : \cle{deux comparatifs}, \cle{un superlatif} et \cle{une expression d'opinion}.
\item \textbf{Le mini-débat oral (10 min).} Chacun présente son sport en une minute maximum. Les autres membres du groupe écoutent et notent les arguments. Il est interdit de lire : tu peux seulement regarder tes notes.
\item \textbf{Confrontation des arguments.} Pose des questions à tes camarades avec des comparatifs : \esp{¿Por qué dices que el fútbol es más emocionante que el voleibol?} Réponds en justifiant ton opinion.
\item \textbf{Rédaction collective (15 min).} Le groupe rédige ensemble l'article de \cle{80 mots} : \esp{« El mejor deporte para los jóvenes cameruneses »}. L'article doit comparer les quatre sports défendus et conclure par un superlatif.
\item \textbf{Affichage et vote.} Les articles sont affichés au mur. La classe circule, lit et vote pour le meilleur argumentaire. Chaque vote doit être justifié à l'oral par une phrase avec un comparatif ou un superlatif : \esp{Voto por el grupo 3 porque su artículo es más claro y más convincente.}
\end{enumerate}

\begin{attention}[Pièges à éviter pendant le débat]
\begin{itemize}
\item Ne dis pas \esp{más mejor} : c'est un double marquage incorrect. Dis \esp{mejor} seul ou \esp{el mejor}.
\item Attention à l'accord : \esp{el deporte más divertido}, mais \esp{la actividad más divertida}.
\item \esp{Tan... como} exige un adjectif entre les deux : \esp{el voleibol es tan divertido como el fútbol}, jamais \esp{tan divertido que}.
\item Après \esp{creo que} et \esp{pienso que}, utilise l'indicatif : \esp{Creo que el fútbol es el mejor deporte.}
\item N'oublie pas les accents : \esp{más}, \esp{también}, \esp{emoción}, \esp{opinión}.
\end{itemize}
\end{attention}

\courssub{Ressources à mobiliser}

\begin{propriete}[Rappel : comparatifs et superlatifs]
\begin{center}
\begin{tabularx}{\linewidth}{|X|X|X|}
\hline
\rowcolor{popblueL} Structure & Ejemplo & Traduction \\
\hline
\esp{más + adj. + que} & \esp{El fútbol es más popular que el tenis.} & Le foot est plus populaire que le tennis. \\
\hline
\esp{menos + adj. + que} & \esp{La natación es menos cara que el golf.} & La natation est moins chère que le golf. \\
\hline
\esp{tan + adj. + como} & \esp{El voleibol es tan divertido como el baloncesto.} & Le volley est aussi amusant que le basket. \\
\hline
\esp{el/la más + adj.} & \esp{El fútbol es el deporte más practicado.} & Le foot est le sport le plus pratiqué. \\
\hline
\esp{el mejor / la mejor} & \esp{Para mí, la carrera es la mejor.} & Pour moi, la course est la meilleure. \\
\hline
\end{tabularx}
\end{center}
Comparatifs irréguliers : \esp{bueno $\rightarrow$ mejor} ; \esp{malo $\rightarrow$ peor} ; \esp{grande $\rightarrow$ mayor} (âge) ; \esp{pequeño $\rightarrow$ menor} (âge).
\end{propriete}

\begin{aretenir}[Expressions d'opinion et lexique utile]
Opinion : \esp{creo que}, \esp{pienso que}, \esp{en mi opinión}, \esp{para mí}, \esp{me parece que}, \esp{estoy de acuerdo}, \esp{no estoy de acuerdo}.\\
Lexique : \esp{el equipo} (l'équipe), \esp{el partido} (le match), \esp{entrenar} (s'entraîner), \esp{ganar / perder} (gagner / perdre), \esp{la salud} (la santé), \esp{emocionante} (passionnant), \esp{divertido} (amusant), \esp{barato} (bon marché), \esp{popular} (populaire), \esp{el campeón / la campeona} (le champion / la championne).
\end{aretenir}

\courssub{Proposition de production et corrigé commenté}

\begin{exoresolu}[Article modèle : El mejor deporte para los jóvenes cameruneses]
Rédige avec ton groupe un article de 80 mots comparant quatre sports et concluant par un superlatif.
\tcblower
\textbf{Production modèle (environ 80 mots) :}

\esp{En nuestro liceo de Yaoundé, los jóvenes practican muchos deportes. El fútbol es más popular que el baloncesto, y para muchos es el deporte más emocionante del mundo. Sin embargo, la carrera es más barata que el fútbol: no necesitas equipo ni campo. El voleibol es tan divertido como el baloncesto y menos peligroso que el fútbol. En nuestra opinión, todos los deportes son buenos para la salud, pero creemos que el fútbol es el mejor deporte para los jóvenes cameruneses porque une a todo el país.}

\textbf{Traduction :} « Dans notre lycée de Yaoundé, les jeunes pratiquent beaucoup de sports. Le football est plus populaire que le basketball, et pour beaucoup c'est le sport le plus passionnant du monde. Cependant, la course est moins chère que le football : tu n'as besoin ni d'équipement ni de terrain. Le volleyball est aussi amusant que le basketball et moins dangereux que le football. À notre avis, tous les sports sont bons pour la santé, mais nous pensons que le football est le meilleur sport pour les jeunes Camerounais parce qu'il unit tout le pays. »

\textbf{Commentaire des choix de langue :}
\begin{itemize}
\item Deux comparatifs d'inégalité : \esp{más popular que}, \esp{más barata que} (accord féminin avec \esp{la carrera}).
\item Un comparatif d'égalité : \esp{tan divertido como}, bien construit avec l'adjectif entre les deux termes.
\item Deux superlatifs : \esp{el deporte más emocionante del mundo} (note l'usage de \esp{de} après le superlatif, et non \esp{en}) et \esp{el mejor deporte}, superlatif irrégulier sans \esp{más}.
\item Expressions d'opinion variées : \esp{en nuestra opinión}, \esp{creemos que}, suivies de l'indicatif.
\item Connecteurs pour structurer : \esp{sin embargo} (cependant), \esp{pero} (mais), \esp{porque} (parce que).
\item Vocabulaire du thème : \esp{practicar deportes}, \esp{la salud}, \esp{equipo}, \esp{campo}.
\end{itemize}
\end{exoresolu}

\begin{exemplebox}[Exemples d'interventions orales pour le débat]
\esp{Defiendo el baloncesto porque es más rápido que el fútbol y, en mi opinión, es el deporte más espectacular.} « Je défends le basket parce qu'il est plus rapide que le foot et, à mon avis, c'est le sport le plus spectaculaire. »\\
\esp{No estoy de acuerdo: la carrera es menos cara que el baloncesto y tan buena para la salud como la natación.} « Je ne suis pas d'accord : la course est moins chère que le basket et aussi bonne pour la santé que la natation. »\\
\esp{Para mí, Samuel Eto'o es el mejor campeón camerunés de todos los tiempos.} « Pour moi, Samuel Eto'o est le meilleur champion camerounais de tous les temps. »
\end{exemplebox}

\courssub{Bilan de l'activité}

Cette activité t'a permis de réinvestir l'ensemble des acquis de la leçon dans une tâche finale authentique et motivante. Vérifie ta progression :

\begin{itemize}
\item Je sais \textbf{comparer} deux sports avec \esp{más... que}, \esp{menos... que} et \esp{tan... como}, sans oublier l'accord de l'adjectif.
\item Je sais former le \textbf{superlatif} avec \esp{el/la más + adjectif} et employer les formes irrégulières \esp{el mejor}, \esp{el peor}.
\item Je sais \textbf{exprimer mon opinion} et réagir à celle des autres : \esp{creo que}, \esp{en mi opinión}, \esp{estoy de acuerdo}, \esp{no estoy de acuerdo}.
\item Je sais \textbf{rédiger un court article} structuré (présentation, comparaison, conclusion avec superlatif) d'environ 80 mots.
\item Je sais \textbf{prendre la parole} devant la classe pendant une minute, avec des notes mais sans lire.
\end{itemize}

\begin{experience}[Pour aller plus loin : le micro-trottoir sportif]
Avec un camarade, prépare une interview pour la radio scolaire : l'un joue le journaliste, l'autre un jeune sportif camerounais. Le journaliste pose trois questions avec des comparatifs (\esp{¿Es el fútbol más difícil que el voleibol? ¿Cuál es el mejor equipo de tu barrio?}) et le sportif répond en défendant ses choix. Enregistre l'interview ou présente-la devant la classe.
\end{experience}

\begin{savaistu}[Le sport, langue universelle]
Le football est une passion partagée par tout le monde hispanophone : en Espagne, en Argentine, au Mexique et en Guinée équatoriale, seul pays africain hispanophone, on commente les matchs dans la même langue qu'à Yaoundé ou à Douala quand les Lions Indomptables jouent. Connaître le vocabulaire du sport en espagnol, c'est pouvoir parler avec des millions de jeunes sur les cinq continents : \esp{¿Cuál es tu deporte favorito?}
\end{savaistu}

\newpage

\coursec{Méthodes et exercices résolus}

\coursec{El deporte, el ocio y el bienestar (continuación) : expresión oral y escrita}

\courssub{Texte support : diálogo después del partido}

\begin{textebox}[Diálogo : Después del partido]
\esp{— ¡Qué partido, Andrés! Nuestro equipo ha ganado dos a uno contra el colegio de Bafoussam.}\\
\esp{— Sí, Lucía, el segundo tiempo fue mejor que el primero. El delantero marcó el gol más bonito del campeonato.}\\
\esp{— Es verdad. Pero el árbitro fue peor que el del año pasado.}\\
\esp{— Tienes razón. Mañana tenemos entrenamiento a las cinco.}
\end{textebox}

\begin{definition}[Traduction et glossaire]
« Quel match, Andrés ! Notre équipe a gagné deux à un contre le collège de Bafoussam. — Oui, Lucie, la deuxième mi-temps a été meilleure que la première. L'attaquant a marqué le plus beau but du championnat. — C'est vrai. Mais l'arbitre était pire que celui de l'an dernier. — Tu as raison. Demain nous avons entraînement à cinq heures. »

\textbf{Vocabulaire clé :} \esp{el partido} (le match), \esp{ganar} (gagner), \esp{el marcador / el resultado} (le score), \esp{dos a uno} (deux à un), \esp{el primer / segundo tiempo} (la 1ère / 2ème mi-temps), \esp{el delantero} (l'attaquant), \esp{marcar un gol} (marquer un but), \esp{el campeonato} (le championnat), \esp{el árbitro} (l'arbitre), \esp{el entrenamiento} (l'entraînement), \esp{mejor / peor} (meilleur / pire — comparatifs irréguliers).
\end{definition}

\courssub{Léxique : el deporte, el ocio y el bienestar}

\begin{exemplebox}[Vocabulaire thématique — Sports et loisirs]
\begin{center}
\begin{tabularx}{\linewidth}{|X|X|X|}
\hline
\rowcolor{popblueL} \textbf{Français} & \textbf{Español} & \textbf{Remarque} \\
\hline
le football / le soccer & \esp{el fútbol} & accent sur \emph{ú} ; \esp{jugar \textbf{al} fútbol} \\
le basketball & \esp{el baloncesto} / \esp{el básquet} &  \\
le volleyball & \esp{el voleibol} / \esp{el vóley} & \\
la natation & \esp{la natación} & accent sur \emph{ó} \\
l'athlétisme & \esp{el atletismo} & \\
le cyclisme & \esp{el ciclismo} & \\
le yoga & \esp{el yoga} & \\
la marche / la randonnée & \esp{el senderismo} / \esp{caminar} & \\
le stade & \esp{el estadio} & \\
le terrain (de jeu) & \esp{el campo} / \esp{la cancha} & \emph{cancha} : Am. Lat. \\
l'entraînement & \esp{el entrenamiento} & \\
l'arbitre & \esp{el árbitro / la árbitra} & accent sur \emph{á} \\
le score / le résultat & \esp{el marcador} / \esp{el resultado} & \\
faire match nul & \esp{empatar} & \\
gagner / perdre & \esp{ganar} / \esp{perder} & \\
le spectateur / la spectatrice & \esp{el espectador} / \esp{la espectadora} & \\
le champion / la championne & \esp{el campeón} / \esp{la campeona} & accent sur \emph{ó} \\
la médaille & \esp{la medalla} & \\
la forme (physique) & \esp{la forma} & \esp{estar en forma} \\
le bien-être & \esp{el bienestar} & \\
les loisirs & \esp{el ocio} / \esp{los pasatiempos} & \\
\hline
\end{tabularx}
\end{center}
\end{exemplebox}

\begin{savaistu}[Le sport au Cameroun et dans l'espace hispanophone]
Au Cameroun, le \esp{fútbol} est roi : les \esp{Indomptables} (Lions Indomptables) unissent le pays. Samuel Eto'o, Roger Milla, Vincent Aboubakar sont des \esp{leyendas}. On joue sur les \esp{canchas} de quartier (terrains en terre battue) ou au stade d'Olembé. En Guinée équatoriale (seul pays africain hispanophone), le football domine aussi (équipe nationale : \esp{Nzalang Nacional}). En Espagne, \esp{La Liga} (Real Madrid, Barça) et la \esp{Roja} (championne du monde 2010) sont des références. En Amérique latine, le football est une religion (Argentine, Brésil, Uruguay, Mexique) ; le \esp{béisbol} domine aux Caraïbes (Cuba, Rép. Dominicaine, Venezuela) ; le \esp{basque pelota} (\esp{pelota vasca}) est typique du Pays basque.
\end{savaistu}

\courssub{Grammaire 1 : Les comparatifs (comparativos)}

\begin{propriete}[Règle : Les trois degrés de comparaison]
En espagnol, on compare à l'aide de structures fixes. L'adjectif (ou l'adverbe) s'accorde avec le \textbf{premier terme} de la comparaison.

\begin{center}
\begin{tabularx}{\linewidth}{|X|X|X|}
\hline
\rowcolor{popblueL} \textbf{Type} & \textbf{Structure} & \textbf{Exemple traduit} \\
\hline
Supériorité & \esp{más + adj./adv. + que} & \esp{El fútbol es \textbf{más popular que} el baloncesto.} (Le foot est plus populaire que le basket.) \\
\hline
Infériorité & \esp{menos + adj./adv. + que} & \esp{El yoga es \textbf{menos peligroso que} el ciclismo.} (Le yoga est moins dangereux que le cyclisme.) \\
\hline
Égalité (adj./adv.) & \esp{tan + adj./adv. + como} & \esp{La natación es \textbf{tan cansada como} el atletismo.} (La natation est aussi fatigante que l'athlétisme.) \\
\hline
Égalité (nom) & \esp{tanto/a/os/as + nom + como} & \esp{Tengo \textbf{tanta afición como} tú.} (J'ai autant de passion que toi.) \\
\hline
\end{tabularx}
\end{center}

\textbf{Accord de \esp{tanto} :} \esp{tanto tiempo}, \esp{tanta gente}, \esp{tantos goles}, \esp{tantas medallas}.
\textbf{Préposition :} Toujours \esp{que} (sup./inf.) ou \esp{como} (égalité). Jamais \esp{de}.
\end{propriete}

\begin{propriete}[Comparatifs irréguliers (obligatoires)]
Ces quatre adjectifs/adverbes n'acceptent \textbf{jamais} \esp{más} ni \esp{menos} devant eux.

\begin{center}
\begin{tabular}{|l|l|l|l|}
\hline
\rowcolor{popblueL} \textbf{Positif} & \textbf{Comparatif} & \textbf{Sens} & \textbf{Exemple} \\
\hline
\esp{bueno / bien} & \esp{mejor} & meilleur / mieux & \esp{Juego \textbf{mejor que} ayer.} \\
\hline
\esp{malo / mal} & \esp{peor} & pire / plus mal & \esp{Este equipo es \textbf{peor que} el otro.} \\
\hline
\esp{grande} & \esp{mayor} & plus grand / plus âgé & \esp{Mi hermano es \textbf{mayor que} yo.} \\
\hline
\esp{pequeño} & \esp{menor} & plus petit / plus jeune & \esp{Este estadio es \textbf{menor que} el de Yaoundé.} \\
\hline
\end{tabular}
\end{center}

\emph{Attention :} \esp{mayor/menor} s'emploient pour l'âge et la taille ; pour « plus grand » (taille physique d'un objet), \esp{más grande} est possible mais \esp{mayor} est fréquent.
\end{propriete}

\begin{methode}[Comparer deux sports ou deux loisirs — Démarche en 4 étapes]
\begin{enumerate}
\item \textbf{Analyser} : Quel type de comparaison ? (Supériorité, infériorité, égalité). Sur quoi porte-t-elle ? (Qualité = adjectif/adverbe ; Quantité = nom ; Action = verbe).
\item \textbf{Sélectionner la structure} : \esp{más/menos... que} | \esp{tan... como} | \esp{tanto... como}.
\item \textbf{Vérifier l'irrégularité} : Si l'adjectif est \esp{bueno, malo, grande, pequeño} $\rightarrow$ utiliser \esp{mejor, peor, mayor, menor}.
\item \textbf{Accorder et écrire} : L'adjectif s'accorde avec le sujet (1er terme). Vérifier \esp{que / como} et les accents.
\end{enumerate}
\end{methode}

\begin{exoresolu}[Entraînement : Comparatifs en contexte]
Énoncé : Complète avec la forme comparative correcte (\esp{más… que}, \esp{menos… que}, \esp{tan… como}, \esp{tanto… como}, \esp{mejor}, \esp{peor}, \esp{mayor}, \esp{menor}).\par
1. En Camerún, el fútbol es \dots{} popular \dots{} el balonmano. (supériorité)\par
2. Para mí, nadar es \dots{} divertido \dots{} correr. (égalité)\par
3. Mi hermana tiene \dots{} medallas \dots{} yo. (égalité de quantité — \emph{elle a autant de médailles que moi})\par
4. El maratón es \dots{} duro \dots{} los 100 metros. (supériorité)\par
5. Mi abuelo camina \dots{} (bien) que yo. (supériorité sur adverbe)\par
6. El árbitro de ayer fue \dots{} (malo) que el de la semana pasada.\par
7. Mi entrenador es \dots{} (grande) que el tuyo. (plus âgé)\par
\tcblower
Solution détaillée :\par
1. \esp{más popular que} (adj. régulier, supériorité).\par
2. \esp{tan divertido como} (adj., égalité).\par
3. \esp{tantas medallas como} (\esp{medallas} fém. pl. $\rightarrow$ \esp{tantas}).\par
4. \esp{más duro que} (adj., supériorité).\par
5. \esp{mejor que} (irrégulier \esp{bien $\rightarrow$ mejor}).\par
6. \esp{peor que} (irrégulier \esp{malo $\rightarrow$ peor}).\par
7. \esp{mayor que} (irrégulier \esp{grande $\rightarrow$ mayor} pour l'âge).\par
\textbf{Traductions :} 1. Au Cameroun, le foot est plus populaire que le handball. 2. Pour moi, nager est aussi amusant que courir. 3. Ma sœur a autant de médailles que moi. 4. Le marathon est plus dur que le 100 m. 5. Mon grand-père marche mieux que moi. 6. L'arbitre d'hier était pire que celui de la semaine dernière. 7. Mon entraîneur est plus âgé que le tien.
\end{exoresolu}

\courssub{Grammaire 2 : Les superlatifs (superlativos)}

\begin{propriete}[Superlatif relatif : le plus / le moins \emph{d'un groupe}]
Structure : \esp{artículo + (nombre) + más / menos + adjetivo + \textbf{de}}.
\begin{itemize}
\item Le complément du groupe (le « dans » ou « de » français) s'introduit par \textbf{\esp{de}} (jamais \esp{en}).
\item Si l'adjectif est placé \textbf{après le nom}, l'article se répète : \esp{el jugador \textbf{el} más rápido \textbf{del} equipo} (ou \esp{el jugador más rápido del equipo}).
\item Contractions obligatoires : \esp{de + el = del} ; \esp{de + los = de los} (pas de contraction), \esp{a + el = al}.
\end{itemize}

\begin{center}
\begin{tabularx}{\linewidth}{|X|X|}
\hline
\rowcolor{popblueL} \textbf{Exemple} & \textbf{Traduction} \\
\hline
\esp{El fútbol es \textbf{el deporte más popular de} Camerún.} & Le football est le sport le plus populaire \textbf{du} Cameroun. \\
\hline
\esp{Marta es \textbf{la jugadora menos rápida del} equipo.} & Marta est la joueuse la moins rapide \textbf{de l'}équipe. \\
\hline
\esp{Es \textbf{el gol más bonito del} campeonato.} & C'est le plus beau but \textbf{du} championnat. \\
\hline
\end{tabularx}
\end{center}
\end{propriete}

\begin{propriete}[Superlatif absolu : très, extrêmement (sans comparaison)]
Structure : \esp{adjetivo + \textbf{-ísimo / -ísima / -ísimos / -ísimas}}.
L'adjectif perd son accent écrit s'il en a un ; accord en genre et nombre.

\begin{center}
\begin{tabular}{|l|l|l|l|}
\hline
\rowcolor{popblueL} \textbf{Adjectif} & \textbf{Règle / Changement} & \textbf{Superlatif} & \textbf{Traduction} \\
\hline
\esp{rápido} & ajout \emph{-ísimo} & \esp{rapidísimo} & très rapide \\
\hline
\esp{fácil} & accent tombe, \emph{-ísimo} & \esp{facilísimo} & très facile \\
\hline
\esp{grande} & \emph{e} tombe, \emph{-ísimo} & \esp{grandísimo} & très grand / immense \\
\hline
\esp{feliz} & \emph{z $\rightarrow$ c} + \emph{ísimo} & \esp{felicísimo} & très heureux \\
\hline
\esp{interesante} & ajout \emph{-ísimo} & \esp{interesantísimo} & très intéressant \\
\hline
\esp{caro} & ajout \emph{-ísimo} & \esp{carísimo} & très cher \\
\hline
\end{tabular}
\end{center}

\textbf{Irréguliers (mêmes que comparatifs) :} \esp{el mejor} (le meilleur), \esp{el peor} ( le pire), \esp{el mayor} (le plus âgé/grand), \esp{el menor} (le plus jeune/petit).
\end{propriete}

\begin{methode}[Former et utiliser le superlatif — Démarche en 3 étapes]
\begin{enumerate}
\item \textbf{Identifier} : Relatif (comparaison dans un groupe : « le plus... de... ») ou Absolu (intensité forte : « très... », « extrêmement... »).
\item \textbf{Construire} :
\begin{itemize}
\item Relatif : \esp{el/la/los/las + (nom) + más/menos + adj. + de} (attention \esp{del}, \esp{de la}, \esp{de los}, \esp{de las}).
\item Absolu : Adjectif (masc. sg.) + \esp{-ísimo/a(s)} (vérifier orthographe : \esp{z$\rightarrow$c}, chute \emph{e} final, perte accent).
\end{itemize}
\item \textbf{Placer l'adjectif} : Avant le nom (souvent pour subjectif/émotion : \esp{el mejor jugador}) ou après (souvent pour classificatif/objectif : \esp{el jugador más rápido}). Au superlatif relatif, l'adjectif suit souvent le nom.
\end{enumerate}


\end{methode}

\begin{exoresolu}[Superlatifs : du relatif à l'absolu]
Énoncé : Traduis en espagnol.\par
1. Samuel Eto'o est le meilleur joueur de l'histoire du Cameroun.\par
2. Le marathon est l'épreuve la plus difficile des Jeux Olympiques.\par
3. Mon petit frère est le plus jeune de l'équipe.\par
4. Ce stade (stadio) est immense / très grand. (superlatif absolu)\par
5. La finale était extrêmement intéressante. (superlatif absolu)
\tcblower
Solution détaillée :\par
1. \esp{Samuel Eto'o es \textbf{el mejor jugador de la historia de} Camerún.} (Irrégulier \esp{mejor} + \esp{de}).\par
2. \esp{El maratón es \textbf{la prueba más difícil de} los Juegos Olímpicos.} (Relatif : \esp{la prueba más difícil de}... Accents : \esp{maratón, difícil}.)\par
3. \esp{Mi hermanito es \textbf{el menor del} equipo.} (Irrégulier \esp{menor} + contraction \esp{del}).\par
4. \esp{Este estadio es \textbf{grandísimo}.} (Absolu : \esp{grande $\rightarrow$ grandísimo}).\par
5. \esp{La final fue \textbf{interesantísima}.} (Absolu : \esp{interesante $\rightarrow$ interesantísima}, accord fém. sg.).\par
\textbf{Piège récurrent :} « du / de la / des » = \esp{de / de la / de los / de las} (jamais \esp{en}).
\end{exoresolu}

\courssub{Compréhension orale et écrite : Exploiter un dialogue sportif}

\begin{methode}[Répondre aux questions sur un dialogue — Méthode Bac]
\begin{enumerate}
\item \textbf{Lire les questions} avant le texte (repérer : \esp{Quién? Qué? Cuándo? Dónde? Por qué? Cuál?}).
\item \textbf{Identifier la situation} : Qui parle ? (noms, relation). Où ? (stade, école, parc). Quand ? (après match, avant entraînement).
\item \textbf{Surligner les mots-clés} : \esp{el marcador, ganar, perder, empatar, el árbitro, el entrenador, la afición, lesionado}.
\item \textbf{Répondre en phrases complètes} en espagnol, en reprenant le vocabulaire de la question. Ne pas recopier bêtement le texte.
\item \textbf{Opinion / Inférence} : Pour \esp{¿Qué piensa...?}, citer l'extrait exact : \esp{Andrés dice que... / Según Andrés...}.
\end{enumerate}
\end{methode}

\begin{exoresolu}[Compréhension détaillée : « Después del partido »]
Texte (repris de la boîte Texte support) :\\
\esp{— ¡Qué partido, Andrés! Nuestro equipo ha ganado dos a uno contra el colegio de Bafoussam.}\\
\esp{— Sí, Lucía, el segundo tiempo fue mejor que el primero. El delantero marcó el gol más bonito del campeonato.}\\
\esp{— Es verdad. Pero el árbitro fue peor que el del año pasado.}\\
\esp{— Tienes razón. Mañana tenemos entrenamiento a las cinco.}

Questions :\par
1. ¿Quiénes son los interlocutores y dónde están?\par
2. ¿Cuál es el resultado del partido y contra quién jugaron?\par
3. ¿Qué opinión tiene Andrés del segundo tiempo? Justifica con el texto.\par
4. ¿Qué superlativo se usa para describir el gol? Traduce.\par
5. ¿Cuándo y qué tienen mañana?
\tcblower
Solution détaillée :\par
1. \esp{Los interlocutores son Lucía y Andrés. Están hablando después de un partido (probablemente en el estadio o en el colegio).}\par
2. \esp{El resultado es dos a uno (2-1). Jugaron contra el colegio de Bafoussam.}\par
3. \esp{Andrés piensa que el segundo tiempo fue mejor que el primero.} (Texte : \esp{« el segundo tiempo fue mejor que el primero »}).\par
4. \esp{El superlativo es « el gol más bonito del campeonato » = « le plus beau but du championnat ».}\par
5. \esp{Mañana tienen entrenamiento a las cinco (de la tarde).}
\end{exoresolu}

\begin{experience}[Activité orale : Interview « Après-match »]
\textbf{Consigne :} En binôme. L'un est journaliste (Radio Sport Yaoundé), l'autre est le capitaine de l'équipe victorieuse (ou le coach).
\textbf{Préparation (5 min) :} Le journaliste prépare 5 questions (\esp{¿Cómo fue el partido? ¿Cuál fue la clave? ¿Qué opina del árbitro? ¿El mejor momento? ¿Próximo reto?}). Le capitaine prépare ses réponses en utilisant \textbf{au moins 2 comparatifs et 1 superlatif}.
\textbf{Jeu de rôle (3 min) :} Enregistrement possible (application OnBuch+). Échange naturel, tutoiement ou vouvoiement selon contexte.
\textbf{Retour classe :} Écoute de 2-3 paires. Correction collective : accords, \esp{mejor/peor}, \esp{de} au superlatif, prononciation \esp{fútbol, campeón, árbitro}.
\end{experience}

\courssub{Traduction : Thème et Version (Sport et Loisirs)}

\begin{methode}[Traduire du français vers l'espagnol (Thème) — Check-list]
\begin{enumerate}
\item \textbf{Temps verbal} : Passé composé $\rightarrow$ \esp{Pretérito perfecto} (lien au présent) ou \esp{Indefinido} (action passée close). Imparfait $\rightarrow$ \esp{Imperfecto} (description, habitude).
\item \textbf{Verbes « sport »} : \esp{Jugar \textbf{al} + deporte} (foot, basket...) ; \esp{Practicar / Hacer + deporte} (natation, athlétisme, sport générique) ; \esp{Entrenar} (s'entraîner).
\item \textbf{Faux amis} : \esp{Asistir a} = assister à (être présent) $\neq$ aider (\esp{ayudar}). \esp{Una prueba} = une épreuve (sportive) $\neq$ un test scolaire seulement. \esp{Un récord} = un record.
\item \textbf{Structures comparatives/superlatives} : Vérifier \esp{más... que}, \esp{tan... como}, \esp{el más... de}, \esp{-ísimo}.
\item \textbf{Grammaire fine} : \esp{Ser/Estar} (\esp{es deportista} / \esp{está en forma}), \esp{Por/Para} (\esp{juega por diversión, entrena para ganar}), Subjonctif si opinion/émotion (\esp{Me gusta que juegue bien}).
\item \textbf{Relecture} : Accents (\esp{fútbol, campeón, maratón, rápído $\rightarrow$ rápido}), \esp{¿ ¡}, accords, contractions \esp{al, del}.
\end{enumerate}


\end{methode}

\begin{exoresolu}[Thème : Phrases à traduire en espagnol]
Énoncé : Traduis en espagnol.\par
1. Hier, nous avons regardé un match de football très excitant au stade de Yaoundé.\par
2. Ma cousine pratique la natation depuis trois ans ; elle nage mieux que son frère.\par
3. Pour moi, le basketball est moins fatigant que le football.\par
4. C'est le joueur le plus rapide de l'équipe nationale.\par
5. Nous allons nous entraîner demain matin pour être en forme.
\tcblower
Solution détaillée :\par
1. \esp{Ayer vimos (ou \textbf{hemos visto}) un partido de fútbol \textbf{emocionantísimo} (ou \textbf{muy emocionante}) en el estadio de Yaoundé.} (\esp{ver} un match ; \esp{emocionante} $\rightarrow$ \esp{emocionantísimo} ou \esp{más emocionante}).\par
2. \esp{Mi prima practica (ou \textbf{nada}) la natación desde hace tres años; nada \textbf{mejor que} su hermano.} (\esp{desde hace} + présent = « depuis » ; \esp{mejor} irrégulier).\par
3. \esp{Para mí, el baloncesto es \textbf{menos cansado que} el fútbol.} (Comparatif infériorité).\par
4. \esp{Es \textbf{el jugador más rápido de la} selección nacional.} (Superlatif relatif + \esp{de la}).\par
5. \esp{Mañana por la mañana vamos a entrenar para \textbf{estar en forma}.} (\esp{ir a + inf.} futur proche ; \esp{para} + but ; \esp{estar en forma} locution figée).
\end{exoresolu}

\begin{exoresolu}[Version : Comprendre un texte espagnol]
Énoncé : Traduis en français.\par
\esp{« Mi deporte favorito es el tenis. Juego al tenis dos veces por semana en el club de mi ciudad. Mi entrenador dice que mi revés ha mejorado mucho, pero mi saque es peor que el de mi rival. El próximo torneo es el más importante del año. Quiero ganar para ser el campeón. »}\par
\tcblower
Solution :\par
« Mon sport préféré est le tennis. Je joue au tennis deux fois par semaine au club de ma ville. Mon entraîneur dit que mon revers s'est beaucoup amélioré, mais mon service est pire que celui de mon rival. Le prochain tournoi est le plus important de l'année. Je veux gagner pour être le champion. »
\textbf{Points grammaticaux :} \esp{juego al tenis} (contraction \esp{al}), \esp{ha mejorado} (parfait), \esp{peor que} (comparatif irrégulier), \esp{el más importante del año} (superlatif relatif + \esp{del}), \esp{para ser} (but).
\end{exoresolu}

\begin{exercice}[Entraînement Thème / Version]
\textbf{Thème (Fr $\rightarrow$ Es) :}
1. Le sport est bon pour la santé physique et mentale.
2. Mon équipe a perdu le match trois à zéro (3-0) contre Douala.
3. Courir est plus difficile que marcher, mais c'est meilleur pour le cœur.
4. Elle est la meilleure athlète de la région.
5. Nous avons assisté à une finale incroyable samedi dernier.

\textbf{Version (Es $\rightarrow$ Fr) :}
\esp{« Los jóvenes de mi barrio practican fútbol en la calle. No tienen campo, pero tienen mucha pasión. El mejor jugador del barrio se llama Kevin. Es más rápido que los demás y tira penaltis perfectos. El domingo hay un torneo: el premio es un balón nuevo. Todos quieren ganar. »}
\end{exercice}

\begin{corrige}[Exercice Thème / Version]
\textbf{Thème :}
1. \esp{El deporte es bueno para la salud física y mental.} (\esp{ser} pour caractéristique intrinsèque).
2. \esp{Nuestro equipo perdió el partido tres a cero contra (el equipo de) Douala.} (\esp{perder} + score \esp{tres a cero}).
3. \esp{Correr es más difícil que caminar, pero es \textbf{mejor} para el corazón.} (\esp{mejor} irrégulier, pas \esp{más bueno}).
4. \esp{Ella es \textbf{la mejor atleta de la} región.} (Superlatif irrégulier \esp{mejor} + \esp{de la}).
5. \esp{Asistimos a una final \textbf{increíble} (ou \textbf{increíblemente} emocionante) el sábado pasado.} (\esp{asistir a} = assister à ; \esp{increíble} adjectif accordé avec \esp{final}).

\textbf{Version :}
« Les jeunes de mon quartier jouent au football dans la rue. Ils n'ont pas de terrain, mais ils ont beaucoup de passion. Le meilleur joueur du quartier s'appelle Kevin. Il est plus rapide que les autres et tire des penalties parfaits. Dimanche il y a un tournoi : le prix est un ballon neuf. Tous veulent gagner. »
\textbf{Vocabulaire :} \esp{la calle} (la rue), \esp{el campo} (le terrain), \esp{la pasión} (la passion), \esp{los demás} (les autres), \esp{tirar penaltis} (tirer des penalties), \esp{el premio} (le prix/récompense), \esp{un balón} (un ballon).
\end{corrige}

\courssub{Production écrite guidée : « Mi deporte favorito » / « Mi campeón favorito » (80 mots)}

\begin{methode}[Rédiger un texte court (75–90 mots) — Structure et exigences]
\begin{enumerate}
\item \textbf{Plan en 3 paragraphes (ou 3 blocs logiques)} :
\begin{itemize}
\item \textbf{Introduction (15–20 mots) :} Nom du sport / du champion. Fréquence, lieu, partenaires. \esp{Mi deporte favorito es... / Mi campeón favorito es... Practico... desde hace... / Juego con... en...}
\item \textbf{Développement (40–50 mots) :} Raisons (plaisir, santé, valeurs), détails concrets, \textbf{comparatif obligatoire} (vs autre sport / autre joueur), \textbf{superlatif obligatoire} (relatif ou absolu). Connecteurs : \esp{porque, además, también, pero, sin embargo}.
\item \textbf{Conclusion (15–20 mots) :} Opinion personnelle, rêve, superlatif final. \esp{Para mí, es el mejor... / Quiero ser como... / Es lo máximo.}
\end{itemize}
\item \textbf{Check-list grammaticale (points Bac)} : \checkmark 1 comparatif (\esp{más/menos/que, tan/como, mejor/peor}). \checkmark 1 superlatif (\esp{el más... de, -ísimo}). \checkmark Verbes variés (\esp{practicar, jugar, entrenar, ganar, correr, admirar, preferir}). \checkmark Connecteurs. \checkmark Accents (\esp{fútbol, campeón, maratón, rápído $\rightarrow$ rápido, difícil}). \checkmark \esp{Ser/Estar}, \esp{Por/Para}.
\item \textbf{Comptage} : Compter les mots (mots séparés par espace). Visée : 80 $\pm$ 10.
\end{enumerate}


\end{methode}

\begin{exoresolu}[Modèle rédigé : « Mi deporte favorito » (niveau A2+/B1)]
Énoncé : Présente ton sport préféré (nom, pratique, raisons, comparaison, superlatif). Environ 80 mots.
\tcblower
Solution détaillée — Modèle :\par
\esp{Mi deporte favorito es el fútbol. Practico este deporte desde hace seis años con mis amigos del barrio, aquí en Yaoundé. Jugamos todos los sábados por la tarde en el campo de tierra del colegio. Me encanta el fútbol porque es \textbf{más divertido que} la natación y porque juego en equipo. \textbf{Además}, es excelente para la salud: corro mucho y \textbf{estoy en forma}. El año pasado, mi equipo ganó el torneo intercolegial y yo marqué el \textbf{gol más bonito del campeonato}. Para mí, el fútbol es \textbf{el mejor deporte del mundo} y Samuel Eto'o es \textbf{el mejor jugador africano de la historia}.}\par
\textbf{Analyse (86 mots) :}
\begin{itemize}
\item \textbf{Intro} : Sport, durée (\esp{desde hace}), lieu, fréquence, partenaires.
\item \textbf{Développement} : Raisons (\esp{porque}), comparatif \esp{más divertido que}, connecteur \esp{además}, \esp{estar en forma}, passé (\esp{ganó, marqué}), superlatif relatif \esp{el gol más bonito del campeonato}.
\item \textbf{Conclusion} : Superlatifs absolus/relatifs \esp{el mejor deporte del mundo, el mejor jugador africano de la historia}.
\item \textbf{Grammaire} : \esp{jugar al fútbol} (implicite), \esp{ser/estar} corrects, accents (\esp{fútbol, campeón, más, mejor}).
\end{itemize}


\end{exoresolu}

\begin{exercice}[À vous de jouer : Production écrite]
\textbf{Sujet A :} Rédige un texte de 80 mots environ pour présenter \textbf{ton sport préféré}. Suis le plan et la check-list de la Méthode 5.
\textbf{Sujet B :} Rédige un texte de 80 mots environ pour présenter \textbf{ton champion(ne) préféré(e)} (footballeur, basketteur, athlète, local ou international). Inclus : nom, sport, pays, qualités, comparatif avec un autre, superlatif, ton admiration.
\textbf{Contraintes :} 75–90 mots. Au moins 1 comparatif, 1 superlatif. Relis-toi (accords, accents, \esp{ser/estar}).
\end{exercice}

\begin{attention}[Les 5 erreurs fatales au Bac — Récapitulatif]
\begin{enumerate}
\item \textbf{\esp{más bueno / más malo} $\rightarrow$ INTERDIT.} \cle{Irréguliers} : \esp{bueno/bien $\rightarrow$ mejor} ; \esp{malo/mal $\rightarrow$ peor}. Jamais \esp{más mejor}.
\item \textbf{Superlatif + \esp{en} $\rightarrow$ FAUX.} « Le meilleur \emph{du} monde » = \esp{el mejor \textbf{del} mundo}. Préposition \textbf{\esp{de}} (contractions \esp{del, de la, de los, de las}).
\item \textbf{\esp{Asistir} = Aider $\rightarrow$ FAUX AMI.} \esp{Asistir a un partido} = Assister à un match (spectateur). « Aider » = \esp{Ayudar}.
\item \textbf{Accents oubliés = Faute comptée.} \esp{Fútbol, campeón, maratón, rápido, difícil, interés, jóvenes, sábado, miércoles, entrenador, árbitro, selección, nación}. Apprends-les par cœur.
\item \textbf{\esp{Jugar a el fútbol} $\rightarrow$ FAUX.} Contraction obligatoire : \esp{Jugar \textbf{al} fútbol} (\esp{a + el = al}). Idem : \esp{Ir al estadio}. « Faire du sport » = \esp{Hacer deporte} / \esp{Practicar deporte} (pas \esp{de}).
\end{enumerate}
\end{attention}

\begin{savaistu}[Éthique sportive et valeurs : \esp{Juego Limpio}]
Au-delà de la grammaire, le programme MINESEC insiste sur les valeurs : \esp{el juego limpio} (fair-play), \esp{el respeto} (respect adversaire, arbitre, règles), \esp{la superación} (dépassement de soi), \esp{el trabajo en equipo} (esprit d'équipe). Au Cameroun, les \esp{valores olímpicos} sont enseignés via le sport scolaire (FENASCO). En Espagne, la \esp{Real Federación Española} promeut \esp{« Gana con honor, pierde con dignidad »}. Utilise ce vocabulaire pour enrichir tes rédactions orales/écrites : \esp{un deportista ejemplar, competir con honor, respetar al rival}.

\end{savaistu}


\newpage

\coursec{Exercices}

\courssub{Je vérifie mes connaissances}

\begin{exercice}[QCM : le bon comparatif]
Entoure la bonne réponse.
\begin{enumerate}
\item El fútbol es \cle{más} popular \ldots\ el balonmano en Camerún. \hfill a) que \quad b) como \quad c) de
\item Mi hermana nada \ldots\ bien \ldots\ yo. \hfill a) más / que \quad b) tan / como \quad c) menos / de
\item El boxeo es \ldots\ peligroso \ldots\ la natación. \hfill a) tan / como \quad b) menos / que \quad c) más / como
\item Este jugador es \ldots\ del equipo. \hfill a) el más alto \quad b) más alto que \quad c) tan alto
\item El ciclismo es \ldots\ caro \ldots\ el atletismo : ¡la bicicleta cuesta mucho! \hfill a) menos / que \quad b) más / que \quad c) tan / como
\end{enumerate}
\end{exercice}

\begin{exercice}[Vrai ou faux ? Justifie ta réponse]
Dis si les phrases suivantes sont vraies (\emph{vrai}) ou fausses (\emph{faux}) et justifie en français.
\begin{enumerate}
\item \esp{Tan... como} exprime la supériorité.
\item \esp{Mejor} est le comparatif de \esp{bueno}.
\item On dit \esp{más bueno} pour dire « meilleur ».
\item Dans \esp{el más rápido del mundo}, \esp{de} se traduit par « de » et introduit le complément du superlatif.
\item \esp{Peor} signifie « pire / plus mauvais ».
\end{enumerate}
\end{exercice}

\begin{exercice}[Vocabulaire : sport et bien-être]
Complète chaque phrase avec le mot espagnol qui convient, choisi dans la liste : \esp{entrenamiento} — \esp{empate} — \esp{saludable} — \esp{aficionados} — \esp{árbitro}.
\begin{enumerate}
\item El partido terminó con un \ldots\ : dos goles contra dos.
\item Los \ldots\ del equipo camerunés cantan en el estadio de Yaoundé.
\item El \ldots\ pita el final del partido.
\item Hacer deporte es muy \ldots\ para el cuerpo y la mente.
\item Los jugadores van al \ldots\ todos los días a las seis de la mañana.
\end{enumerate}
\end{exercice}

\begin{exercice}[Comparatifs et superlatifs à compléter]
Complète le tableau avec la forme correcte de l'adjectif.
\begin{center}
\begin{tabular}{|l|l|l|}
\hline
\rowcolor{popblueL} Adjectif & Comparatif de supériorité & Superlatif relatif \\
\hline
rápido & más rápido \ldots & el más rápido \ldots \\
\hline
bueno & \ldots & \ldots \\
\hline
malo & \ldots & \ldots \\
\hline
grande & \ldots & \ldots \\
\hline
joven & \ldots & \ldots \\
\hline
\end{tabular}
\end{center}
\end{exercice}

\courssub{Je m'entraîne}

\begin{exercice}[Texte à trous : un match de la CAN]
Complète le texte suivant avec : \esp{más} — \esp{menos} — \esp{tan} — \esp{que} — \esp{como} — \esp{mejor} — \esp{peor} — \esp{el más}.
\begin{textebox}[Un partido de la CAN en Douala]
Ayer vi un partido de la Copa de África en el estadio de Douala. El equipo camerunés jugó \ldots\ que su adversario : corrió \ldots\ y defendió \ldots\ . El delantero camerunés fue \ldots\ rápido \ldots\ todos los jugadores del campo, y marcó dos goles. El portero contrario no fue \ldots\ bueno \ldots\ el nuestro : cometió \ldots\ errores \ldots\ nuestro portero. Para mí, el centrocampista camerunés fue \ldots\ jugador del partido. Al final, Camerún ganó tres a uno y los aficionados celebraron la victoria en las calles de la ciudad.
\end{textebox}
\end{exercice}

\begin{exercice}[Reformulation : du comparatif au superlatif]
Transforme chaque phrase : passe du comparatif au superlatif relatif, ou l'inverse, sans changer le sens. Exemple : \esp{Este jugador es más rápido que los demás.} $\rightarrow$ \esp{Es el más rápido del equipo.}
\begin{enumerate}
\item El fútbol es más popular que los otros deportes en Camerún.
\item La natación es el deporte menos peligroso de la lista.
\item Samuel Eto'o es más famoso que los otros jugadores africanos.
\item El ciclismo es el deporte más caro de mi barrio.
\item Mi hermana es la más joven de su equipo de voleibol.
\end{enumerate}
\end{exercice}

\begin{exercice}[Thème : phrases sur le sport]
Traduis en espagnol.
\begin{enumerate}
\item Le basket-ball est moins populaire que le football au Cameroun.
\item Mon frère court aussi vite que moi.
\item La natation est le meilleur sport pour la santé.
\item Ce joueur est plus fort que son adversaire, mais moins rapide.
\item Le match d'hier était pire que le match de samedi.
\end{enumerate}
\end{exercice}

\begin{exercice}[Appariements : questions et réponses]
Associe chaque question (1-5) à la réponse qui convient (a-e).
\begin{enumerate}
\item \esp{¿Cuál es tu deporte favorito?}
\item \esp{¿Entrenas todos los días?}
\item \esp{¿Quién es el mejor jugador del equipo?}
\item \esp{¿Prefieres el fútbol o el baloncesto?}
\item \esp{¿Quién ganó el partido de ayer?}
\end{enumerate}
\begin{itemize}
\item[a)] \esp{Prefiero el fútbol : es más emocionante que el baloncesto.}
\item[b)] \esp{Mi deporte favorito es el atletismo.}
\item[c)] \esp{Ganó nuestro equipo, dos a cero.}
\item[d)] \esp{No, entreno solo tres veces por semana.}
\item[e)] \esp{Para mí, el mejor es el capitán : corre más que todos.}
\end{itemize}
\end{exercice}

\courssub{Je me prépare au Bac}

\begin{exercice}[Compréhension de l'écrit et questions de langue]
Lis le texte suivant, puis réponds aux questions.
\begin{textebox}[El deporte en mi ciudad]
En mi ciudad, Douala, el deporte es muy importante para los jóvenes. El fútbol es más popular que los otros deportes : los chicos juegan en las calles, en los patios de las escuelas y en los estadios. Muchos sueñan con ser tan famosos como los grandes campeones africanos.

Pero el fútbol no es el único deporte. El baloncesto es cada vez más apreciado, sobre todo entre las chicas, y la natación es tan sana como el atletismo, aunque menos practicada porque hay pocas piscinas. Mi primo, por ejemplo, nada todos los sábados : dice que es el mejor deporte para el cuerpo porque es menos peligroso que el boxeo y más completo que el ciclismo.

El deporte no es solo una competición : es también un placer, una escuela de disciplina y una manera de estar sano. Por eso, los profesores animan a los alumnos a practicar una actividad física regular. En mi instituto, el equipo de fútbol entrena tres veces por semana y el ambiente es fantástico. Para mí, hacer deporte con los amigos es el mejor momento de la semana.
\end{textebox}
\textbf{Questions de compréhension} (réponds en espagnol, avec des phrases complètes) :
\begin{enumerate}
\item ¿Dónde juegan al fútbol los jóvenes de Douala ? Cita dos lugares.
\item ¿Por qué la natación es menos practicada que el fútbol ?
\item Según el primo del narrador, ¿por qué la natación es el mejor deporte para el cuerpo ?
\item Según el texto, el deporte es « una escuela de disciplina ». Explica esta idea con tus palabras.
\end{enumerate}
\textbf{Questions de langue} :
\begin{enumerate}
\item Relève dans le texte deux comparatifs et deux superlatifs, puis traduis-les en français.
\item Trouve dans le texte le contraire de : \esp{menos practicada} ; \esp{peor}.
\item Conjugue les verbes entre parenthèses : \esp{Los chicos (jugar) en las calles} ; \esp{Mi primo (nadar) todos los sábados}.
\end{enumerate}
\end{exercice}

\begin{exercice}[Thème et version]
\textbf{Thème.} Traduis en espagnol :
\begin{quote}
« Le football est plus populaire que le basket au Cameroun, mais Samuel Eto'o est le joueur le plus célèbre du monde africain. Mon équipe s'entraîne tous les jours, et notre capitaine court aussi vite qu'un professionnel. Pour moi, le sport est le meilleur loisir du week-end. »
\end{quote}
\textbf{Version.} Traduis en français :
\begin{textebox}[Para traducir]
La natación es tan sana como el atletismo, pero menos peligrosa que el boxeo ; para mí, es el mejor deporte del verano. Mi hermana nada más rápido que yo, pero yo corro más que ella. Los dos somos los más deportistas de la familia.
\end{textebox}
\end{exercice}

\begin{exercice}[Production écrite et expression orale]
\textbf{A. Production écrite guidée.} Rédige un texte d'environ \cle{80 mots} en espagnol sur le sujet suivant : \esp{Mi campeón favorito} ou \esp{Mi deporte favorito}.
Ta production doit contenir :
\begin{itemize}
\item une présentation du champion ou du sport (nom, pays, discipline) ;
\item au moins \cle{deux comparatifs} (\esp{más... que}, \esp{menos... que} ou \esp{tan... como}) ;
\item au moins \cle{un superlatif} (\esp{el mejor}, \esp{el más...}) ;
\item une conclusion personnelle (pourquoi tu l'aimes).
\end{itemize}
\textbf{B. Expression orale (travail en binôme).} Un élève joue le rôle d'un journaliste sportif, l'autre celui d'un jeune champion camerounais. Le journaliste pose au moins cinq questions sur l'entraînement, les loisirs et le bien-être ; le champion répond en utilisant obligatoirement deux comparatifs et un superlatif. Puis échangez les rôles.
\end{exercice}

\newpage

\coursec{Corrigés des exercices}

\begin{corrige}[Exercice 1 — QCM : le bon comparatif]
\begin{enumerate}
\item \textbf{a) que} : \esp{El fútbol es más popular que el balonmano en Camerún.} « Le football est plus populaire que le handball au Cameroun. » Après \esp{más} + adjectif, le second terme de la comparaison est introduit par \esp{que}.
\item \textbf{b) tan / como} : \esp{Mi hermana nada tan bien como yo.} « Ma sœur nage aussi bien que moi. » \esp{Tan... como} exprime l'égalité.
\item \textbf{b) menos / que} : \esp{El boxeo es menos peligroso que la natación.} « La boxe est moins dangereuse que la natation. » Infériorité : \esp{menos} + adjectif + \esp{que}.
\item \textbf{a) el más alto} : \esp{Este jugador es el más alto del equipo.} « Ce joueur est le plus grand de l'équipe. » Superlatif relatif : article défini + \esp{más} + adjectif.
\item \textbf{b) más / que} : \esp{El ciclismo es más caro que el atletismo.} « Le cyclisme est plus cher que l'athlétisme. » La justification (le vélo coûte cher) impose la supériorité.
\end{enumerate}
\end{corrige}

\begin{corrige}[Exercice 2 — Vrai ou faux ?]
\begin{enumerate}
\item \textbf{Faux.} \esp{Tan... como} exprime l'\emph{égalité} (« aussi... que ») : \esp{Corro tan rápido como tú.} « Je cours aussi vite que toi. » La supériorité s'exprime avec \esp{más... que}.
\item \textbf{Vrai.} \esp{Mejor} est le comparatif irrégulier de \esp{bueno} : \esp{Este equipo es mejor que el otro.} « Cette équipe est meilleure que l'autre. »
\item \textbf{Faux.} On ne dit jamais \esp{más bueno} pour « meilleur » : la forme correcte est \esp{mejor}. (\esp{Más bueno} existe seulement au sens de « plus gentil, plus bon de caractère ».)
\item \textbf{Vrai.} Dans le superlatif relatif, le complément est introduit par \esp{de} : \esp{el más rápido del mundo} « le plus rapide du monde » (\esp{del} = \esp{de} + \esp{el}).
\item \textbf{Vrai.} \esp{Peor} est le comparatif irrégulier de \esp{malo} : « pire, plus mauvais » : \esp{Mi equipo es peor que el tuyo.} « Mon équipe est pire que la tienne. »
\end{enumerate}
\end{corrige}

\begin{corrige}[Exercice 3 — Vocabulaire : sport et bien-être]
\begin{enumerate}
\item \esp{El partido terminó con un \textbf{empate} : dos goles contra dos.} « Le match s'est terminé par un match nul : deux buts partout. »
\item \esp{Los \textbf{aficionados} del equipo camerunés cantan en el estadio de Yaoundé.} « Les supporters de l'équipe camerounaise chantent au stade de Yaoundé. »
\item \esp{El \textbf{árbitro} pita el final del partido.} « L'arbitre siffle la fin du match. » (Attention à l'accent : \esp{árbitro}.)
\item \esp{Hacer deporte es muy \textbf{saludable} para el cuerpo y la mente.} « Faire du sport est très sain pour le corps et l'esprit. »
\item \esp{Los jugadores van al \textbf{entrenamiento} todos los días a las seis de la mañana.} « Les joueurs vont à l'entraînement tous les jours à six heures du matin. »
\end{enumerate}
\end{corrige}

\begin{corrige}[Exercice 4 — Comparatifs et superlatifs à compléter]
\begin{center}
\begin{tabular}{|l|l|l|}
\hline
\rowcolor{popblueL} Adjectif & Comparatif de supériorité & Superlatif relatif \\
\hline
rápido & más rápido que & el más rápido de \\
\hline
bueno & mejor que & el mejor de \\
\hline
malo & peor que & el peor de \\
\hline
grande & mayor que & el mayor de \\
\hline
joven & menor que & el menor de \\
\hline
\end{tabular}
\end{center}
\textbf{Justifications :}
\begin{itemize}
\item \esp{Bueno} et \esp{malo} ont des formes irrégulières : \esp{mejor} / \esp{peor} (jamais \esp{más bueno}, \esp{más malo} au sens de « meilleur / pire »).
\item \esp{Grande} → \esp{mayor} (« plus grand, plus âgé ») ; \esp{joven} → \esp{menor} (« plus jeune »). Ces formes sont invariables en genre : \esp{mi hermana mayor} « ma sœur aînée ».
\item Au superlatif, le complément est introduit par \esp{de} : \esp{el mejor del mundo} « le meilleur du monde ».
\end{itemize}
\end{corrige}

\begin{corrige}[Exercice 5 — Texte à trous : un match de la CAN]
\begin{textebox}[Un partido de la CAN en Douala — corrigé]
Ayer vi un partido de la Copa de África en el estadio de Douala. El equipo camerunés jugó \textbf{mejor} que su adversario : corrió \textbf{más} y defendió \textbf{mejor}. El delantero camerunés fue \textbf{el más} rápido \textbf{de} todos los jugadores del campo, y marcó dos goles. El portero contrario no fue \textbf{tan} bueno \textbf{como} el nuestro : cometió \textbf{más} errores \textbf{que} nuestro portero. Para mí, el centrocampista camerunés fue \textbf{el mejor} jugador del partido. Al final, Camerún ganó tres a uno y los aficionados celebraron la victoria en las calles de la ciudad.
\end{textebox}
\textbf{Remarque :} le mot \esp{menos} n'était pas utilisé ici ; \esp{peor} non plus (le sens impose \esp{mejor} : l'équipe camerounaise a \emph{mieux} joué). Pour « le plus rapide de tous les joueurs », on accepte \esp{el más rápido de todos} (superlatif + complément en \esp{de}).
\textbf{Traduction :} « Hier, j'ai vu un match de la Coupe d'Afrique au stade de Douala. L'équipe camerounaise a mieux joué que son adversaire : elle a couru plus et mieux défendu. L'attaquant camerounais a été le plus rapide de tous les joueurs du terrain et a marqué deux buts. Le gardien adverse n'a pas été aussi bon que le nôtre : il a commis plus d'erreurs que notre gardien. Pour moi, le milieu de terrain camerounais a été le meilleur joueur du match. À la fin, le Cameroun a gagné trois à un et les supporters ont célébré la victoire dans les rues de la ville. »
\end{corrige}

\begin{corrige}[Exercice 6 — Reformulation : du comparatif au superlatif]
\begin{enumerate}
\item \esp{El fútbol es el deporte más popular de Camerún.} « Le football est le sport le plus populaire du Cameroun. »
\item \esp{La natación es menos peligrosa que los otros deportes de la lista.} « La natation est moins dangereuse que les autres sports de la liste. » (Attention à l'accord : \esp{peligrosa}, féminin de \esp{la natación}.)
\item \esp{Samuel Eto'o es el jugador africano más famoso.} / \esp{Samuel Eto'o es el más famoso de los jugadores africanos.} « Samuel Eto'o est le plus célèbre des joueurs africains. »
\item \esp{El ciclismo es más caro que los otros deportes de mi barrio.} « Le cyclisme est plus cher que les autres sports de mon quartier. »
\item \esp{Mi hermana es más joven que las otras jugadoras de su equipo de voleibol.} « Ma sœur est plus jeune que les autres joueuses de son équipe de volley. » (On peut aussi écrire : \esp{es menor que las otras...}.)
\end{enumerate}
\textbf{Méthode :} comparatif \esp{más X que los demás} $\leftrightarrow$ superlatif \esp{el más X de + groupe}. Le sens reste identique ; seule la structure change.
\end{corrige}

\begin{corrige}[Exercice 7 — Thème : phrases sur le sport]
\begin{enumerate}
\item \esp{El baloncesto es menos popular que el fútbol en Camerún.} (Infériorité : \esp{menos... que} ; « basket-ball » = \esp{baloncesto}, pas de calque de « basket ».)
\item \esp{Mi hermano corre tan rápido como yo.} (Égalité : \esp{tan... como} ; « vite » = \esp{rápido}, adverbe invariable ici.)
\item \esp{La natación es el mejor deporte para la salud.} (Superlatif irrégulier : \esp{el mejor}, jamais \esp{el más bueno}.)
\item \esp{Este jugador es más fuerte que su adversario, pero menos rápido.} (Deux comparatifs dans la même phrase : \esp{más... que} puis \esp{menos} + adjectif.)
\item \esp{El partido de ayer fue peor que el partido del sábado.} (Comparatif irrégulier \esp{peor} ; au passé, \esp{fue} ou \esp{era} selon l'aspect — \esp{fue} convient pour un match précis.)
\end{enumerate}
\textbf{Points de vigilance :} ne pas oublier l'article devant les noms de sports (\esp{el baloncesto}, \esp{la natación}) ; ne pas traduire « au Cameroun » par \esp{a Camerún} mais \esp{en Camerún}.
\end{corrige}

\begin{corrige}[Exercice 8 — Appariements : questions et réponses]
\begin{center}
\begin{tabular}{|l|l|}
\hline
\rowcolor{popblueL} Question & Réponse \\
\hline
1. ¿Cuál es tu deporte favorito? & b) Mi deporte favorito es el atletismo. \\
\hline
2. ¿Entrenas todos los días? & d) No, entreno solo tres veces por semana. \\
\hline
3. ¿Quién es el mejor jugador del equipo? & e) Para mí, el mejor es el capitán : corre más que todos. \\
\hline
4. ¿Prefieres el fútbol o el baloncesto? & a) Prefiero el fútbol : es más emocionante que el baloncesto. \\
\hline
5. ¿Quién ganó el partido de ayer? & c) Ganó nuestro equipo, dos a cero. \\
\hline
\end{tabular}
\end{center}
\textbf{Justifications :} la question 1 porte sur le sport favori (réponse b : \esp{mi deporte favorito}) ; la question 2 sur la fréquence de l'entraînement (réponse d : \esp{tres veces por semana}) ; la question 3 sur le meilleur joueur (réponse e : \esp{el mejor es el capitán}) ; la question 4 propose un choix avec \esp{o} (réponse a : \esp{prefiero el fútbol}) ; la question 5 demande le vainqueur (réponse c : \esp{ganó nuestro equipo}).
\end{corrige}

\begin{corrige}[Exercice 9 — Compréhension de l'écrit et questions de langue]
\textbf{Questions de compréhension} (réponses modèles en phrases complètes) :
\begin{enumerate}
\item \esp{Los jóvenes de Douala juegan al fútbol en las calles, en los patios de las escuelas y en los estadios.} « Les jeunes de Douala jouent au football dans les rues, dans les cours des écoles et dans les stades. » (Deux lieux suffisent.)
\item \esp{La natación es menos practicada porque hay pocas piscinas.} « La natation est moins pratiquée parce qu'il y a peu de piscines. »
\item \esp{Según su primo, la natación es el mejor deporte para el cuerpo porque es menos peligrosa que el boxeo y más completa que el ciclismo.} « Selon son cousin, la natation est le meilleur sport pour le corps parce qu'elle est moins dangereuse que la boxe et plus complète que le cyclisme. »
\item Réponse attendue, par exemple : \esp{El deporte es « una escuela de disciplina » porque enseña a respetar las reglas, a entrenar con regularidad y a esforzarse todos los días.} « Le sport est une école de discipline parce qu'il apprend à respecter les règles, à s'entraîner régulièrement et à fournir des efforts chaque jour. »
\end{enumerate}
\textbf{Questions de langue :}
\begin{enumerate}
\item Comparatifs : \esp{más popular que los otros deportes} « plus populaire que les autres sports » ; \esp{tan famosos como los grandes campeones} « aussi célèbres que les grands champions » ; \esp{tan sana como el atletismo} « aussi saine que l'athlétisme » ; \esp{menos peligroso que el boxeo} « moins dangereux que la boxe » ; \esp{más completo que el ciclismo} « plus complet que le cyclisme ». Superlatifs : \esp{el mejor deporte para el cuerpo} « le meilleur sport pour le corps » ; \esp{el mejor momento de la semana} « le meilleur moment de la semaine ». (Deux de chaque suffisent.)
\item Contraire de \esp{menos practicada} : \esp{más popular} (ou \esp{más apreciado / más practicado}) ; contraire de \esp{peor} : \esp{mejor} (\esp{el mejor deporte}).
\item \esp{Los chicos \textbf{juegan} en las calles} (présent, 3\up{e} personne du pluriel, verbe à changement radical o $\rightarrow$ ue) ; \esp{Mi primo \textbf{nada} todos los sábados} (présent, 3\up{e} personne du singulier).
\end{enumerate}
\textbf{Barème indicatif (sur 20) :} compréhension 10 pts (2,5 pts par question : 1 pt pour l'information exacte, 1,5 pt pour la correction de la langue) ; questions de langue 10 pts (4 pts pour les comparatifs/superlatifs relevés et traduits, 2 pts pour les contraires, 4 pts pour les conjugaisons).
\textbf{Points de vigilance :} répondre par des phrases complètes en reprenant les mots de la question ; accorder les adjectifs (\esp{peligrosa} avec \esp{la natación}) ; ne pas confondre \esp{mejor} (comparatif) et \esp{el mejor} (superlatif).
\end{corrige}

\begin{corrige}[Exercice 10 — Thème et version]
\textbf{Thème (proposition de traduction) :}
\begin{quote}
\esp{El fútbol es más popular que el baloncesto en Camerún, pero Samuel Eto'o es el jugador más famoso del mundo africano. Mi equipo entrena todos los días, y nuestro capitán corre tan rápido como un profesional. Para mí, el deporte es el mejor ocio del fin de semana.}
\end{quote}
\textbf{Points de vigilance (thème) :} « s'entraîne » = \esp{entrena} ou \esp{se entrena} (les deux sont acceptés) ; « aussi vite qu'un professionnel » = \esp{tan rápido como un profesional} ; « le joueur le plus célèbre » = \esp{el jugador más famoso} (article répété devant le nom) ; « du monde africain » = \esp{del mundo africano} (contraction \esp{de} + \esp{el} = \esp{del}).
\textbf{Version (traduction) :}
\begin{quote}
« La natation est aussi saine que l'athlétisme, mais moins dangereuse que la boxe ; pour moi, c'est le meilleur sport de l'été. Ma sœur nage plus vite que moi, mais moi je cours plus qu'elle. Nous sommes tous les deux les plus sportifs de la famille. »
\end{quote}
\textbf{Points de vigilance (version) :} \esp{tan sana como} = « aussi saine que » (égalité) ; \esp{más rápido} = « plus vite » (adverbe) ; \esp{los más deportistas de la familia} = « les plus sportifs de la famille » (superlatif au pluriel, article \esp{los}).
\textbf{Barème indicatif (sur 20) :} thème 10 pts (2 pts par segment de sens : choix du comparatif/superlatif correct, articles, prépositions) ; version 10 pts (fidélité du sens 6 pts, qualité du français 4 pts).
\end{corrige}

\begin{corrige}[Exercice 11 — Production écrite et expression orale]
\textbf{A. Proposition de rédaction modèle (environ 80 mots) :}
\begin{textebox}[Mi campeón favorito — modelo]
Mi campeón favorito es Samuel Eto'o, el famoso futbolista camerunés. Para mí, es el mejor jugador africano de todos los tiempos : marcó más goles que muchos delanteros europeos y jugó tan bien como los más grandes campeones del mundo. Eto'o es más rápido y más inteligente que la mayoría de los jugadores de su generación. Además, ayuda a los jóvenes de su país : es un ejemplo de disciplina y de generosidad. Me gusta porque demuestra que con trabajo y esfuerzo, un chico africano puede llegar a ser el más grande.
\end{textebox}
\textbf{Traduction :} « Mon champion préféré est Samuel Eto'o, le célèbre footballeur camerounais. Pour moi, c'est le meilleur joueur africain de tous les temps : il a marqué plus de buts que beaucoup d'attaquants européens et a joué aussi bien que les plus grands champions du monde. Eto'o est plus rapide et plus intelligent que la majorité des joueurs de sa génération. De plus, il aide les jeunes de son pays : c'est un exemple de discipline et de générosité. Je l'aime parce qu'il montre qu'avec du travail et de l'effort, un jeune Africain peut devenir le plus grand. »
\textbf{Vérification des consignes :} présentation (nom, pays, discipline) ; deux comparatifs (\esp{más goles que}, \esp{tan bien como}, \esp{más rápido... que}) ; un superlatif (\esp{el mejor jugador africano}, \esp{el más grande}) ; conclusion personnelle (\esp{Me gusta porque...}).
\textbf{Barème indicatif (sur 10) :} respect des consignes et de la longueur 2 pts ; comparatifs et superlatif corrects 3 pts ; richesse du vocabulaire sportif 2 pts ; correction grammaticale et orthographique (accents !) 2 pts ; cohérence et conclusion personnelle 1 pt.
\textbf{Points de vigilance :} ne pas écrire \esp{más bueno} mais \esp{mejor} ; accorder les adjectifs ; soigner les accents (\esp{fútbol}, \esp{rápido}, \esp{también}) ; éviter le calque « le plus meilleur » (\esp{el más mejor} est impossible : dire \esp{el mejor}).
\textbf{B. Expression orale — trame de dialogue possible :}
\begin{itemize}
\item \esp{— ¿Cuántas veces por semana entrenas? — Entreno cinco veces por semana, más que mis compañeros.}
\item \esp{— ¿Cuál es tu ejercicio favorito? — La carrera : es el ejercicio más completo del entrenamiento.}
\item \esp{— ¿Qué haces en tu tiempo libre? — Escucho música y juego con mis amigos ; es tan relajante como dormir.}
\item \esp{— ¿Es importante la comida para ti? — Sí, como de manera más saludable que antes : muchas frutas y verduras.}
\item \esp{— ¿Quién es tu modelo? — Mi entrenador : es el mejor del campeonato y me anima todos los días.}
\end{itemize}
Critères d'évaluation orale : prononciation claire, au moins deux comparatifs et un superlatif correctement employés, fluidité et interaction naturelle entre les deux élèves.
\end{corrige}

\newpage

\coursec{Fiche bilan}

\begin{popfigure}
\begin{tikzpicture}[mindmap, grow cyclic, every node/.style={concept, text=white, font=\small},
root concept/.append style={concept color=popblue, font=\bfseries},
level 1/.append style={level distance=3.1cm, sibling angle=60},
level 1 concept/.append style={font=\footnotesize}]
\node[root concept]{Deporte, ocio y bienestar}
child[concept color=poporange]{node[align=center]{Léxico\\ \esp{el partido, entrenar, ganar}}}
child[concept color=popgreen]{node[align=center]{Comparativos\\ \esp{más... que, tan... como}}}
child[concept color=poppurple]{node[align=center]{Superlativos\\ \esp{el mejor, el más}}}
child[concept color=poppink]{node[align=center]{Traducción\\ thème / version}}
child[concept color=popgold]{node[align=center]{Producción\\ 80 palabras}}
child[concept color=popblue!70]{node[align=center]{Diálogo\\ después del partido}};
\end{tikzpicture}
\legende{Carte mentale de la leçon : le sport, les loisirs et le bien-être (expression orale et écrite).}
\end{popfigure}

\begin{bilan}
\textbf{1. Lexique clé (à connaître par cœur)}
\begin{itemize}
\item \esp{el partido} : le match ; \esp{el entrenamiento} : l'entraînement ; \esp{entrenar(se)} : s'entraîner.
\item \esp{ganar / perder / empatar} : gagner / perdre / faire match nul ; \esp{el marcador} : le score.
\item \esp{el jugador / la jugadora} : le joueur / la joueuse ; \esp{el equipo} : l'équipe ; \esp{el árbitro} : l'arbitre.
\item \esp{el campeón / la campeona} : le champion ; \esp{la afición} : les supporters ; \esp{animar} : encourager.
\item \esp{el ocio} : les loisirs ; \esp{el tiempo libre} : le temps libre ; \esp{el bienestar} : le bien-être.
\item \esp{estar en forma} : être en forme ; \esp{hacer deporte} : faire du sport ; \esp{cansarse} : se fatiguer.
\end{itemize}

\textbf{2. Grammaire : comparatifs et superlatifs}
\begin{center}
\begin{tabularx}{\linewidth}{|l|X|X|}
\hline
\rowcolor{popblueL} Type & Español & Ejemplo \\
\hline
Supériorité & \esp{más + adj. + que} & \esp{El fútbol es más popular que el tenis.} \\
\hline
Infériorité & \esp{menos + adj. + que} & \esp{El tenis es menos caro que el golf.} \\
\hline
Égalité & \esp{tan + adj. + como} & \esp{El balonmano es tan divertido como el baloncesto.} \\
\hline
Superlatif relatif & \esp{el/la/los/las más + adj. + de} & \esp{Es el jugador más rápido del equipo.} \\
\hline
Irréguliers & \esp{mejor, peor, mayor, menor} & \esp{Es el mejor campeón de África.} \\
\hline
\end{tabularx}
\end{center}
\begin{itemize}
\item Piège : \esp{mejor} et \esp{peor} remplacent \esp{más bueno} / \esp{más malo} (formes fautives à éviter).
\item \esp{mayor / menor} parlent surtout de l'âge : \esp{Mi hermano mayor} = mon frère aîné.
\item Le superlatif garde l'article défini : \esp{la más alta de la clase}.
\end{itemize}

\textbf{3. Méthodes express}
\begin{itemize}
\item \textbf{Dialogue oral} : saluer, réagir (\esp{¡Qué partido!}, \esp{¡Fue increíble!}), donner son avis avec un comparatif, conclure.
\item \textbf{Rédaction (80 mots)} : 1) présenter le sport ou le champion ; 2) dire pourquoi on l'aime (2 comparatifs minimum) ; 3) parler de sa pratique ; 4) conclure sur le bien-être.
\item \textbf{Traduction} : repérer la structure de comparaison en français, choisir \esp{más/menos/tan}, vérifier l'accord de l'adjectif.
\end{itemize}
\end{bilan}

\begin{aretenir}[Checklist avant le Bac]
\begin{itemize}
\item[$\square$] Je connais le vocabulaire du match : \esp{partido, marcador, equipo, árbitro, empatar}.
\item[$\square$] Je sais former les trois comparatifs : \esp{más... que}, \esp{menos... que}, \esp{tan... como}.
\item[$\square$] Je n'écris jamais \esp{más bueno} : je dis \esp{mejor}.
\item[$\square$] Je sais former le superlatif : \esp{el más + adjectif + de}.
\item[$\square$] Je fais toujours l'accord de l'adjectif : \esp{la jugadora más rápida}.
\item[$\square$] Je sais réagir dans un dialogue : \esp{¡Qué partido más emocionante!}
\item[$\square$] Je sais rédiger 80 mots sur mon sport préféré en quatre paragraphes courts.
\item[$\square$] Je place au moins deux comparatifs dans ma production écrite.
\item[$\square$] Je vérifie les accents : \esp{fútbol, campeón, afición, emocionante}.
\item[$\square$] Je distingue \esp{ganar} (gagner) et \esp{ganarse la vida} (gagner sa vie).
\item[$\square$] Je relis ma traduction pour éviter les calques du français.
\end{itemize}
\end{aretenir}

\findecours

\end{document}
